\documentclass[12pt,a4paper]{article}
\usepackage[utf8]{inputenc}[spanish]
\usepackage{amsmath}
\usepackage{amsfonts}
\usepackage{amssymb}
\usepackage{lmodern}
\usepackage{amsmath}
\usepackage{amsbsy}
\usepackage{amsthm}
\usepackage{enumerate}
\usepackage{graphicx}
\usepackage{mathtools}
\usepackage{stackrel}
\usepackage{multicol}
\usepackage{hyperref}
\usepackage{cleveref}
\usepackage[left=2cm,right=2cm,top=2cm,bottom=2cm]{geometry}
\providecommand{\abs}[1]{\lvert#1\rvert}

% Párrafo
\newcommand{\PN}{\par\noindent}
% Comillas izquierdas
\newcommand{\CL}{\textquotedblleft}
% Comillas derechas
\newcommand{\CR}{\textquotedblright}


% Conjuntos
\newcommand{\SIGMA}{\Sigma^{\ast}}
\newcommand{\TAU}{T^{\tau}}
\newcommand{\FT}{F^{\tau}}


% Supremo e ínfimo
\newcommand{\SU}{\mathsf{s}}
\newcommand{\IN}{\mathsf{i}}


% Poset
\newcommand{\POSET}{(P, \leq)}
% Poset primo
\newcommand{\POSETPRIMO}{(P^{\prime}, \leq^{\prime})}

% Reticulado
\newcommand{\RET}{(L, \SU, \IN)}
% Reticulado primo
\newcommand{\RETPRIMO}{(L^{\prime}, \SU^{\prime}, \IN^{\prime})}
% Reticulado cocientado
\newcommand{\RETCOCIENTADO}{(L/\theta, \mathsf{\tilde{s}}, \mathsf{\tilde{\imath}})}

% Reticulado acotado
\newcommand{\ACOTADO}{(L, \SU, \IN, 0, 1)}
% Reticulado acotado primo
\newcommand{\ACOTADOPRIMO}{(L^{\prime}, \SU^{\prime}, \IN^{\prime}, 0^{\prime}, 1^{\prime})}
% Reticulado acotado cocientado
\newcommand{\ACOTADOCOCIENTADO}{(L/\theta, \mathsf{\tilde{s}}, \mathsf{\tilde{\imath}}, 0/\theta, 1/\theta)}

% Reticulado complementado
\newcommand{\COMPLEMENTADO}{(L, \SU, \IN, ^{c}, 0, 1)}
% Reticulado complementado primo
\newcommand{\COMPLEMENTADOPRIMO}{(L^{\prime}, \SU^{\prime}, \IN^{\prime}, ^{c^{\prime}} 0^{\prime}, 1^{\prime})}
% Reticulado complementado cocientado
\newcommand{\COMPLEMENTADOCOCIENTADO}{(L/\theta, \mathsf{\tilde{s}}, \mathsf{\tilde{\imath}}, ^{c}/\theta, 0/\theta, 1/\theta)}

% Álgebra de Boole
\newcommand{\BOOLE}{(B, \SU, \IN, ^{c}, 0, 1)}


% Teoremas
\newtheorem{theorem}[equation]{Theorem}
\newtheorem{lemma}[equation]{Lemma}
\newtheorem{proposition}[equation]{Proposition}
\newtheorem{corollary}[equation]{Corollary}

\author{Agustín Curto, agucurto95@gmail.com}
\title{Resumen de teoremas para el final \\ de Lógica}
\date{2017}

\begin{document}
	\clearpage\maketitle
	\thispagestyle{empty}
	\tableofcontents

	\vspace{5cm}
	\PN \textbf{Nota:} Este resumen se corresponde con la materia dictada en el año 2017. El autor no se
	responsabiliza de posibles cambios que pudiesen realizarse en los temas dictados en la misma, así como tampoco de
	errores involuntarios que pudiesen existir en dicho resumen.

	\vspace{\fill}
	\begin{center}
		Por favor, mejorá este documento en github
		\includegraphics[width=1cm]{graphics/github.png} \\
		https://github.com/ResumenesFaMAF/resumenLogica
	\end{center}

	\pagebreak

	\section{Estructuras algebráicas ordenadas}

  % % Lemma 89: Con prueba. Lemma 1.
  % \begin{lemma}
  %   \PN Sean $(P, \leq)$ y $(P^{\prime}, \leq^{\prime})$ posets. Supongamos que $F$ es un isomorfismo de $(P, \leq)$ en
  %   $(P^{\prime}, \leq^{\prime})$, entonces:
  %
  %   \begin{enumerate}[a)]
  %     \item Para cada $S \subseteq P$ y cada $a \in P$, se tiene que $a$ es \textbf{cota superior} (resp.
  %       \textbf{inferior}) de $S$ si y solo si $F(a)$ es \textbf{cota superior} (resp. \textbf{inferior}) de $F(S)$.
  %     \item Para cada $S \subseteq P$, se tiene que $\exists \ \sup (S)$ si y solo si $\exists \ \sup (F(S))$ y en el
  %       caso de que existan tales elementos se tiene que $F(\sup (S)) = \sup (F(S))$.
  %     \item $P$ tiene $1$ (resp. $0$) si y solo si $P^{\prime }$ tiene $1$ (resp. $0$) y en tal caso tales elementos
  %       están conectados por $F$.
  %     \item Para cada $m \in P$, $m$ es \textbf{maximal} (resp. \textbf{minimal}) si y solo si $F(m)$ es
  %       \textbf{maximal} (resp. \textbf{minimal}).
  %     \item Para $a, b \in P$, tenemos que $a \prec b$ si y solo si $F(a) \prec^{\prime} F(b)$.
  %   \end{enumerate}
  % \end{lemma}
  % \begin{proof}
  %   \begin{enumerate}[a)]
  %     \item Probaremos solo el caso de la \textbf{cota superior}.
  %       \PN \begin{tabular}{|c|} \hline $\Rightarrow$ \\\hline \end{tabular} Supongamos que $a$ es \textbf{cota
  %       superior} de $S$, veamos entonces que $F(a)$ es \textbf{cota superior} de $F(S)$. Sean:
  %       \begin{itemize}
  %         \item $x \in F(S)$
  %         \item $s \in S$ tal que $x = F(s)$.
  %       \end{itemize}
  %
  %       \PN Ya que $s \leq a$, tenemos que $x = F(s) \leq^{\prime} F(a)$. Luego, $F(a)$ es \textbf{cota superior}.
  %
  %       \PN \begin{tabular}{|c|} \hline $\Leftarrow$ \\\hline \end{tabular} Supongamos ahora que $F(a)$ es \textbf{cota
  %       superior} de $F(S)$ y veamos entonces que $a$ es cota superior de $S$.
  %
  %       \PN Sea $s \in S$, ya que $F(s) \leq^{\prime} F(a)$, tenemos que $s = F^{-1}(F(s)) \leq^{\prime} F^{-1}(F(a)) =
  %       a$. Por lo tanto, $a$ es \textbf{cota superior}.
  %
  %     \item \begin{tabular}{|c|} \hline $\Rightarrow$ \\\hline \end{tabular} Supongamos existe $\sup (S)$. Veamos
  %       que $F(\sup (S))$ es el supremo de $F(S)$. Por el iniciso (a) $F(\sup (S))$ es cota superior de $F(S)$. Veamos
  %       que es la menor de las cotas superiores. Supongamos $b^{\prime}$ cota superior de $F(S)$, entonces
  %       $F^{-1}(b^{\prime})$ es cota superior de $S$, es decir, $\sup (S) \leq F^{-1}(b^{\prime})$, produciendo
  %       $F(\sup (S)) \leq^{\prime} b^{\prime}$. Por lo tanto, $F(\sup (S))$ es el supremo de $F(S)$.
  %
  %       \PN \begin{tabular}{|c|} \hline $\Leftarrow$ \\\hline \end{tabular} Supongamos existe $\sup (F(S))$. Veamos
  %       que $F^{-1}(\sup (F(S)))$ es el supremo de $S$. Nuevamente, por el iniciso (a) $F^{-1}(\sup (F(S)))$ es cota
  %       superior de $S$. Veamos que es la menor de las cotas superiores. Supongamos $b$ cota superior de $S$, entonces
  %       $F(b)$ es cota superior de $F(S)$, es decir, $\sup (F(S)) \leq F(b)$, produciendo $F^{-1}(\sup (F(S))) \leq b$.
  %       Por lo tanto, $F^{-1}(\sup (F(S)))$ es el supremo de $S$.
  %
  %     \item Se desprende del inciso (b) tomando $S = P$.
  %     \item Probaremos solo el caso \textbf{maximal}.
  %       \PN \begin{tabular}{|c|} \hline $\Rightarrow$ \\\hline \end{tabular} Supongamos que $m$ es maximal de
  %       $(P, \leq)$. Veamos que $F(m)$ es maximal de $(P^{\prime}, \leq^{\prime})$. Supongamos que $F(m)$ no es maximal
  %       de $(P^{\prime}, \leq^{\prime})$, es decir, $F(m) <^{\prime} b^{\prime} \ \forall b^{\prime} \in P^{\prime}$.
  %       Dado que $F$ es isomorfismo:
  %       \begin{eqnarray*}
  %   			F^{-1}(F(m)) < F^{-1}(b^{\prime}) \\
  %   			m < F^{-1}(b^{\prime})
  %   		\end{eqnarray*}
  %       \PN Lo cual es un absurdo, dado que $m$ es maximal de $(P, \leq)$. Por lo tanto, $F(m)$ es maximal de
  %       $(P^{\prime}, \leq^{\prime})$.
  %
  %       \PN \begin{tabular}{|c|} \hline $\Leftarrow$ \\\hline \end{tabular} Supongamos que $F(m)$ es maximal de
  %       $(P^{\prime}, \leq^{\prime})$. Veamos que $m$ es maximal de $(P, \leq)$. Supongamos que $m$ no es maximal
  %       de $(P, \leq)$, es decir, $m < b \ \forall b \in P$. Dado que $F$ es isomorfismo:
  %       \[
  %   			F(m) < F(b)
  %   		\]
  %       \PN Lo cual es un absurdo, dado que $F(m)$ es maximal de $(P^{\prime}, \leq^{\prime})$. Por lo tanto, $m$ es
  %       maximal de $(P, \leq)$.
  %
  %     \item \begin{tabular}{|c|} \hline $\Rightarrow$ \\\hline \end{tabular} Supongamos $a \prec b$, veamos que $F(a)
  %       \prec^{\prime} F(b)$. Debemos ver:
  %       \begin{enumerate}[1)]
  %         \item $F(a) <^{\prime} F(b)$
  %         \item $\nexists z^{\prime}$ tal que $F(a) < z^{\prime} < F(b)$
  %       \end{enumerate}
  %
  %       \PN Ya que $a \prec b$, por definición tenemos: \begin{tabular}{|c|} \hline $a < b$ y $\nexists z$ tal que
  %       $a < z < b$ \\\hline \end{tabular} $(\star)$
  %
  %       \PN Dado que la función $F$ es un isomorfismo, se cumple (1). Veamos que se cumple (2), supongamos que $\exists
  %       z^{\prime}$ tal que $F(a) < z^{\prime} < F(b)$. Luego, nuevamente utilizando que $F$ es isomorfismo, tenemos:
  %       \begin{eqnarray*}
  %   			F^{-1}(F(a)) < &F^{-1}(z^{\prime})& < F^{-1}(F(b)) \\
  %   			a < &F^{-1}(z^{\prime})& < b
  %   		\end{eqnarray*}
  %       \PN Lo cual, contradice $(\star)$, el absurdo vino de suponer que $\exists z^{\prime}$ tal que
  %       $F(a) < z^{\prime} < F(b)$, por lo tanto $\nexists z^{\prime}$ tal que $F(a) < z^{\prime} < F(b)$.
  %
  %       \PN Finalmente, dado que se cumplen los puntos (1) y (2), se cumple también $F(a) \prec^{\prime} F(b)$.
  %
  %       \PN \begin{tabular}{|c|} \hline $\Leftarrow$ \\\hline \end{tabular} Supongamos $F(a) \prec^{\prime} F(b)$,
  %       veamos que $a \prec b$.
  %
	% 	    \PN Ya que $F^{-1}: (P^{\prime}, \leq^{\prime}) \rightarrow (P, \leq)$ es isomorfismo, por lo ya visto tenemos:
  %   		\begin{eqnarray*}
  %   			F^{-1}(F(a)) &\prec& F^{-1}(F(b)) \\
  %   			a &\prec& b
  %   		\end{eqnarray*}
  %   \end{enumerate}
  % \end{proof}
  %
  % % Lemma 90: Con prueba. Lemma 2.
  % \begin{lemma}
  %   \PN Dado un reticulado $(L, \leq)$ y elementos $x, y, z, w \in L$, se cumplen las siguientes propiedades:
  %   \begin{multicols}{2}
  %     \begin{enumerate}
  %       \item $x \leq x \ \SU \ y$
  %       \item $x \ \IN \ y \leq x$
  %       \item $x \ \SU \ x = x \ \IN \ x = x$
  %       \item $x \ \SU \ y = y \ \SU \ x$
  %       \item $x \ \IN \ y = y \ \IN \ x$
  %       \item $x \leq y \Leftrightarrow x \ \SU \ y = y \Leftrightarrow x \ \IN \ y = x$
  %       \item $x \ \SU \ (x \ \IN \ y) = x$
  %       \item $x \ \IN(x \ \SU \ y) = x$
  %       \item $(x \ \SU \ y) \ \SU \ z = x \ \SU \ (y \ \SU \ z)$
  %       \item $(x \ \IN \ y) \ \IN \ z = x \ \IN \ (y \ \IN \ z)$
  %       \item Si $x \leq z$ e $y \leq w$ entonces:
  %         \begin{itemize}
  %           \item $x \ \SU \ y \leq z \ \SU \ w$
  %           \item $x \ \IN \ y \leq z \ \IN \ w$
  %         \end{itemize}
  %       \item $(x \ \IN \ y) \ \SU \ (x \ \IN \ z) \leq x \ \IN \ (y \ \SU \ z)$
  %     \end{enumerate}
  %   \end{multicols}
  % \end{lemma}
  % \begin{proof}
  %   \begin{enumerate}
  %     \item \begin{tabular}{|c|} \hline $x \leq x \ \SU \ y$ \\\hline \end{tabular}
  %       \begin{eqnarray*}
  %         \textup{Si } x \leq y &\Rightarrow& x \leq x \ \SU \ y = y \Rightarrow x \leq y \\
  %         \textup{Si } y \leq x &\Rightarrow& x \leq x \ \SU \ y = x \Rightarrow x \leq x
  %       \end{eqnarray*}
  %     \item \begin{tabular}{|c|} \hline $x \ \IN \ y \leq x$ \\\hline \end{tabular}
  %       \begin{eqnarray*}
  %         \textup{Si } x \leq y \Rightarrow x \ \IN \ y &\leq& x \\
  %         x &\leq& x \\
  %         \textup{Si } y \leq x \Rightarrow x \ \IN \ y &\leq& x \\
  %         y &\leq& x
  %       \end{eqnarray*}
  %     \item \begin{tabular}{|c|} \hline $x \ \SU \ x = x \ \IN \ x = x$ \\\hline \end{tabular}
  %       \begin{eqnarray*}
  %         x \ \SU \ x &=& x \ \IN \ x = x \\
  %         x &=& x = x
  %       \end{eqnarray*}
  %       % TODO
  %     \item \begin{tabular}{|c|} \hline $x \ \SU \ y = y \ \SU \ x$ \\\hline \end{tabular}
  %       \begin{eqnarray*}
  %         \textup{Si } x \leq y \Rightarrow x \ \SU \ y &=& y \ \SU \ x \\
  %         y &=& y \\
  %         \textup{Si } y \leq x \Rightarrow x \ \SU \ y &=& y \ \SU \ x \\
  %         x &=& x
  %       \end{eqnarray*}
  %     \item \begin{tabular}{|c|} \hline $x \ \IN \ y = y \ \IN \ x$ \\\hline \end{tabular}
  %       \begin{eqnarray*}
  %         \textup{Si } x \leq y \Rightarrow x \ \IN \ y &=& y \ \IN \ x \\
  %         x &=& x \\
  %         \textup{Si } y \leq x \Rightarrow x \ \IN \ y &=& y \ \IN \ x \\
  %         y &=& y
  %       \end{eqnarray*}
  %     \item \begin{tabular}{|c|} \hline $x \leq y \Leftrightarrow x \ \SU \ y = y \Leftrightarrow x \ \IN
  %       \ y = x$ \\\hline \end{tabular}
  %       % TODO
  %     \item Ya que $x\mathsf{\;i\;}y\leq x$, (6) nos dice que $(x\mathsf{\;i\;}y)\; \SU\;x=x$, por lo cual $x\;\SU\;(x\mathsf{\;i\;}y)=x$.
  %     \item Similar a (7).
  %     \item Notese que $x\;\SU\;(y\;\SU\;z)$ es cota superior de $ \{x,y,z\}$ ya que onviamente $x\leq x\;\SU\;(y\;\SU\;z)$ y ademas
  %       $\displaystyle \begin{array}{rcl} y & \leq & (y\;\SU\;z)\leq x\;\SU\;(y\;\SU\;z) \\ z & \leq & (y\;\SU\;z)\leq x\;\SU\;(y\;\SU\;z) \end{array} $
  %
  %       Ya que $x\;\SU\;(y\;\SU\;z)$ es cota superior de $\{x,y\}$, tenemos que $x\;\SU\;y\leq x\;\SU\ (y\;\SU\;z)$, por lo cual $x\;\SU\;(y\;\SU\;z)$ es cota superior del conjunto $\{x\; \SU\;y,z\}$, lo cual dice que $(x\;\SU\;y)\;\SU\;z\leq x\;\SU\;(y\;\SU\;z)$. Analogamente se puede probar que $x\; \SU\;(y\;\SU\;z)\leq (x\;\SU\;y)\;\SU\;z$.
  %     \item Similar a (9).
  %     \item Ya que
  %       $\displaystyle \begin{array}{rcl} x & \leq & z\leq z\;\SU\;w \\ y & \leq & w\leq z\;\SU\;w \end{array} $
  %
  %       tenemos que $z\;\SU\;w$ es cota superior de $\{x,y\}$ lo cual dice que $x\;\SU\;y\leq z\;\SU\;w$. La otra desigualdad es analoga.
  %     \item Ya que
  %       $\displaystyle \begin{array}{rcl} (x\mathsf{\;i\;}y),(x\mathsf{\;i\;}z) & \leq & x \\ (x\mathsf{\;i\;}y),(x\mathsf{\;i\;}z) & \leq & y\;\SU\;z \end{array} $
  %
  %       tenemos que $(x\;\IN\;y),(x\mathsf{\;i\;}z)\leq x\mathsf{\;i\;}(y\; \SU\;z)$, por lo cual $(x\mathsf{\;i\;}y)\;\SU\;(x\mathsf{\;i\; }z)\leq x\mathsf{\;i\;}(y\;\SU\;z)$.
  %   \end{enumerate}
  % \end{proof}

  % Lemma 91: Con prueba. Lemma 3.
  \begin{lemma}
    \PN Sea $(L, \leq)$ un reticulado, dados elementos $x_{1}, \dotsc, x_{n} \in L$, con $n \geq 2$, se tiene
    \[
      \begin{array}{rcl}
        (\dotsc (x_{1} \ \SU \ x_{2}) \ \SU \ \dotsc) \ \SU \ x_{n} &=& \sup (\{x_{1}, \dotsc, x_{n}\}) \\
        (\dotsc (x_{1} \ \IN \ x_{2}) \ \IN \ \dotsc) \ \IN \ x_{n} &=& \inf (\{x_{1}, \dotsc, x_{n}\})
      \end{array}
    \]
  \end{lemma}
  \begin{proof}
    \PN Probaremos por inducción en $n$.

    \vspace{3mm}
    \PN \underline{Caso Base:} \begin{tabular}{|c|} \hline $n = 2$ \\\hline \end{tabular}
    \[
      \begin{array}{rcl}
        x_{1} \ \SU \ x_{2} &=& \sup(\{x_{1}, x_{2}\}) \\
        x_{1} \ \IN \ x_{2} &=& \inf(\{x_{1}, x_{2}\})
      \end{array}
    \]

    \PN Lo cual vale, dado que es la definición.

    \vspace{3mm}
		\PN \underline{Caso Inductivo:} \begin{tabular}{|c|} \hline $n > 2$ \\\hline \end{tabular}

    \PN Supongamos ahora que vale para $n$ y veamos entonces que vale para $n+1$. Sean $x_{1}, \dotsc, x_{n+1} \in L$,
    por hipótesis inductiva tenemos que:
    \begin{eqnarray*}
      (\dotsc (x_{1} \ \SU \ x_{2}) \ \SU \ \dotsc) \ \SU \ x_{n} &=& \sup(\{x_{1}, \dotsc, x_{n}\}) \ \ (\star_{1}) \\
      (\dotsc (x_{1} \ \IN \ x_{2}) \ \IN \ \dotsc) \ \IN \ x_{n} &=& \inf(\{x_{1}, \dotsc, x_{n}\}) \ \ \ (\star_{2})
    \end{eqnarray*}


    \PN Veamos entonces que:
    \begin{eqnarray*}
      ((\dotsc(x_{1} \ \SU \ x_{2}) \ \SU \ \dotsc) \ \SU \ x_{n}) \ \SU \ x_{n+1} &=& \sup(\{x_{1}, \dotsc, x_{n+1}\})
        \ \ (\dag_{1}) \\
      ((\dotsc(x_{1} \ \IN \ x_{2}) \ \IN \ \dotsc) \ \IN \ x_{n}) \ \IN \ x_{n+1} &=& \inf(\{x_{1}, \dotsc, x_{n+1}\})
        \ \ \ (\dag_{2})
    \end{eqnarray*}

    \PN Para ello debemos ver $((\dotsc(x_{1} \ \SU \ x_{2}) \ \SU \ \dotsc) \ \SU \ x_{n}) \ \SU \ x_{n+1}$ es cota
    superior de $\{x_{1}, \dotsc, x_{n+1}\}$ y que es la menor de las cotas superiores. Además, que $((\dotsc(x_{1} \
    \IN \ x_{2}) \ \IN \ \dotsc) \ \IN \ x_{n}) \ \IN \ x_{n+1}$ es cota inferior de $\{x_{1}, \dotsc, x_{n+1}\}$ y que
    es la mayor de las cotas inferiores.

    \vspace{5mm}
    \PN Es fácil ver que $((\dotsc(x_{1} \ \SU \ x_{2}) \ \SU \ \dotsc) \ \SU \ x_{n}) \ \SU \ x_{n+1}$ es cota superior
    de $ \{x_{1}, \dotsc, x_{n+1}\}$. Supongamos que $z$ es otra cota superior de $\{x_{1}, \dotsc, x_{n+1}\}$. Ya que
    $z$ es también cota superior del conjunto $\{x_{1}, \dotsc, x_{n}\}$, por $(\star_{1})$ tenemos que:
    \[
      (\dotsc (x_{1} \ \SU \ x_{2}) \ \SU \ \dotsc) \ \SU \ x_{n} \leq z
    \]

    \PN Además, dado que $x_{n+1} \leq z$, tenemos que:
    \[
      ((\dotsc (x_{1} \ \SU \ x_{2}) \ \SU \ \dotsc) \ \SU \ x_{n}) \ \SU \ x_{n+1} \leq z
    \]

    \PN Por lo tanto, vale $(\dag_{1})$.

    \vspace{5mm}
    \PN Nuevamente, es fácil ver que $((\dotsc(x_{1} \ \IN \ x_{2}) \ \IN \ \dotsc) \ \IN \ x_{n}) \ \IN \ x_{n+1}$ es
    cota inferior de $ \{x_{1}, \dotsc, x_{n+1}\}$. Supongamos que $z^{\prime}$ es otra cota inferior de $\{x_{1},
    \dotsc, x_{n+1}\}$. Ya que $z^{\prime}$ es también cota inferior del conjunto $\{x_{1}, \dotsc, x_{n}\}$, por
    $(\star_{2})$ tenemos que:
    \[
      z^{\prime} \leq (\dotsc (x_{1} \ \IN \ x_{2}) \ \IN \ \dotsc) \ \IN \ x_{n}
    \]

    \PN Además, dado que $z^{\prime} \leq x_{n+1}$, tenemos que:
    \[
      z^{\prime} \leq ((\dotsc (x_{1} \ \IN \ x_{2}) \ \IN \ \dotsc) \ \IN \ x_{n}) \ \IN \ x_{n+1}
    \]

    \PN Por lo tanto, vale $(\dag_{2})$.

  \end{proof}

  % % Theorem 92: Con prueba. Theorem 4.
  % \begin{theorem}
  %   \PN Sea $(L, \SU, \IN)$ un reticulado, la relación binaria definida por:
  %   \[
  %     x \leq y \Leftrightarrow x \ \SU \ y=y
  %   \]
  %
  %   \PN es un orden parcial sobre $L$ para el cual se cumple:
  %   \[
  %     \begin{array}{rcl}
  %       \sup (\{x, y\}) &=& x \ \SU \ y \\
  %       \inf (\{x, y\}) &=& x \ \IN \ y
  %     \end{array}
  %   \]
  % \end{theorem}
  % \begin{proof}
  %   \begin{itemize}
  %     \item \underline{Reflexiva:}
  %     \item \underline{Antisimétrica:}
  %     \item \underline{Transitiva:} Supongamos que $x \leq y$ e $y \leq z$, entonces:
  %       \[
  %         x \ \SU \ z = x \ \SU \ (y \ \SU \ z) = (x \ \SU \ y) \ \SU \ z = y \ \SU \ z = z
  %       \]
  %
  %       \PN por lo cual $x \leq z$. Veamos ahora que $\sup (\{x,y\})=x\ \SU \y$. Es claro que $x\ \SU \y$ es una cota superior del conjunto $\{x,y\}$. Supongamos $x,y\leq z$. Entonces
  %     $\displaystyle (x\ \SU \y)\ \SU \z=x\ \SU \(y\ \SU \z)=x\; \mathsf{s\;}z=z, $
  %
  %     por lo que $x\ \SU \y\leq z$. Es decir que $x\ \SU \y$ es la menor cota superior.
  %     Para probar que $\inf (\{x,y\})=x\mathsf{\;i\;}y$, probaremos que para todo $ u,v\in L$,
  %
  %     $\displaystyle u\leq v\text{ si y solo si }u\mathsf{\;i\;}v=u, $
  %
  %     lo cual le permitira al lector aplicar un razonamiento similar al usado en el caso de la operacion $\SU$. Supongamos que $u\;\SU\;v=v$. Entonces $u\mathsf{\;i\;}v=u\mathsf{\;i\;}(v\;\SU\;v)=u$. Reciprocamente si $u\mathsf{\;i\;}v=u$, entonces $u\;\SU\;v=(u\mathsf{ \;i\;}v)\;\SU\;v=v$, por lo cual $u\leq v$.
  %   \end{itemize}
  % \end{proof}
  %
  % % Lemma 93: Con prueba. Lemma 5.
  % \begin{lemma}
  %   \PN Si $F: (L, \SU, \IN) \rightarrow (L^{\prime}, \SU^{\prime}, \IN^{\prime})$ es un
  %   homomorfismo biyectivo, entonces $F$ es un isomorfismo.
  % \end{lemma}
  % \begin{proof}
  %   \PN Debemos probar que $F^{-1}$ es un homomorfismo. Sean $F(x), F(y)$ dos elementos cualesquiera de $L^{\prime}$,
  %   tenemos que:
  %   \begin{eqnarray*}
  %     F^{-1}(F(x) \ \SU^{\prime} \ F(y)) &=& F^{-1}(F(x \ \SU \ y)) \\
  %     &=& x \ \SU \ y \\
  %     &=& F^{-1}(F(x)) \ \SU \ F^{-1}(F(y))
  %   \end{eqnarray*}
  %
  %   \PN Luego, $F^{-1}$ es homomorfismo y por lo tanto $F$ es isomorfismo.
  % \end{proof}
  %
  % % Lemma 94: Con prueba. Lemma 6.
  % \begin{lemma}
  %   \PN Sean $(L, \SU, \IN)$ y $(L^{\prime}, \SU^{\prime}, \IN^{\prime})$ reticulados y sea
  %   $F: (L, \SU, \IN) \rightarrow (L^{\prime}, \SU^{\prime}, \IN^{\prime})$ un homomorfismo,
  %   entonces $I_{F}$ es un subuniverso de $(L^{\prime}, \SU^{\prime}, \IN^{\prime})$.
  % \end{lemma}
  % \begin{proof}
  %   Ya que $L$ es no vacio tenemos que $I_{F}$ tambien es no vacio. Sean $a,b\in I_{F}$. Sean $x,y\in L$ tales que $F(x)=a$ y $F(y)=b$. Se tiene que
  %   \begin{eqnarray*}
  %     a\;\SU^{\prime }\ b &=& F(x)\;\SU^{\prime }\ F(y)=F(x\mathsf{ \;s\;}y)\in I_{F} \\
  %     a\;\IN^{\prime }\ b & =& F(x)\;\IN^{\prime }\ F(y)=F(x\mathsf{ \;i\;}y)\in I_{F}
  %   \end{eqnarray*}
  %
  %   \PN por lo cual $I_{F}$ es cerrada bajo $\SU^{\prime }$ e $\IN ^{\prime }$.
  % \end{proof}
  %
  % % Lemma 95: Con prueba. Lemma 7.
  % \begin{lemma}
  %   \PN Sean $(L,\SU,\IN)$ y $(L^{\prime },\SU^{\prime }, \IN^{\prime })$ reticulados y sean $(L\leq )$ y $(L^{\prime },\leq ^{\prime })$ los posets asociados. Sea $F:L\rightarrow L^{\prime }$ una funcion. Entonces $F$ es un isomorfismo de $(L,\SU,\IN)$ en $ (L^{\prime },\SU^{\prime },\IN^{\prime })$ si y solo si $F$ es un isomorfismo de $(L,\leq )$ en $(L^{\prime },\leq ^{\prime })$.
  % \end{lemma}
  % \begin{proof}
  %   Supongamos $F$ es un isomorfismo de $(L,\SU,\IN)$ en $ (L^{\prime },\SU^{\prime },\IN^{\prime })$. Sean $x,y\in L$, tales que $x\leq y$. Tenemos que $y=x\mathsf{\;s\;}y$ por lo cual $F(y)=F(x \mathsf{\;s\;}y)=F(x)\mathsf{\;s^{\prime }\;}F(y)$, produciendo $F(x)\leq ^{\prime }F(y)$. En forma similar se puede ver que $F^{-1}$ es tambien un homomorfismo de $(L^{\prime },\leq ^{\prime })$ en $(L,\leq )$. Si $F$ es un isomorfismo de $(L,\leq )$ en $(L^{\prime },\leq ^{\prime })$, entonces el Lema 89 nos dice que $F$ y $F^{-1}$ respetan las operaciones de supremo e infimo por lo cual $F$ es un isomorfismo de $(L,\SU,\mathsf{ i})$ en $(L^{\prime },\SU^{\prime },\IN^{\prime })$. $\Box$
  % \end{proof}
  %
  % % Lemma 96: Con prueba. Lemma 8.
  % \begin{lemma}
  %   $(L/\theta ,\mathsf{\tilde{s}},\mathsf{\tilde{\imath}})$ es un reticulado. El orden parcial $\tilde{\leq}$ asociado a este reticulado cumple
  %   $\displaystyle x/\theta \tilde{\leq}y/\theta \text{ sii }y\theta (x\mathsf{\;s\;}y) $
  % \end{lemma}
  % \begin{proof}
  %   Veamos que la estructura $(L/\theta ,\mathsf{\tilde{s}},\mathsf{\tilde{\imath }})$ cumple (I4). Sean $x/\theta $, $y/\theta $, $z/\theta $ elementos cualesquiera de $L/\theta $. Tenemos que
  %
  %   $\displaystyle \begin{array}{ccl} (x/\theta \mathsf{\;\tilde{s}\;}y/\theta )\;\mathsf{\tilde{s}}\;z/\theta & = & (x\mathsf{\;s\;}y)/\theta \;\mathsf{\tilde{s}}\;z/\theta \\ & = & ((x\mathsf{\;s\;}y)\;\SU\;z)/\theta \\ & = & (x\mathsf{\;s\;}(y\;\SU\;z))/\theta \\ & = & x/\theta \;\mathsf{\tilde{s}}\;(y\;\SU\;z)/\theta \\ & = & x/\theta \mathsf{\;\tilde{s}\;}(y/\theta \;\mathsf{\tilde{s}} \;z/\theta ) \end{array} $
  %
  %   En forma similar se puede ver que la estructura $(L/\theta ,\mathsf{\tilde{s} },\mathsf{\tilde{\imath}})$ cumple el resto de las identidades que definen reticulado.
  % \end{proof}
  %
  % % Corollary 97: Con prueba. Corollary 9.
  % \begin{corollary}
  %   Sea $(L,\SU,\IN)$ un reticulado en el cual hay un elemento maximo $1$ (resp. minimo $0$). Entonces si $\theta $ es una congruencia sobre $(L,\SU,\IN)$, $1/\theta $ (resp. $0/\theta $) es un elemento maximo (resp. minimo) de $(L/\theta ,\mathsf{\tilde{s}},\mathsf{ \tilde{\imath}})$.
  % \end{corollary}
  % \begin{proof}
  %   Ya que $1\theta (x\mathsf{\;s\;}1)$, para cada $x\in L$, tenemos que $ x/\theta \tilde{\leq}1/\theta $, para cada $x\in L$. $\Box$
  % \end{proof}
  %
  % % Lemma 98: Con prueba. Lemma 10.
  % \begin{lemma}
  %   Si $F:(L,\SU,\IN)\rightarrow (L^{\prime },\SU^{\prime }, \IN^{\prime })$ es un homomorfismo de reticulados, entonces $\ker F$ es una congruencia sobre $(L,\SU,\IN)$.
  % \end{lemma}
  % \begin{proof}
  %   Dejamos al lector ver que $\ker F$ es una relacion de equivalencia. Supongamos $x\ker Fx^{\prime }$ y $y\ker Fy^{\prime }$. Entonces
  %
  %   $\displaystyle F(x\mathsf{\;s\;}y)=F(x)\mathsf{\;s^{\prime }\;}F(y)=F(x^{\prime })\mathsf{ \;s^{\prime }\;}F(y^{\prime })=F(x^{\prime }\mathsf{\;s\;}y^{\prime }) $
  %
  %   lo cual nos dice que $(x\mathsf{\;s\;}y)\ker F(x^{\prime }\mathsf{\;s\;} y^{\prime })$. En forma similar tenemos que $(x\mathsf{\;i\;}y)\ker F(x^{\prime }\mathsf{\;i\;}y^{\prime })$. $\Box$
  % \end{proof}
  %
  % % Lemma 99: Con prueba. Lemma 11.
  % \begin{lemma}
  %   Sea $(L,\SU,\IN)$ un reticulado y sea $\theta $ una congruencia sobre $(L,\SU,\IN)$. Entonces $\pi _{\theta }$ es un homomorfismo de $(L,\SU,\IN)$ en $(L/\theta ,\mathsf{\tilde{ s}},\mathsf{\tilde{\imath}})$. Ademas $\ker \pi _{\theta }=\theta $.
  % \end{lemma}
  % \begin{proof}
  %   Sean $x,y\in L$. Tenemos que
  %
  %   $\displaystyle \pi _{\theta }(x\mathsf{\;s\;}y)=(x\mathsf{\;s\;}y)/\theta =x/\theta \mathsf{ \;\tilde{s}\;}y/\theta =\pi _{\theta }(x)\mathsf{\;\tilde{s}\;}\pi _{\theta }(y) $
  %
  %   por lo cual $\pi _{\theta }$ preserva la operacion supremo. Para la operacion infimo es similar. $\Box$
  % \end{proof}
  %
  % % Lemma 100: Con prueba. Lemma 12.
  % \begin{lemma}
  %   Si $F:(L,\SU,\IN,0,1)\rightarrow (L^{\prime },\SU ^{\prime },\IN^{\prime },0^{\prime },1^{\prime })$ un homomorfismo biyectivo, entonces $F$ es un isomorfismo
  % \end{lemma}
  % \begin{proof}
  %   Similar a la prueba del Lemma 93. $\Box$
  % \end{proof}
  %
  % % Lemma 101: Con prueba. Lemma 13.
  % \begin{lemma}
  %   Si $F:(L,\SU,\IN,0,1)\rightarrow (L^{\prime },\SU ^{\prime },\IN^{\prime },0^{\prime },1^{\prime })$ es un homomorfismo, entonces $I_{F}$ es un subuniverso de $(L^{\prime },\SU ^{\prime },\IN^{\prime },0^{\prime },1^{\prime })$.
  % \end{lemma}
  % \begin{proof}
  %   Ya que $F$ es un homomorfismo de $(L,\SU,\IN)$ en $ (L^{\prime },\SU^{\prime },\IN^{\prime })$ tenemos que $I_{F}$ es subuniverso de $(L^{\prime },\SU^{\prime },\IN^{\prime })$ lo cual ya que $0^{\prime },1^{\prime }\in I_{F}$ implica que $I_{F}$ es un subuniverso de $(L^{\prime },\SU^{\prime },\IN^{\prime },0^{\prime },1^{\prime })$. $\Box$
  % \end{proof}
  %
  % % Lemma 102: Sin prueba. Lemma 14.
  % \begin{lemma}
  %   \PN Si $F: (L, \SU, \IN, 0, 1) \rightarrow (L^{\prime}, \SU^{\prime}, \IN^{\prime},
  %   0^{\prime}, 1^{\prime})$ es un homomorfismo de reticulados acotados, entonces $\ker F$ es una congruencia sobre $(L,
  %   \SU, \IN, 0, 1)$.
  % \end{lemma}
  %
  % % Lemma 103: Sin prueba. Lemma 15.
  % \begin{lemma}
  %   \PN Sea $(L, \SU, \IN, 0, 1)$ un reticulado acotado y $\theta$ una congruencia sobre $(L, \SU,
  %   \IN, 0, 1)$, entonces:
  %   \begin{enumerate}[a)]
  %     \item $(L/\theta, \mathsf{\tilde{s}}, \mathsf{\tilde{\imath}}, 0/\theta, 1/\theta)$ es un reticulado acotado.
  %     \item $\pi_{\theta}$ es un homomorfismo de $(L, \SU, \IN, 0, 1)$ en $(L/\theta, \mathsf{\tilde{s}},
  %       \mathsf{\tilde{\imath}}, 0/\theta, 1/\theta)$ cuyo núcleo es $\theta$.
  %   \end{enumerate}
  % \end{lemma}
  %
  % % Lemma 104: Sin prueba. Lemma 16.
  % \begin{lemma}
  %   \PN Si $F:(L, \SU, \IN, ^{c}, 0, 1) \rightarrow (L^{\prime}, \SU^{\prime}, \IN^{\prime},
  %   ^{c^{\prime}}, 0^{\prime}, 1^{\prime})$ un homomorfismo biyectivo, entonces $F$ es un isomorfismo.
  % \end{lemma}
  %
  % % Lemma 105: Sin prueba. Lemma 17.
  % \begin{lemma}
  %   \PN Si $F: (L, \SU, \IN, ^{c}, 0, 1) \rightarrow (L^{\prime}, \SU^{\prime },
  %   \IN^{\prime}, ^{c^{\prime}}, 0^{\prime}, 1^{\prime})$ es un homomorfismo, entonces $I_{F}$ es un subuniverso
  %   de $(L^{\prime}, \SU^{\prime}, \IN^{\prime}, ^{c^{\prime}}, 0^{\prime}, 1^{\prime})$.
  % \end{lemma}
  %
  % % Lemma 106: Sin prueba. Lemma 18.
  % \begin{lemma}
  %   \PN Si $F: (L, \SU, \IN, ^{c}, 0, 1) \rightarrow (L^{\prime}, \SU^{\prime},
  %   \IN^{\prime}, ^{c^{\prime}}, 0^{\prime}, 1^{\prime})$ es un homomorfismo de reticulados complementados,
  %   entonces $\ker F$ es una congruencia sobre $(L, \SU, \IN, ^{c}, 0, 1)$.
  % \end{lemma}
  %
  % % Lemma 107: Sin prueba. Lemma 19.
  % \begin{lemma}
  %   \PN Sea $(L, \SU, \IN, ^{c}, 0, 1)$ un reticulado complementado y sea $\theta$ una congruencia sobre
  %   $(L, \SU, \IN, ^{c}, 0, 1)$.
  %   \begin{enumerate}[a)]
  %     \item $(L/\theta, \mathsf{\tilde{s}}, \mathsf{\tilde{\imath}}, ^{\tilde{c}}, 0/\theta, 1/\theta)$ es un reticulado
  %       complementado.
  %     \item $\pi_{\theta}$ es un homomorfismo de $(L, \SU, \IN, ^{c}, 0, 1)$ en $(L/\theta,
  %       \mathsf{\tilde{s}}, \mathsf{\tilde{\imath}}, ^{\tilde{c}}, 0/\theta, 1/\theta)$ cuyo núcleo es $\theta$.
  %   \end{enumerate}
  % \end{lemma}
  %
  % % Lemma 108: Con prueba. Lemma 20.
  % \begin{lemma}
  %   Sea $(L,\SU,\IN)$ un reticulado. Son equivalentes:
  %   (1) $x\mathsf{\;i\;}(y\;\SU\;z)=(x\mathsf{\;i\;}y)\;\SU \;(x\mathsf{\;i\;}z)$, cualesquiera sean $x,y,z\in L$
  %   (2) $x\;\SU\;(y\mathsf{\;i\;}z)=(x\mathsf{\;s\;}y)\mathsf{\;i\; }(x\;\SU\;z)$, cualesquiera sean $x,y,z\in L$.
  % \end{lemma}
  % \begin{proof}
  %   (1)$\Rightarrow $(2). Notese que
  %
  %   $\displaystyle \begin{array}{rcl} (x\mathsf{\;s\;}y)\mathsf{\;i\;}(x\;\SU\;z) & =& ((x\mathsf{\;s\;}y) \mathsf{\;i\;}x)\;\SU\;((x\mathsf{\;s\;}y)\mathsf{\;i\;}z) \\ & =& (x\;\SU\;(z\mathsf{\;i\;}(x\mathsf{\;s\;}y)) \\ & =& (x\;\SU\;((z\;\mathsf{i\;}x)\mathsf{\;s\;}(z\;\mathsf{i\;}y)) \\ & =& (x\;\SU\;(z\;\mathsf{i\;}x))\mathsf{\;s\;}(z\;\mathsf{i\;}y) \\ & =& x\mathsf{\;s\;}(z\;\mathsf{i\;}y) \\ & =& x\mathsf{\ s\ }(y\ \mathsf{i\ }z) \end{array} $
  %
  %   (2)$\Rightarrow $(1) es similar. $\Box$
  % \end{proof}
  %
  % % Lemma 109: Con prueba. Lemma 21.
  % \begin{lemma}
  %   Si $(L,\SU,\IN,0,1)$ un reticulado acotado y distributivo, entonces todo elemento tiene a lo sumo un complemento.
  % \end{lemma}
  % \begin{proof}
  %   Supongamos $x\in L$ tiene complementos $y,z$. Se tiene
  %
  %   $\displaystyle y=y\;\mathsf{i\;}1=y\;\mathsf{i\;}(x\ \SU \z)=(y\;\mathsf{i\;}x)\; \mathsf{s\;}(y\;\mathsf{i\;}z)=0\ \SU \(y\;\mathsf{i\;}z)=y\;\mathsf{ i\;}z, $
  %
  %   por lo cual $y\leq z$. En forma analoga se muestra que $z\leq y$. $\Box$
  % \end{proof}
  %
  % % Lemma 110: Con prueba. Lemma 22.
  % \begin{lemma}
  %   Si $S$ es no vacio, entonces $[S)$ es un filtro. Mas aun si $F$ es un filtro y $F\supseteq S$, entonces $F\supseteq \lbrack S)$.
  % \end{lemma}
  % \begin{proof}
  %   Ya que $S\subseteq \lbrack S)$, tenemos que $[S)\neq \varnothing $. Claramente $[S)$ cumple la propiedad (3). Veamos cumple la (2). Si $y\geq s_{1}\; \mathsf{i\;}s_{2}\;\mathsf{i\;}\dotsc\;\mathsf{i\;}s_{n}$ y $z\geq t_{1}\; \mathsf{i\;}t_{2}\;\mathsf{i\;}$\dotsc$\;\mathsf{i\;}t_{m}$, con $ s_{1},s_{2}, \dotsc, s_{n}$, $t_{1},t_{2}, \dotsc, t_{m}\in S$, entonces
  %
  %   $\displaystyle y\;\mathsf{i\;}z\geq s_{1}\;\mathsf{i\;}s_{2}\;\mathsf{i\;}\dotsc\;\mathsf{i\;} s_{n}\;\mathsf{i\;}t_{1}\;\mathsf{i\;}t_{2}\;\mathsf{i\;}\dotsc\;\mathsf{i\;} t_{m}, $
  %
  %   lo cual prueba (2).
  % \end{proof}
  %
  % % Lemma 111: Sin prueba. Lemma 23.
  % \begin{lemma}
  %   \PN (\textbf{Zorn}) Sea $(P, \leq)$ un poset y supongamos que cada cadena de $P$ tiene una cota superior, entonces
  %   existe un elemento maximal en $P$. Un filtro $F$ de un reticulado $(L, \SU, \IN)$ será llamado primo
  %   cuando se cumplan:
  %   \begin{enumerate}
  %     \item $F \neq L$
  %     \item $x \ \SU \ y \in F \Rightarrow x \in F$ ó $y \in F$.
  %   \end{enumerate}
  % \end{lemma}
  %
  % % Theorem 112: Con prueba. Theorem 24.
  % \begin{theorem}
  %   (Teorema del Filtro Primo) Sea $(L,\SU,\IN)$ un reticulado distributivo y $F$ un filtro. Supongamos $x_{0}\in L-F$. Entonces hay un filtro primo $P$ tal que $x_{0}\notin P$ y $F\subseteq P$.
  % \end{theorem}
  % \begin{proof}
  %   Sea
  %
  %   $\displaystyle \mathcal{F}=\{F_{1}:F_{1}\text{ es un filtro, }x_{0}\notin F_{1}\text{ y } F\subseteq F_{1}\}. $
  %
  %   Notese que $\mathcal{F}\neq \varnothing $, por lo cual $(\mathcal{F},\subseteq )$ es un poset. Veamos que cada cadena en $(\mathcal{F},\subseteq )$ tiene una cota superior. Sea $C$ una cadena. Si $C=\varnothing $, entonces cualquier elemento de $\mathcal{F}$ es cota de $C$. Supongamos entonces $C\neq \varnothing $. Sea
  %   $\displaystyle G=\{x\in L:x\in F_{1},\text{para algun }F_{1}\in C\}. $
  %
  %   Veamos que $G$ es un filtro. Es claro que $G$ es no vacio. Supongamos que $ x,y\in G$. Sean $F_{1},F_{2}\in \mathcal{F}$ tales que $x\in F_{1}$ y $y\in F_{2}$. Si $F_{1}\subseteq F_{2}$, entonces ya que $F_{2}$ es un filtro tenemos que $x\;\mathsf{i\;}y\in F_{2}\subseteq G$. Si $F_{2}\subseteq F_{1}$ , entonces tenemos que $x\;\mathsf{i\;}y\in F_{1}\subseteq G$. Ya que $C$ es una cadena, tenemos que siempre $x\;\mathsf{i\;}y\in G$. En forma analoga se prueba la propiedad restante por lo cual tenemos que $G$ es un filtro. Ademas $x_{0}\notin G$, por lo que $G\in \mathcal{F}$ es cota superior de $C$ . Por el lema de Zorn, $(\mathcal{F},\subseteq )$ tiene un elemento maximal $ P$. Veamos que $P$ es un filtro primo. Supongamos $x\ \SU \y\in P$ y $ x,y\notin P$. Entonces ya que $P$ es maximal tenemos que
  %   $\displaystyle x_{0}\in \lbrack P\cup \{x\})\cap \lbrack P\cup \{y\}) $
  %
  %   Ya que $x_{0}\in \lbrack P\cup \{x\})$, tenemos que hay elementos $ p_{1}, \dotsc, p_{n}\in P$, tales que
  %   $\displaystyle x_{0}\geq p_{1}\;\mathsf{i\;}\dotsc\;\mathsf{i\;}p_{n}\;\mathsf{i\;}x $
  %
  %   Ya que $x_{0}\in \lbrack P\cup \{y\})$, tenemos que hay elementos $ q_{1}, \dotsc, q_{m}\in P$, tales que
  %   $\displaystyle x_{0}\geq q_{1}\;\mathsf{i\;}\dotsc\;\mathsf{i\;}q_{m}\;\mathsf{i\;}y $
  %
  %   Si llamamos $p$ al siguiente elemento de $P$
  %   $\displaystyle p_{1}\;\mathsf{i\;}\dotsc\;\mathsf{i\;}p_{n}\;\mathsf{i\;}q_{1}\;\mathsf{i\;} \dotsc\;\mathsf{i\;}q_{m} $
  %
  %   tenemos que
  %   $\displaystyle \begin{array}{rcl} x_{0} & \geq & p\;\mathsf{i\;}x \\ x_{0} & \geq & p\;\mathsf{i\;}y \end{array} $
  %
  %   Se tiene que $x_{0}\geq (p\;\mathsf{i\;}x)\ \SU \(p\;\mathsf{i\;} y)=p\;\mathsf{i\;}(x\ \SU \y)\in P$, lo cual es absurdo ya que $ x_{0}\notin P$. $\Box$
  % \end{proof}
  %
  % % Corollary 113: Con prueba. Corollary 25.
  % \begin{corollary}
  %   Sea $(L,\SU,\IN,0,1)$ un reticulado acotado distributivo. Si $ \varnothing \neq S\subseteq L$ es tal que $s_{1}\;\mathsf{i\;}s_{2}\;\mathsf{ i\;}\dotsc\;\mathsf{i\;}s_{n}\neq 0$, para cada $s_{1}, \dotsc, s_{n}\in S$, entonces hay un filtro primo que contiene a $S$.
  % \end{corollary}
  % \begin{proof}
  %   Notese que $[S)\neq L$ por lo cual se puede aplicar el Teorema del filtro primo. $\Box$
  % \end{proof}
  %
  % % Lemma 114: Con prueba. Lemma 26.
  % \begin{lemma}
  %   Sea $(B,\SU,\IN,^{c},0,1)$ un algebra de Boole. Entonces para un filtro $F\subseteq B$ las siguientes son equivalentes:
  %   (1) $F$ es primo
  %   (2) $x\in F$ o $x^{c}\in F$, para cada $x\in B$.
  % \end{lemma}
  % \begin{proof}
  %   (1)$\Rightarrow $(2). Ya que $x\ \SU \x^{c}=1\in F$, (2) se cumple si $F$ es primo.
  %
  %   (2)$\Rightarrow $(1). Supongamos que $x\ \SU \y\in F$ y que $x\not\in F$. Entonces por (2), $x^{c}\in F$ y por lo tanto tenemos que
  %
  %   $\displaystyle y\geq x^{c}\;\mathsf{i\;}y=(x^{c}\;\mathsf{i\;}x)\ \SU \(x^{c}\; \mathsf{i\;}y)=x^{c}\;\mathsf{i\;}(x\ \SU \y)\in F, $
  %
  %   lo cual dice que $y\in F$. $\Box$
  % \end{proof}
  %
  % % Lemma 115: Con prueba. Lemma 27.
  % \begin{lemma}
  %   Sea $(B,\SU,\IN,^{c},0,1)$ un algebra de Boole. Supongamos que $b\neq 0$ y $a=\inf A$, con $A\subseteq B$. Entonces si $b\;\mathsf{i\;}a=0$ , existe un $e\in A$ tal que $b\;\mathsf{i\;}e^{c}\neq 0$.
  % \end{lemma}
  % \begin{proof}
  %   Supongamos que para cada $e\in A$, tengamos que $b\;\mathsf{i\;}e^{c}=0$. Entonces tenemos que para cada $e\in A$,
  %
  %   $\displaystyle b=b\;\mathsf{i\;}(e\ \SU \e^{c})=(b\;\mathsf{i\;}e)\ \SU \(b\; \mathsf{i\;}e^{c})=b\;\mathsf{i\;}e, $
  %
  %   lo cual nos dice que $b$ es cota inferior de $A$. Pero entonces $b\leq a$, por lo cual $b=b\;\mathsf{i\;}a=0$, contradiciendo la hipotesis. $\Box$
  % \end{proof}
  %
  % % Theorem 116: Con prueba. Theorem 28.
  % \begin{theorem}
  %   (Rasiova y Sikorski) Sea $(B,\SU,\IN,^{c},0,1)$ un algebra de Boole. Sea $x\in B$, $x\neq 0$. Supongamos que $A_{1},A_{2},\dotsc$ son subconjuntos de $B$ tales que existe $\inf (A_{j})$, para cada $j=1,2\dotsc.$ Entonces hay un filtro primo $P$ el cual cumple:
  %   (a) $x\in P$
  %   (b) $P\supseteq A_{j}\Rightarrow P\ni \inf (A_{j})$, para cada $ j=1,2,\dotsc.$
  % \end{theorem}
  % \begin{proof}
  %   Sean $a_{j}=\inf (A_{j})$, $j=1,2,\dotsc$. Construiremos inductivamente una sucesion $b_{0},b_{1},\dotsc$ de elementos de $B$ tal que:
  %
  %   (1) $b_{0}=x$
  %   (2) $b_{0}\;\mathsf{i\;}$\dotsc$\;\mathsf{i\;}b_{n}\neq 0$, para cada $ n\geq 0$
  %   (3) $b_{j}=a_{j}$ o $b_{j}^{c}\in A_{j}$, para cada $j\geq 1$.
  %   Definamos $b_{0}=x$. Supongamos ya definimos $b_{0}, \dotsc, b_{n}$, veamos como definir $b_{n+1}$. Si $(b_{0}\;\mathsf{i\;}\dotsc\;\mathsf{i\;} b_{n})\;\mathsf{i\;}a_{n+1}\neq 0$, entonces definamos $b_{n+1}=a_{n+1}$. Si $(b_{0}\;\mathsf{i\;}\dotsc\;\mathsf{i\;}b_{n})\;\mathsf{i\;}a_{n+1}=0$, entonces por el lema anterior, tenemos que hay un $e\in A_{n+1}$ tal que $ (b_{0}\;\mathsf{i\;}\dotsc\;\mathsf{i\;}b_{n})\;\mathsf{i\;}e^{c}\neq 0$, lo cual nos permite definir $b_{n+1}=e^{c}$.
  %
  %   Usando (2) se puede probar que el conjunto $S=\{b_{0},b_{1},\dotsc\}$ satisface la hipotesis del primer corolario del Teorema del filtro primo, por lo cual hay un filtro primo $P$ tal que $\{b_{0},b_{1},\dotsc\}\subseteq P$. Es facil chequear que $P$ satisface las propiedades (a) y (b). $\Box$
  % \end{proof}

	\section{Términos y fórmulas}

  % Lemma 117. Con prueba. Lemma 29.
  \begin{lemma}
    Supongamos $t\in T_{k}^{\tau }$, con $k\geq 1$. Entonces ya sea $t\in Var \cup \mathcal{C}$ o $t=f(t_{1}, \dotsc,t_{n})$, con $f\in \mathcal{F }_{n}$, $n\geq 1,\;t_{1}, \dotsc,t_{n}\in T_{k-1}^{\tau }$.
  \end{lemma}
  \begin{proof}
    Por induccion en $k$.

    CASO $k=1$: Es directo ya que por definicion

    $\displaystyle T_{1}^{\tau }=Var \cup \mathcal{C}\cup \{f(t_{1}, \dotsc,t_{n}):f\in \mathcal{F}_{n},n\geq 1,t_{1}, \dotsc,t_{n}\in T_{0}^{\tau }\}. $

    CASO $k\Rightarrow k+1$: Sea $t\in T_{k+1}^{\tau }$. Por definicion de $ T_{k+1}^{\tau }$ tenemos que $t\in T_{k}^{\tau }$ o $t=f(t_{1}, \dotsc,t_{n})$ con $f\in \mathcal{F}_{n}$, $n\geq 1$ y $t_{1}, \dotsc,t_{n}\in T_{k}^{\tau }$. Si se da que $t\in T_{k}^{\tau }$, entonces podemos aplicar hipotesis inductiva y usar que $T_{k-1}^{\tau }\subseteq T_{k}^{\tau }$. Esto completa el caso. $\Box$
  \end{proof}

  % Lemma 118. Con prueba. Lemma 30.
  \begin{lemma}
    \PN Sea $b \in Bal$. Se tiene:
    \begin{enumerate}[(1)]
      \item $\lvert b\rvert_{(}-\lvert b\rvert_{)}=0$
      \item Si $x$ es tramo inicial propio de $b$, entonces $\lvert x\rvert_{(}-\lvert x\rvert_{)} >0$
      \item Si $x$ es tramo final propio de $b$, entonces $\lvert x\rvert_{(}-\lvert x\rvert_{)}< 0$
    \end{enumerate}
  \end{lemma}
  \begin{proof}
    Probaremos por induccion en $k$, que valen (1), (2) y (3) para cada $b\in Bal_{k}$. El caso $k=1$ es trivial. Supongamos $b\in Bal_{k+1}$. Si $b\in Bal_{k}$, se aplica directamente HI. Supongamos entonces que $ b=(b_{1}...b_{n})$, con $b_{1}, \dotsc,b_{n}\in Bal_{k}$, $n\geq 1$. Por HI, $ b_{1}, \dotsc,b_{n}$ cumplen (1) por lo cual $b$ cumple (1). Veamos que $b$ cumple (2). Sea $x$ un tramo inicial propio de $b$. Notese que $x$ es de la forma $x=(b_{1}...b_{i}x_{1}$ con $0\leq i\leq n-1$ y $x_{1}$ un tramo inicial de $b_{i+1}$ (en el caso $i=0$ interpretamos $b_{1}...b_{i}= \varepsilon )$. Pero entonces ya que

    $\displaystyle \lvert x\rvert_{(}-\lvert x\rvert_{)}=1+\left( \sum_{j=1}^{i}\lvert b_{j}\rvert_{(}-\lvert b_{j}\rvert_{)}\right) +\lvert x_{1}\rvert_{(}-\lvert x_{1}\rvert_{)} $

    tenemos que por HI, se da que $\lvert x\rvert_{(}-\lvert x\rvert_{)} >0$. En forma analoga se puede ver que $b$ cumple (3). $\Box$
  \end{proof}

  % Lemma 119. Sin prueba. Lemma 31.
  \begin{lemma}
    \PN $del(xy) = del(x)del(y) \ \forall x, y \in \SIGMA$
  \end{lemma}

  % Lemma 120. Sin prueba. Lemma 32.
  \begin{lemma}
    \PN Supongamos que $\Sigma$ es tal que $\TAU \subseteq \SIGMA$, entonces $del(t) \in Bal$, para cada $t \in \TAU -
    (Var \cup \mathcal{C})$.
  \end{lemma}

  % Lemma 121. Con prueba. Lemma 33.
  \begin{lemma}
    \PN Sean $s, t \in \TAU$ y supongamos que hay palabras $x,y,z$, con $y \neq \varepsilon$ tales que $s=xy$ y $t=yz$,
    entonces $x=z=\varepsilon$ ó $s, t \in \mathcal{C}$. En particular si un término es tramo inicial o final de otro
    término, entonces dichos terminos son iguales.
  \end{lemma}
  \begin{proof}
    Supongamos $s\in \mathcal{C}$. Ya que $y\neq \varepsilon $ tenemos que $t$ debe comenzar con un simbolo que ocurre en un nombre de cte, lo cual dice que $t$ no puede ser ni una variable ni de la forma $g(t_{1}, \dotsc,t_{m})$, es decir $t\in \mathcal{C}$. Supongamos $s\in Var$. Si $x\neq \varepsilon $ tenemos que $t$ debe comenzar con alguno de los siguientes simbolos

    $\displaystyle \mathit{0}\;\mathit{1\;}...\;\mathit{9}\;\mathbf{0}\;\mathbf{1}\ ...\; \mathbf{9} $

    lo cual es absurdo. O sea que $x=\varepsilon $ y por lo tanto $t$ debe comenzar con $\mathsf{X}$. Pero esto dice que $t\in Var$ de lo que sigue facilmente que $z=\varepsilon $. Supongamos entonces que $s$ es de la forma $ f(s_{1}, \dotsc,s_{n})$. Ya que $)$ debe ocurrir en $t$, tenemos que $t$ es de la forma $g(t_{1}, \dotsc,t_{m})$. O sea que $del(s),del(t)\in Bal$. Ya que $)$ ocurre en $y$, $del(y)\neq \varepsilon $. Tenemos tambien que
    $\displaystyle \begin{array}{rcl} del(s) & =& del(x)del(y) \\ del(t) & =& del(y)del(z) \end{array} $

    La primera igualdad, por (3) del Lema 118, nos dice que
    $\displaystyle \lvert del(y)\rvert_{(}-\lvert del(y)\rvert_{)}\leq 0, $

    y la segunda que
    $\displaystyle \lvert del(y)\rvert_{(}-\lvert del(y)\rvert_{)}\geq 0, $

    por lo cual
    $\displaystyle \lvert del(y)\rvert_{(}-\lvert del(y)\rvert_{)}=0 $

    Pero entonces ya que $del(y)$ es tramo final de $del(s)$, (3) del Lema 118 nos dice que $del(x)=\varepsilon $. Similarmente obtenemos que $del(z)=\varepsilon $. Ya que que $t$ termina con $)$ tenemos que $z=\varepsilon $. O sea que $f(s_{1}, \dotsc,s_{n})=xg(t_{1}, \dotsc,t_{m})$ con $del(x)=\varepsilon $, de lo que se saca que $f=xg$ ya que $($ no ocurre en $ x$. De la definicion de tipo se desprende que $x=\varepsilon $. $\Box$
  \end{proof}

  % Theorem 122. Con prueba. Lemma 34.
  \begin{theorem}
    \PN (Lectura unica de terminos). Dado $t \in \TAU$ se da una de las siguientes:
    \begin{enumerate}[1)]
      \item $t \in Var \cup \mathcal{C}$
      \item Hay únicos $n \geq 1, \ f \in \mathcal{F}_{n}, \ t_{1}, \dotsc, t_{n} \in \TAU$ tales que $t = f(t_{1},
      \dotsc, t_{n})$.
    \end{enumerate}
  \end{theorem}
  \begin{proof}
    En virtud del Lema 117 solo nos falta probar la unicidad en el punto (2). Supongamos que

    $\displaystyle t=f(t_{1}, \dotsc,t_{n})=g(s_{1}, \dotsc,s_{m}) $

    con $n,m\geq 1,\;f\in \mathcal{F}_{n}$, $g\in \mathcal{F}_{m}$, $ t_{1}, \dotsc,t_{n},s_{1}, \dotsc,s_{m}\in \TAU$. Notese que $f=g$. O sea que $ n=m=a(f)$. Notese que $t_{1}$ es tramo inicial de $s_{1}$ o $s_{1}$ es tramo inicial de $t_{1}$, lo cual por el lema anterior nos dice que $t_{1}=s_{1}$. Con el mismo razonamiento podemos probar que debera suceder $ t_{2}=s_{2}, \dotsc,t_{n}=s_{n}$. $\Box$
  \end{proof}

  % Lemma 123. Con prueba. Lemma 35.
  \begin{lemma}
    \PN Sean $r, s, t \in \TAU$.
    \begin{enumerate}[(a)]
      \item Si $s \neq t = f(t_{1}, \dotsc, t_{n})$ y $s$ ocurre en $t$, entonces dicha ocurrencia sucede dentro de
        algún $t_{j}$, $j = 1, \dotsc, n$.
      \item Si $r, s$ ocurren en $t$, entonces dichas ocurrencias son disjuntas o una ocurre dentro de otra. En
        particular, las distintas ocurrencias de $r$ en $t$ son disjuntas.
      \item Si $t^{\prime}$ es el resultado de reemplazar una ocurrencia de $s$ en $t$ por $r$, entonces $t^{\prime} \in
        \TAU$.
    \end{enumerate}
  \end{lemma}
  \begin{proof}
    (a) Supongamos la ocurrencia de $s$ comienza en algun $t_{j}$. Entonces el Lema 121 nos conduce a que dicha ocurrencia debera estar contenida en $t_{j}$. Veamos que la ocurrencia de $s$ no puede ser a partir de un $i\in \{1, \dotsc,\lvert f\rvert \}$. Supongamos lo contrario. Tenemos entonces que $s$ debe ser de la forma $g(s_{1}, \dotsc,s_{m})$ ya que no puede estar en $Var \cup \mathcal{C}$. Notese que $i\neq 1$ ya que en caso contrario $s$ seria un tramo inicial propio de $t$. Pero entonces $g$ debe ser un tramo final propio de $f$, lo cual es absurdo. Ya que $s$ no puede comenzar con parentesis o coma, hemos contemplado todos los posibles casos de comienzo de la ocurrencia de $s$ en $t$.

    (b) y (c) pueden probarse por induccion, usando (a). $\Box$
  \end{proof}

  % Lemma 124. Con prueba. Lemma 36.
  \begin{lemma}
    \PN Supongamos $\varphi \in F_{k}^{\tau}$, con $k \geq 1$, entonces $\varphi$ es de alguna de las siguientes formas:
    \begin{itemize}
      \item $\varphi = (t \equiv s)$, con $t, s \in \TAU$.
      \item $\varphi = r(t_{1}, \dotsc, t_{n})$, con $r \in \mathcal{R}_{n}, t_{1}, \dotsc, t_{n} \in \TAU$.
      \item $\varphi = (\varphi_{1} \eta \varphi_{2})$, con $\eta \in \{\wedge, \vee, \rightarrow, \leftrightarrow\}, \
        \varphi_{1}, \varphi_{2} \in F_{k-1}^{\tau}$.
      \item $\varphi = \lnot \varphi_{1}$, con $\varphi_{1} \in F_{k-1}^{\tau}$.
      \item $\varphi = Qv \varphi_{1}$, con $Q \in \{\forall, \exists\}, \ v \in Var$ y $\varphi_{1} \in F_{k-1}^{\tau}$.
    \end{itemize}
  \end{lemma}
  \begin{proof}
    Induccion en $k$. $\Box$
  \end{proof}

  % Lemma 125. Con prueba. Lemma 37.
  \begin{lemma}
    \PN Sea $\tau$ un tipo.
    \begin{enumerate}[a)]
      \item Supongamos que $\Sigma$ es tal que $FT \subseteq \SIGMA$, entonces $del(\varphi) \in Bal$, para cada
        $\varphi \in FT$.
      \item Sea $\varphi \in F_{k}^{\tau}$, con $k \geq 0$, existen $x \in (\{\lnot\} \cup \{Qv: Q \in \{\forall,
        \exists\}$ y $v \in Var\})^{\ast}$ y $\varphi_{1} \in FT$ tales que $\varphi = x \varphi_{1}$ y
        $\varphi_{1}$ es de la forma $(\psi_{1} \eta \psi_{2})$ o atómica. En particular toda fórmula termina con el
        símbolo $)$.
    \end{enumerate}
  \end{lemma}
  \begin{proof}
    (b) Induccion en $k$. El caso $k=0$ es trivial. Supongamos (b) vale para cada $\varphi \in F_{k}^{\tau }$ y sea $\varphi \in F_{k+1}^{\tau }$. Hay varios casos de los cuales haremos solo dos

    CASO $\varphi =(\varphi_{1}\eta \varphi_{2})$, con $\varphi_{1},\varphi_{2}\in F_{k}^{\tau }$ y $\eta \in \{\vee ,\wedge ,\rightarrow ,\leftrightarrow \}$.

    Podemos tomar $x=\varepsilon $ y $\varphi_{1}=\varphi $.

    CASO $\varphi =Qx_{i}\psi $, con $\psi \in F_{k}^{\tau }$, $i\geq 1$ y $Q\in \{\forall ,\exists \}$.

    Por HI hay $\bar{x}\in (\{\lnot \}\cup \{Qv:Q\in \{\forall ,\exists \}$ y $v\in Var\})^{\ast }$ y $\psi_{1}\in \FT$ tales que $ \psi =x\psi_{1}$ y $\psi_{1}$ es de la forma $(\gamma_{1}\eta \gamma_{2}) $ o atomica. Entonces es claro que $x=Qx_{i}\bar{x}$ y $\varphi_{1}=\psi_{1}$ cumplen (b). $\Box$
  \end{proof}

  % Lemma 126. Con prueba. Lemma 38.
  \begin{lemma}
    \PN Ninguna fórmula es tramo final propio de una fórmula atómica, es decir, si $\varphi = x \psi$, con $\varphi \in
    F_{0}^{\tau}$ y $\psi \in FT$, entonces $x = \varepsilon$.
  \end{lemma}
  \begin{proof}
    Si $\varphi $ es de la forma $(t\equiv s)$, entonces $\lvert del(y)\rvert_{(}-\lvert del(y)\rvert_{)}< 0$ para cada tramo final propio $y$ de $\varphi $, lo cual termina el caso ya que $del(\psi )$ es balanceada. Supongamos entonces $\varphi =r(t_{1}, \dotsc,t_{n})$. Notese que $\psi $ no puede ser tramo final de $t_{1}, \dotsc,t_{n})$ ya que $del(\psi )$ es balanceada y $\lvert del(y)\rvert_{(}-\lvert del(y)\rvert_{)}< 0$ para cada tramo final $y$ de $t_{1}, \dotsc,t_{n})$. Es decir que $\psi =y(t_{1}, \dotsc,t_{n})$, para algun tramo final $y$ de $r$. Ya que en $\psi $ no ocurren cuantificadores ni nexos ni el simbolo $\equiv $ el Lema 124 nos dice $\psi =\tilde{r}(s_{1}, \dotsc,s_{m})$, con $ \tilde{r}\in \mathcal{R}_{m}$, $m\geq 1$ y $s_{1}, \dotsc,s_{m}\in \TAU$. Ahora es facil usando un argumento paresido al usado en la prueba del Teorema 122 concluir que $m=n$, $s_{i}=t_{i}$, $ i=1, \dotsc,n$ y $\tilde{r}$ es tramo final de $r$. Por (3) de la definicion de tipo tenemos que $\tilde{r}=r$ lo cual nos dice que $\varphi =\psi $ y $ x=\varepsilon $ $\Box$
  \end{proof}

  % Lemma 127. Con prueba. Lemma 39.
  \begin{lemma}
    \PN Si $\varphi = x \psi$, con $\varphi, \psi \in FT$ y $x$ sin paréntesis, entonces $x \in (\{\lnot\} \cup
    \{Qv: Q \in \{\forall, \exists\}$ y $v \in Var\})^{\ast}$.
  \end{lemma}
  \begin{proof}
    Por induccion en el $k$ tal que $\varphi \in F_{k}^{\tau }$. El caso $k=0$ es probado en el lema anterior. Asumamos que el resultado vale cuando $ \varphi \in F_{k}^{\tau }$ y veamos que vale cuando $\varphi \in F_{k+1}^{\tau }$. Mas aun supongamos $\varphi \in F_{k+1}^{\tau }-F_{k}^{\tau }$. Primero haremos el caso en que $\varphi =Qv\varphi_{1}$, con $Q\in \{\forall ,\exists \},\;v\in Var$ y $\varphi_{1}\in F_{k}^{\tau }$ . Supongamos $x\neq \varepsilon $. Ya que $\psi $ no comienza con simbolos de $v$, tenemos que $\psi $ debe ser tramo final de $\varphi_{1}$ lo cual nos dice que hay una palabra $x_{1}$ tal que $x=Qvx_{1}$ y $\varphi_{1}=x_{1}\psi $. Por HI tenemos que $x_{1}\in (\{\lnot \}\cup \{Qv:Q\in \{\forall ,\exists \}$ y $v\in Var\})^{\ast }$ con lo cual $x\in (\{\lnot \}\cup \{Qv:Q\in \{\forall ,\exists \}$ y $v\in Var\})^{\ast }$. El caso en el que $\varphi =\lnot \varphi_{1}$ con $\varphi_{1}\in F_{k}^{\tau }$, es similar. Note que no hay mas casos posibles ya que $\varphi $ no puede comenzar con $($ porque en $x$ no ocurren parentesis por hipotesis $\Box$
  \end{proof}

  % Proposition 128. Con prueba. Lemma 40.
  \begin{proposition}
    \PN Si $\varphi, \psi \in FT$ y $x, y, z$ son tales que $\varphi = xy, \psi = yz$ y $y \neq \varepsilon$,
    entonces $z = \varepsilon$ y $x \in (\{\lnot\} \cup \{Qv: Q \in \{\forall, \exists\}$ y $v \in Var\})^{\ast}$. En
    particular ningún tramo inicial propio de una fórmula es una fórmula.
  \end{proposition}
  \begin{proof}
    Ya que $\varphi $ termina con $)$ tenemos que $del(y)\neq \varepsilon .$ Ya que $del(\varphi ),del(\psi )\in Bal$ y ademas

    $\displaystyle \begin{array}{rcl} del(\varphi ) & =& del(x)del(y) \\ del(\psi ) & =& del(y)del(z) \end{array} $

    tenemos que $del(y)$ es tramo inicial y final de palabras balanceadas, lo cual nos dice que
    $\displaystyle \lvert del(y)\rvert_{(}-\lvert del(y)\rvert_{)}=0 $

    Pero esto por (3) del Lema 118 nos dice que $ del(x)=\varepsilon $. Similarmente obtenemos que $del(z)=\varepsilon $. Pero $\psi $ termina con $)$ lo cual nos dice que $z=\varepsilon $. Es decir que $ \varphi =x\psi $. Por el lema anterior tenemos que $x\in (\{\lnot \}\cup \{Qv:Q\in \{\forall ,\exists \}$ y $v\in Var\})^{\ast }$ $\Box$
  \end{proof}

  % Theorem 129. Con prueba. Lemma 41.
  \begin{theorem}
    \PN (Lectura única de fórmulas) Dada $\varphi \in FT$ se da una y solo una de las siguientes:
    \begin{enumerate}[(1)]
      \item $\varphi = (t \equiv s)$, con $t, s \in \TAU$
      \item $\varphi = r(t_{1}, \dotsc, t_{n})$, con $r \in \mathcal{R}_{n}, t_{1}, \dotsc, t_{n} \in \TAU$
      \item $\varphi = (\varphi_{1} \eta \varphi_{2})$, con $\eta \in \{\wedge, \vee, \rightarrow, \leftrightarrow\}, \
        \varphi_{1}, \varphi_{2} \in \FT$
      \item $\varphi = \lnot \varphi_{1}$, con $\varphi_{1} \in \FT$
      \item $\varphi = Qv \varphi_{1}$, con $Q \in \{\forall, \exists\}, \ \varphi_{1} \in \FT$ y $v \in Var$.
    \end{enumerate}

    \PN Más aún, en todos los puntos tales descomposiciones son únicas.
  \end{theorem}
  \begin{proof}
    Si una formula $\varphi $ satisface (1), entonces $\varphi $ no puede contener simbolos del alfabeto $\{\wedge ,\vee,
    \rightarrow ,\leftrightarrow \}$ lo cual garantiza que $\varphi $ no puede satisfacer (3). Ademas $ \varphi $ no
    puede satisfacer (2) o (4) o (5) ya que $\varphi $ comienza con $($. En forma analoga se puede terminar de ver que
    las propiedades (1), $\dotsc$,(5) son excluyentes.

    La unicidad en las descomposiciones de (4) y (5) es obvia. La de (3) se desprende facilmente del lema anterior y la de los puntos (1) y (2) del lema analogo para terminos. $\Box$
  \end{proof}

  % Lemma 130. Con prueba. Lemma 42.
  \begin{lemma}
    \PN Sea $\tau$ un tipo.
    \begin{enumerate}[(a)]
      \item Las fórmulas atómicas no tienen subfórmulas propias.
      \item Si $\varphi$ ocurre propiamente en $(\psi \eta \varphi)$, entonces tal ocurrencia es en $\psi$ ó en
        $\varphi$.
      \item Si $\varphi$ ocurre propiamente en $\lnot \psi$, entonces tal ocurrencia es en $\psi$.
      \item Si $\varphi$ ocurre propiamente en $Qx_{k} \psi$, entonces tal ocurrencia es en $\psi$.
      \item Si $\varphi_{1}, \varphi_{2}$ ocurren en $\varphi$, entonces dichas ocurrencias son disjuntas o una contiene
        a la otra.
      \item Si $\lambda^{\prime}$ es el resultado de reemplazar alguna ocurrencia de $\varphi$ en $\lambda$ por $\psi$,
        entonces $\lambda^{\prime} \in \FT$.
    \end{enumerate}
  \end{lemma}
  \begin{proof}
    Ejercicio.
  \end{proof}

	\section{Estructuras}

  % Lemma 133. Con prueba. Lemma 43.
  \begin{lemma} \label{lemma_43}
    \PN Sea $\mathbf{A}$ una estructura de tipo $\tau$ y sea $t\in \TAU$. Supongamos que $\vec{a}, \vec{b}$ son
    asignaciones tales que $a_{i} = b_{i}$ cada vez que $x_{i}$ ocurra en $t$, entonces $t^{\mathbf{A}}[\vec{a}] =
    t^{\mathbf{A}}[\vec{b}]$.
  \end{lemma}
  \begin{proof}
    Sea

    - Teo$_{k}$: El lema vale para $t\in T_{k}^{\tau }$.
    Teo$_{0}$ es facil de probar. Veamos Teo$_{k}\Rightarrow $Teo$_{k+1}$. Supongamos $t\in T_{k+1}^{\tau }-T_{k}^{\tau }$
    y sean $\vec{a},\vec{b}$ asignaciones tales que $a_{i}=b_{i}$ cada vez que $x_{i}$ ocurra en $t$. Notese que
    $t=f(t_{1}, \dotsc, t_{n})$, con $f\in \mathcal{F}_{n},\;n\geq 1 \ \text{y} \  t_{1}, \dotsc, t_{n}\in \TAU$. Notese que para
    cada $j=1, \dotsc, n$, tenemos que $a_{i}=b_{i}$ cada vez que $x_{i}$ ocurra en $t_{j}$, lo cual por Teo$_{k}$ nos
    dice que

    $t_{j}^{\mathbf{A}}[\vec{a}]=t_{j}^{\mathbf{A}}[\vec{b}]\text{, }j=1, \dotsc, n $

    Se tiene entonces que
    \[
      \begin{array}{ccl}
        t^{\mathbf{A}}[\vec{a}] &=& i(f)(t_{1}^{\mathbf{A}}[\vec{a}], \dotsc, t_{n}^{ \mathbf{A}}[\vec{a}])\text{ (por def de }t^{\mathbf{A}}[\vec{a}]\text{)} \\
        &=& i(f)(t_{1}^{\mathbf{A}}[\vec{b}], \dotsc, t_{n}^{\mathbf{A}}[\vec{b}]) \\
        &=& t^{\mathbf{A}}[\vec{b}]\text{ (por def de }t^{\mathbf{A}}[\vec{a}] \text{)}
      \end{array}
    \]
  \end{proof}

  % Lemma 134. Con prueba. Lemma 44.
  \begin{lemma} \label{lemma_44}
    \PN \newline
    \begin{enumerate}[(a)]
      \item $Li((t \equiv s)) = \{v \in Var: v$ ocurre en $t$ ó $v$ ocurre en $s\}$
      \item $Li(r(t_{1}, \dotsc, ,t_{n})) = \{v \in Var: v$ ocurre en algún $t_{i}\}$
      \item $Li(\lnot \varphi) = Li(\varphi)$
      \item $Li((\varphi \eta \psi)) = Li(\varphi) \cup Li(\psi)$
      \item $Li(Qx_{j}\varphi) = Li(\varphi)-\{x_{j}\}$
    \end{enumerate}
  \end{lemma}
  \begin{proof}
    \begin{enumerate}[(a)]
      \item son triviales de las definiciones, teniendo en cuenta que si una variable $v$ ocurre en $(t\equiv s)$ (resp. en $r(t_{1}, \dotsc, t_{n})$) entonces $v$ ocurre en $t \ \text{o} \ v$ ocurre en $s$ (resp.$v$ ocurre en algun $ t_{i}$)
      \item son triviales de las definiciones, teniendo en cuenta que si una variable $v$ ocurre en $(t\equiv s)$ (resp. en $r(t_{1}, \dotsc, t_{n})$) entonces $v$ ocurre en $t \ \text{o} \ v$ ocurre en $s$ (resp.$v$ ocurre en algun $ t_{i}$)
      \item es similar a (d)
      \item Supongamos $v\in Li((\varphi \eta \psi ))$, entonces hay un $i$ tal que $ v$ ocurre libremente en $(\varphi \eta \psi )$ a partir de $i$. Por definicion tenemos que ya sea $v$ ocurre libremente en $\varphi $ a partir de $i-1 \ \text{o} \ v$ ocurre libremente en $\psi $ a partir de $i-\left\vert (\varphi \eta \right\vert $, con lo cual $v\in Li(\varphi )\cup Li(\psi )$
      Supongamos ahora que $v\in Li(\varphi )\cup Li(\psi )$. S.p.d.g. supongamos $ v\in Li(\psi )$. Por definicion tenemos que hay un $i$ tal que $v$ ocurre libremente en $\psi $ a partir de $i$. Pero notese que esto nos dice por definicion que $v$ ocurre libremente en $(\varphi \eta \psi )$ a partir de $ i+\left\vert (\varphi \eta \right\vert $ con lo cual $v\in Li((\varphi \eta \psi ))$.
      \item Supongamos $v\in Li(Qx_{j}\varphi )$, entonces hay un $i$ tal que $v$ ocurre libremente en $Qx_{j}\varphi $ a partir de $i$. Por definicion tenemos que $v\neq x_{j} \ \text{y} \ v$ ocurre libremente en $\varphi $ a partir de $ i-\left\vert Qx_{j}\right\vert $, con lo cual $v\in Li(\varphi )-\{x_{j}\}$
      Supongamos ahora que $v\in Li(\varphi )-\{x_{j}\}$. Por definicion tenemos que hay un $i$ tal que $v$ ocurre libremente en $\varphi $ a partir de $i$. Ya que $v\neq x_{j}$ esto nos dice por definicion que $v$ ocurre libremente en $Qx_{j}\varphi $ a partir de $i+\left\vert Qx_{j}\right\vert $, con lo cual $v\in Li(Qx_{j}\varphi )$.
    \end{enumerate}
  \end{proof}

  % Lemma 135. Con prueba. Lemma 45.
  \begin{lemma} \label{lemma_45}
    \PN Supongamos que $\vec{a}, \vec{b}$ son asignaciones tales que si $x_{i} \in Li(\varphi)$, entonces $a_{i} =
    b_{i}$, entonces $\mathbf{A} \models \varphi \lbrack \vec{a}] \Leftrightarrow \mathbf{A} \models \varphi \lbrack
    \vec{b}]$.
  \end{lemma}
  \begin{proof}
    Probaremos por induccion en $k$ que el lema vale para cada $\varphi \in F_{k}^{\tau }.$ El caso $k=0$ se desprende del Lema 131. Veamos que Teo$_{k}$ implica Teo$_{k+1}.$ Sea $\varphi \in F_{k+1}^{\tau }-F_{k}^{\tau }.$ Hay varios casos:

    CASO $\varphi = (\varphi_{1}\wedge \varphi_{2})$.

    Ya que $Li(\varphi_{i})\subseteq Li(\varphi )$, $i=1,2$, Teo$_{k}$ nos dice que $\mathbf{A} \models \varphi_{i}[\vec{a}]$ sii $\mathbf{A} \models \varphi_{i}[\vec{b}]$, para $i=1,2$. Se tiene entonces que

    $\displaystyle \begin{array}{l} \mathbf{A} \models \varphi \lbrack \vec{a}] \\ \ \ \Updownarrow \text{ (por (3) en la def de }\mathbf{A} \models \varphi \lbrack \vec{a}]\text{)} \\ \mathbf{A} \models \varphi_{1}[\vec{a}]\text{ y }\mathbf{A} \models \varphi_{2}[\vec{a}] \\ \ \ \Updownarrow \text{ (por Teo}_{k}\text{)} \\ \mathbf{A} \models \varphi_{1}[\vec{b}]\text{ y }\mathbf{A} \models \varphi_{2}[\vec{b}] \\ \ \ \Updownarrow \text{(por (3) en la def de }\mathbf{A} \models \varphi \lbrack \vec{a}]\text{)} \\ \mathbf{A} \models \varphi \lbrack \vec{b}] \end{array} $

    CASO $\varphi = (\varphi_{1}\wedge \varphi_{2})$.

    Es completamente similar al anterior.

    CASO $\varphi =\lnot \varphi_{1}.$

    Es completamente similar al anterior.

    CASO $\varphi =\forall x_{j}\varphi_{1}.$

    Supongamos $\mathbf{A} \models \varphi \lbrack \vec{a}]$. Entonces por (8) en la def de $\mathbf{A} \models \varphi \lbrack \vec{a}]$ se tiene que $\mathbf{A} \models \varphi_{1}[\downarrow _{j}^{a}(\vec{a})]$, para todo $a\in A$. Notese que $\downarrow _{j}^{a}(\vec{a}) \ \text{y} \ \downarrow _{j}^{a}(\vec{b})$ coinciden en toda $x_{i}$ de $x_{i}\in Li(\varphi_{1})\subseteq Li(\varphi_{1})\cup \{x_{j}\}$, con lo cual por Teo$_{k}$ se tiene que $\mathbf{A} \models \varphi_{1}[\downarrow _{j}^{a}(\vec{b})]$, para todo $a\in A$, lo cual por (8) en la def de $\mathbf{A} \models \varphi \lbrack \vec{a}]$ nos dice que $\mathbf{A} \models \varphi \lbrack \vec{b}]$. La prueba de que $\mathbf{A} \models \varphi \lbrack \vec{b}]$ implica que $ \mathbf{A} \models \varphi \lbrack \vec{a}]$ es similar.

    CASO $\varphi =\exists x_{j}\varphi_{1}$.

    Es similar al anterior.
  \end{proof}

  % Corollary 136. Sin prueba. Lemma 46.
  \begin{corollary} \label{corollary_46}
    \PN Si $\varphi$ es una sentencia, entonces $\mathbf{A} \models \varphi \lbrack \vec{a}] \Leftrightarrow \mathbf{A}
    \models \varphi \lbrack \vec{b}]$, cualesquiera sean las asignaciones $\vec{a}, \vec{b}$.
  \end{corollary}

  % Lemma 137. Con prueba. Lemma 47.
  \begin{lemma} \label{lemma_47}
    \PN \newline
    \begin{enumerate}[(a)]
      \item Si $Li(\varphi) \cup Li(\psi) \subseteq \{x_{i_{1}}, \dotsc, x_{i_{n}}\}$, entonces $\varphi \thicksim \psi$
      si y solo si la sentencia $\forall x_{i_{1}} \dotsc \forall x_{i_{n}}(\varphi \leftrightarrow \psi)$ es
      universalmente válida.
      \item Si $\varphi_{i} \thicksim \psi_{i}, i = 1, 2$, entonces $\lnot \varphi_{1} \thicksim \lnot \psi_{1},
      (\varphi_{1} \eta \varphi_{2}) \thicksim (\psi_{1} \eta \psi_{2}) \ \text{y} \ Qv\varphi_{1} \thicksim Qv \psi_{1}$.
      \item Si $\varphi \thicksim \psi \ \text{y} \ \alpha^{\prime}$ es el resultado de reemplazar en una fórmula $\alpha$
      algunas (posiblemente $0$) ocurrencias de $\varphi$ por $\psi$, entonces $\alpha \thicksim \alpha^{\prime}$.
    \end{enumerate}
  \end{lemma}
  \begin{proof}
    \begin{enumerate}[(a)]
      \item Tenemos que
        \[
          \begin{array}{l}
            \varphi \thicksim \psi \\
            \ \ \Updownarrow \text{ (por (6) de la def de}\models \text{)} \\
            \mathbf{A} \models (\varphi \leftrightarrow \psi )[\vec{a}]\text{, para todo } \mathbf{A}\text{ y toda }\vec{a}\in A^{\mathbb{N}} \\
            \ \ \Updownarrow \\
            \mathbf{A} \models (\varphi \leftrightarrow \psi )[\downarrow _{i_{n}}^{a}(\vec{a })]\text{, para todo }\mathbf{A}\text{, }a\in A\text{ y toda }\vec{a}\in A^{ \mathbb{N}} \\ \ \ \Updownarrow (\text{por (8) de la def de}\models ) \\ \mathbf{A} \models \forall x_{i_{n}}(\varphi \leftrightarrow \psi )[\vec{a}] \text{, para todo }\mathbf{A}\text{ y toda }\vec{a}\in A^{\mathbb{N}} \\ \ \ \Updownarrow \\ \mathbf{A} \models \forall x_{i_{n}}(\varphi \leftrightarrow \psi )[\downarrow _{i_{n-1}}^{a}(\vec{a})]\text{, para todo }\mathbf{A}\text{, }a\in A\text{ y toda }\vec{a}\in A^{\mathbb{N}} \\ \ \ \Updownarrow \text{ (por (8) de la def de}\models \text{)} \\ \mathbf{A} \models \forall x_{i_{n-1}}\forall x_{i_{n}}(\varphi \leftrightarrow \psi )[\vec{a}]\text{, para todo }\mathbf{A}\text{ y toda }\vec{a}\in A^{ \mathbb{N}} \\ \ \ \Updownarrow \\ \ \ \ \ \vdots \\ \ \ \Updownarrow \\ \mathbf{A} \models \forall x_{i_{1}}...\forall x_{i_{n}}(\varphi \leftrightarrow \psi )[\vec{a}]\text{, para todo }\mathbf{A}\text{ y toda }\vec{a}\in A^{ \mathbb{N}} \\ \ \ \Updownarrow \\ \forall x_{i_{1}}...\forall x_{i_{n}}(\varphi \leftrightarrow \psi )\text{ es universalmente valida}
          \end{array}
        \]
      \item Es dejado al lector.
      \item Por induccion en el $k$ tal que $\alpha \in F_{k}^{\tau }$.
    \end{enumerate}
  \end{proof}

  % Lemma 138. Con prueba. Lemma 48.
  \begin{lemma} \label{lemma_48}
    \PN Sea $F: \mathbf{A} \rightarrow \mathbf{B}$ un homomorfismo, entonces:
    \[
      F(t^{\mathbf{A}}[(a_{1}, a_{2}, \dotsc)] = t^{\mathbf{B}}[F(a_{1}), F(a_{2}), \dotsc)]
    \]
    \PN para cada $t \in \TAU, (a_{1}, a_{2}, \dotsc) \in A^{\mathbb{N}}$.
  \end{lemma}
  \begin{proof}
    Sea

    - Teo$_{k}$: Si $F:\mathbf{A} \rightarrow \mathbf{B}$ es un homomorfismo, entonces
    $\displaystyle F(t^{\mathbf{A}}[(a_{1},a_{2},...)]=t^{\mathbf{B}}[F(a_{1}),F(a_{2}),...)] $

    para cada $t\in T_{k}^{\tau }$, $(a_{1},a_{2},...)\in A^{\mathbb{N}}$.
    Teo$_{0}$ es trivial. Veamos que Teo$_{k}$ implica Teo$_{k+1}$. Supongamos que vale Teo$_{k}$ y supongamos $F:\mathbf{A} \rightarrow \mathbf{B}$ es un homomorfismo, $t\in T_{k+1}^{\tau }-T_{k}^{\tau } \ \text{y} \ \vec{a} = (a_{1},a_{2},...)\in A^{\mathbb{N}}$. Denotemos $(F(a_{1}),F(a_{2}),...)$ con $F(\vec{a})$. Por Lema 117, $t=f(t_{1}, \dotsc, t_{n})$, con $n\geq 1 $$\;f\in \mathcal{F}_{n} \ \text{y} \ t_{1}, \dotsc, t_{n}\in T_{k}^{\tau }$. Tenemos entonces

    \[
      \begin{array}{ccl}
        F(t^{\mathbf{A}}[\vec{a}]) &=& F(f(t_{1}, \dotsc, t_{n})^{\mathbf{A}}[\vec{a}]) \\
        &=& F(f^{\mathbf{A}}(t_{1}^{\mathbf{A}}[\vec{a}], \dotsc, t_{n}^{\mathbf{A}}[ \vec{a}])) \\
        &=& f^{\mathbf{B}}(F(t_{1}^{\mathbf{A}}[\vec{a}]), \dotsc, F(t_{n}^{\mathbf{A}}[ \vec{a}])) \\
        &=& f^{\mathbf{B}}(t_{1}^{\mathbf{B}}[F(\vec{a})], \dotsc, t_{n}^{\mathbf{B}}[F( \vec{a})])) \\
        &=& f(t_{1}, \dotsc, t_{n})^{\mathbf{B}}[F(\vec{a})] \\
        &=& t^{\mathbf{B}}[F(\vec{a})]
      \end{array}
    \]
  \end{proof}

  % Lemma 139. Con prueba. Lemma 49.
  \begin{lemma} \label{lemma_49}
    \PN Supongamos que $F: \mathbf{A} \rightarrow \mathbf{B}$ es un isomorfismo. Sea $\varphi \in F^{\tau}$, entonces:
    \[
      \mathbf{A} \models \varphi \lbrack (a_{1}, a_{2}, \dotsc)] \Leftrightarrow \mathbf{B} \models \varphi \lbrack
      (F(a_{1}), F(a_{2}), \dotsc)]
    \]
    \PN para cada $(a_{1}, a_{2}, \dotsc) \in A^{\mathbb{N}}$. En particular $\mathbf{A} \ \text{y} \ \mathbf{B}$ satisfacen las
    mismas sentencias de tipo $\tau$.
  \end{lemma}
  \begin{proof}
    %TODO
  \end{proof}

  % Lemma 140. Con prueba. Lemma 50.
  \begin{lemma} \label{lemma_50}
    \PN Supongamos que $\tau$ es algebraico, si $F: \mathbf{A} \rightarrow \mathbf{B}$ es un homomorfismo biyectivo,
    entonces $F$ es un isomorfismo.
  \end{lemma}
  \begin{proof}
    Solo falta probar que $F^{-1}$ es un homomorfismo. Supongamos que $c\in \mathcal{C}$. Ya que $F(c^{\mathbf{A}})=c^{\mathbf{B}}$, tenemos que $ F^{-1}(c^{\mathbf{B}})=c^{\mathbf{A}}$, por lo cual $F^{-1}$ cumple (1) de la definicion de homomorfismo. Supongamos ahora que $f\in \mathcal{F}_{n}$ y sean $b_{1}, \dotsc, b_{n}\in B$. Sean $a_{1}, \dotsc, a_{n}\in A$ tales que $ F(a_{i})=b_{i}$, $i=1, \dotsc, n$. Tenemos que

    $\displaystyle \begin{array}{ccl} F^{-1}(f^{\mathbf{B}}(b_{1}, \dotsc, b_{n})) & = & F^{-1}(f^{\mathbf{B} }(F(a_{1}), \dotsc, F(a_{n}))) \\ & = & F^{-1}(F(f^{\mathbf{A}}(a_{1}, \dotsc, a_{n})) \\ & = & f^{\mathbf{A}}(a_{1}, \dotsc, a_{n}) \\ & = & f^{\mathbf{A}}(F^{-1}(b_{1}), \dotsc, F^{-1}(b_{n})) \end{array} $

    por lo cual $F^{-1}$ satisface (2) de la definicion de homomorfismo
  \end{proof}

  % Lemma 141. Con prueba. Lemma 51.
  \begin{lemma} \label{lemma_51}
    \PN Supongamos que $\tau$ es algebraico, si $F: \mathbf{A} \rightarrow \mathbf{B}$ es un homomorfismo, entonces
    $I_{F}$ es un subuniverso de $\mathbf{B}$.
  \end{lemma}
  \begin{proof}
    Ya que $A\neq \varnothing $ tenemos que $I_{F}\neq \varnothing .$ Es claro que $ c^{\mathbf{B}}=F(c^{\mathbf{A}})\in I_{F}$ para cada $c\in \mathcal{C}$. Sea $f\in \mathcal{F}_{n}$ y sean $b_{1}, \dotsc, b_{n}\in I_{F}$ Sean $ a_{1}, \dotsc, a_{n}$ tales que $F(a_{i})=b_{i}$ $i=1, \dotsc, n$. Tenemos que

    $\displaystyle f^{\mathbf{B}}(b_{1}, \dotsc, b_{n})=f^{\mathbf{B}}(F(a_{1}), \dotsc, F(a_{n}))=F(f^{ \mathbf{A}}(a_{1}, \dotsc, a_{n}))\in I_{F} $

    por lo cual $I_{F}$ es cerrada bajo $f^{\mathbf{B}}$.
  \end{proof}

  % Lemma 142. Con prueba. Lemma 52.
  \begin{lemma} \label{lemma_52}
    \PN Supongamos que $\tau$ es algebraico, si $F: \mathbf{A} \rightarrow \mathbf{B}$ es un homomorfismo, entonces
    $\ker F$ es una congruencia sobre $\mathbf{A}$.
  \end{lemma}
  \begin{proof}
    Sea $f\in \mathcal{F}_{n}$. Supongamos que $a_{1}, \dotsc, a_{n},b_{1}, \dotsc, b_{n} \in A$ son tales que $a_{1}\ker Fb_{1}, \dotsc, a_{n}\ker Fb_{n}$. Tenemos entonces que

    $\displaystyle \begin{array}{ccl} F(f^{\mathbf{A}}(a_{1}, \dotsc, a_{n})) & = & f^{\mathbf{B} }(F(a_{1}), \dotsc, F(a_{n})) \\ & = & f^{\mathbf{B}}(F(b_{1}), \dotsc, F(b_{n})) \\ & = & F(f^{\mathbf{B}}(b_{1}, \dotsc, b_{n})) \end{array} $

    lo cual nos dice que $f^{\mathbf{A}}(a_{1}, \dotsc, a_{n})\ker Ff^{\mathbf{B} }(b_{1}, \dotsc, b_{n})$
  \end{proof}

  % Lemma 143. Con prueba. Lemma 53.
  \begin{lemma} \label{lemma_53}
    \PN $\pi_{\theta}: \mathbf{A} \rightarrow \mathbf{A}/\theta$ es un homomorfismo cuyo núcleo es $\theta$.
  \end{lemma}
  \begin{proof}
    Sea $c\in \mathcal{C}$. Tenemos que

    \[
      \pi_{\theta}(c^{\mathbf{A}}) = c^{\mathbf{A}}/\theta = c^{\mathbf{A}/\theta}
    \]

    Sea $f\in \mathcal{F}_{n}$, con $n\geq 1$ y sean $a_{1}, \dotsc, a_{n}\in A$. Tenemos que

    \[
      \begin{array}{ccl}
        \pi_{\theta}(f^{\mathbf{A}}(a_{1}, \dotsc, a_{n})) &=& f^{\mathbf{A}}(a_{1}, \dotsc, a_{n})/\theta \\
        &=& f^{\mathbf{A}/\theta}(a_{1}/\theta, \dotsc, a_{n}/\theta) \\
        &=& f^{\mathbf{A}/\theta}(\pi_{\theta}(a_{1}), \dotsc, \pi_{\theta}(a_{n}))
      \end{array}
    \]
    con lo cual $\pi _{\theta }$ es un homomorfismo. Es trivial que $\ker \pi _{\theta }=\theta $.
  \end{proof}

  % Corollary 144. Con prueba. Lemma 54.
  \begin{corollary} \label{corollary_54}
    \PN Para cada $t \in \TAU, \vec{a} \in A^{\mathbb{N}}$, se tiene que $t^{\mathbf{A}/\theta}[(a_{1}/\theta,
    a_{2}/\theta, \dotsc)] = t^{\mathbf{A}}[(a_{1}, a_{2}, \dotsc)]/\theta$.
  \end{corollary}
  \begin{proof}
    Ya que $\pi _{\theta }$ es un homomorfismo, se puede aplicar el Lema 136.
  \end{proof}

  % Theorem 145. Con prueba. Lemma 55.
  \begin{theorem} \label{theorem_55}
    \PN Sea $F: \mathbf{A} \rightarrow \mathbf{B}$ un homomorfismo sobreyectivo, entonces:
    \begin{eqnarray*}
      A/\ker F &\rightarrow& B \\
      a/\ker F &\rightarrow& F(a)
    \end{eqnarray*}
    \PN define sin ambiguedad una función $\bar{F}$ la cual es un isomorfismo de $\mathbf{A}/\ker F$ en $\mathbf{B}$.
  \end{theorem}
  \begin{proof}
    Notese que la definicion de $\bar{F}$ es inambigua ya que si $a/\ker F=a^{\prime }/\ker F$, entonces $F(a)=F(a^{\prime }).$ Ya que $F$ es sobre, tenemos que $\bar{F}$ lo es. Supongamos que $\bar{F}(a/\ker F)=\bar{F} (a^{\prime }/\ker F).$ Claramente entonces tenemos que $F(a)=F(a^{\prime })$ , lo cual nos dice que $a/\ker F=a^{\prime }/\ker F$. Esto prueba que $\bar{F }$ es inyectiva. Para ver que $\bar{F}$ es un isomorfismo, por el Lema 138, basta con ver que $\bar{F}$ es un homomorfismo. Sea $c\in \mathcal{C}$. Tenemos que

    $\displaystyle \bar{F}(c^{\mathbf{A}/\ker F})=\bar{F}(c^{\mathbf{A}}/\ker F)=F(c^{\mathbf{A} })=c^{\mathbf{B}} $

    Sea $f\in \mathcal{F}_{n}$. Sean $a_{1}, \dotsc, a_{n}\in A$. Tenemos que
    \[
      \begin{array}{ccl}
        \bar{F}(f^{\mathbf{A}/\ker F}(a_{1}/\ker F, \dotsc, a_{n}/\ker F)) &=& \bar{F}(f^{\mathbf{A}}(a_{1}, \dotsc,
          a_{n})/\ker F) \\
        &=& F(f^{\mathbf{A}}(a_{1}, \dotsc, a_{n})) \\
        &=& f^{\mathbf{B}}(F(a_{1}), \dotsc, F(a_{n})) \\
        &=& f^{\mathbf{B}}(\bar{F}(a_{1}/\ker F), \dotsc, \bar{F}(a_{n}/\ker F))
      \end{array}
    \]

    con lo cual $\bar{F}$ cunple (2) de la definicion de homomorfismo.
  \end{proof}

  % Lemma 146. Nada. Lemma 56.
  \begin{lemma}
    \PN Este lema no se evalua.
  \end{lemma}

  % Lemma 147. Nada. Lemma 57.
  \begin{lemma}
    \PN Este lema no se evalua.
  \end{lemma}

  % Lemma 148. Con prueba. Lemma 58.
  \begin{lemma} \label{lemma_58}
    \PN Sea $\tau$ un tipo cualquiera y supongamos $t \in \TAU$. Si $t =_{d} t(v_{1}, \dotsc, v_{n})$ entonces se da
    alguna de las siguientes:
    \begin{enumerate}
      \item $t = c$ para algún $c \in \mathcal{C}$.
      \item $t = v_{j}$ para algún j.
      \item $t = f(t_{1}, \dotsc, t_{m})$, con $f \in \mathcal{F}_{m} \ \text{y} \ t_{1}, \dotsc, t_{m} \in T_{k-1}^{\tau}$ tales
      que las variables que ocurren en cada uno de ellos están en $\{v_{1}, \dotsc, v_{n}\}$.
    \end{enumerate}
  \end{lemma}
  \begin{proof}
  \end{proof}

  % Lemma 149. Con prueba. Lemma 59.
  \begin{lemma} \label{lemma_59}
    \PN Sea $\tau$ un tipo cualquiera y supongamos $t \in \TAU$. Si $t =_{d} t(v_{1}, \dotsc, v_{n})$. Sea $\mathbf{A}$
    un modelo de tipo $\tau$. Sean $a_{1}, \dotsc, a_{n} \in A$. Se tiene que:
    \begin{enumerate}
      \item Si $t = c$ entonces $t^{\mathbf{A}}[a_{1}, \dotsc, a_{n}]= c^{\mathbf{A}}$.
      \item Si $t = v_{j}$ entonces $t^{\mathbf{A}}[a_{1}, \dotsc, a_{n}]= a_{j}$.
      \item Si $t = f(t_{1}, \dotsc, t_{m})$, con $f \in \mathcal{F}_{m} \ \text{y} \ t_{1}, \dotsc, t_{m} \in \TAU$, entonces:
      \[
        t^{\mathbf{A}}[a_{1}, \dotsc, a_{n}]= f^{\mathbf{A}}(t_{1}^{\mathbf{A}}[a_{1}, \dotsc, a_{n}], \dotsc,
        t_{m}^{\mathbf{A}}[a_{1}, \dotsc, a_{n}])
      \]
    \end{enumerate}
  \end{lemma}
  \begin{proof}
    \begin{enumerate}[(1)]
      \item trivial
      \item trivial
      \item Sea $\vec{b}$ una asignación tal que a cada $v_{i}$le asigna el valor $a_{i}$. Por definición tenemos EquivElim
      \begin{eqnarray*}
        t^{\mathbf{A}}[a_{1}, \dotsc, a_{n}] &=& t^{\mathbf{A}}[\vec{b}] \\
        &=& f^{\mathbf{A}}(t_{1}^{\mathbf{A}}[\vec{b}], \dotsc, t_{m}^{\mathbf{A}}[\vec{b}]) \\
        &=& f^{\mathbf{A}}(t_{1}^{\mathbf{A}}[a_{1}, \dotsc, a_{n}], \dotsc, t_{m}^{\mathbf{A}}[a_{1}, \dotsc, a_{n}])
      \end{eqnarray*}
    \end{enumerate}
  \end{proof}

  % Lemma 150. Con prueba. Lemma 60.
  \begin{lemma} \label{lemma_60}
    \PN \textbf{(De reemplazo para términos)}. Supongamos $t =_{d} t(w_{1}, \dotsc, w_{k}), s_{1} =_{d} s_{1}(v_{1},
    \dotsc, v_{n}), \dotsc, s_{k} =_{d} s_{k}(v_{1}, \dotsc, v_{n})$. Todas las variables de $t(s_{1}, \dotsc, s_{k})$
    están en $\{v_{1}, \dotsc, v_{n}\}$ y si declaramos $t(s_{1}, \dotsc, s_{k}) =_{d} t(s_{1}, \dotsc, s_{k})(v_{1},
    \dotsc, v_{n})$, entonces para cada estructura $\mathbf{A} \ \text{y} \ a_{1}, \dotsc a_{n} \in A$, se tiene:
    \[
      t(s_{1}, \dotsc, s_{k})^{\mathbf{A}}[a_{1}, \dotsc, a_{n}] = t^{\mathbf{A}}[s_{1}^{\mathbf{A}}[a_{1}, \dotsc,
      a_{n}], \dotsc, s_{k}^{\mathbf{A}}[a_{1}, \dotsc, a_{n}]]
    \]
  \end{lemma}
  \begin{proof}
    Probaremos que valen (a) y (b), por induccion en el $l$ tal que $t\in T_{l}^{\tau }.$ El caso $l=0$ es dejado al lector. Supongamos entonces que valen (a) y (b) siempre que $t\in T_{l}^{\tau }$ y veamos que entonces valen (a) y (b) cuando $t\in T_{l+1}^{\tau }-T_{l}^{\tau }$. Hay $f\in \mathcal{F} _{m} \ \text{y} \ t_{1}, \dotsc, t_{m}\in T_{l}^{\tau }$ tales que $ t_{1}=_{d}t_{1}(w_{1}, \dotsc, w_{k}), \dotsc, t_{m}=_{d}t_{m}(w_{1}, \dotsc, w_{k}) \ \text{y} \  t=f(t_{1}, \dotsc, t_{m})$. Notese que por (a) de la HI tenemos que

    $\displaystyle t_{i}(s_{1}, \dotsc, s_{k})=_{d}t_{i}(s_{1}, \dotsc, s_{k})(v_{1}, \dotsc, v_{n})\text{, } i=1, \dotsc, m $

    lo cual ya que
    $\displaystyle t(s_{1}, \dotsc, s_{k})=f(t_{1}(s_{1}, \dotsc, s_{k}), \dotsc, t_{m}(s_{1}, \dotsc, s_{k})) $

    nos dice que
    $\displaystyle t(s_{1}, \dotsc, s_{k})=_{d}t(s_{1}, \dotsc, s_{k})(v_{1}, \dotsc, v_{n}) $

    obteniendo asi (a). Para probar (b) notemos que por (b) de la hipotesis inductiva
    $\displaystyle t_{j}(s_{1}, \dotsc, s_{k})^{\mathbf{A}}[\vec{a}]=t_{j}^{\mathbf{A}}[s_{1}^{ \mathbf{A}}[\vec{a}], \dotsc, s_{k}^{\mathbf{A}}[\vec{a}]],j=1, \dotsc, m $

    lo cual nos dice que
    $\displaystyle \begin{array}{ccl} t(s_{1}, \dotsc, s_{k})^{\mathbf{A}}[\vec{a}] & = & f(t_{1}(s_{1}, \dotsc, s_{k}), \dotsc, t_{m}(s_{1}, \dotsc, s_{k}))^{\mathbf{A}}[\vec{a}] \\ & = & f^{\mathbf{A}}(t_{1}(s_{1}, \dotsc, s_{k})^{\mathbf{A}}[\vec{a} ], \dotsc, t_{m}(s_{1}, \dotsc, s_{k})^{\mathbf{A}}[\vec{a}]) \\ & = & f^{\mathbf{A}}(t_{1}^{\mathbf{A}}[s_{1}^{\mathbf{A}}[\vec{a} ], \dotsc, s_{k}^{\mathbf{A}}[\vec{a}]], \dotsc, t_{m}^{\mathbf{A}}[s_{1}^{\mathbf{A}}[ \vec{a}], \dotsc, s_{k}^{\mathbf{A}}[\vec{a}]]) \\ & = & t^{\mathbf{A}}[s_{1}^{\mathbf{A}}[\vec{a}], \dotsc, s_{k}^{\mathbf{A}}[\vec{ a}]] \end{array} $
  \end{proof}

  % Lemma 151. Con prueba. Lemma 61.
  \begin{lemma} \label{lemma_61}
    \PN Sea $\tau$ un tipo cualquiera y supongamos $\varphi \in F^{\tau}$. Si $\varphi =_{d} \varphi(v_{1}, \dotsc,
    v_{n})$, entonces se da una y solo una de las siguientes:
    \begin{enumerate}[(1)]
      \item $\varphi = (t \equiv s)$, con $t, s \in \TAU$, únicos y tales que las variables que ocurren en $t$ o en $s$
      están todas en $ \{v_{1}, \dotsc, v_{n}\}$.
      \item $\varphi = r(t_{1}, \dotsc, t_{m})$, con $r \in \mathcal{R}_{m} \ \text{y} \ t_{1}, \dotsc, t_{m} \in \TAU$, únicos y
      tales que las variables que ocurren en cada $t_{i}$ están todas en $\{v_{1}, \dotsc, v_{n}\}$.
      \item $\varphi = (\varphi_{1} \wedge \varphi_{2})$, con $\varphi_{1}, \varphi_{2} \in F^{\tau}$, únicas y tales
      que $Li(\varphi_{1}) \cup Li(\varphi_{2}) \subseteq \{v_{1}, \dotsc, v_{n}\}$.
      \item $\varphi = (\varphi_{1} \vee \varphi_{2})$, con $\varphi_{1}, \varphi_{2} \in F^{\tau}$, únicas y tales que
      $Li(\varphi_{1}) \cup Li(\varphi_{2}) \subseteq \{v_{1}, \dotsc, v_{n}\}$.
      \item $\varphi = (\varphi_{1} \rightarrow \varphi_{2})$, con $\varphi_{1}, \varphi_{2} \in F^{\tau}$, únicas y
      tales que $Li(\varphi_{1}) \cup Li(\varphi_{2}) \subseteq \{v_{1}, \dotsc, v_{n}\}$.
      \item $\varphi = (\varphi_{1} \leftrightarrow \varphi_{2})$, con $ \varphi_{1}, \varphi_{2} \in F^{\tau}$, únicas
      y tales que $Li(\varphi_{1}) \cup Li(\varphi_{2}) \subseteq \{v_{1}, \dotsc, v_{n}\}$.
      \item $\varphi = \lnot \varphi_{1}$, con $\varphi_{1} \in F^{\tau}$, única y tal que $Li(\varphi_{1}) \subseteq
      \{v_{1}, \dotsc, v_{n}\}$.
      \item $\varphi = \forall v_{j} \varphi_{1}$, con $v_{j} \in \{v_{1}, \dotsc, v_{n}\} \ \text{y} \ \varphi_{1} \in
      F^{\tau}$, únicas y tales que $ Li(\varphi_{1}) \subseteq \{v_{1}, \dotsc, v_{n}\}$.
      \item $\varphi = \forall v \varphi_{1}$, con $v \in Var-\{v_{1}, \dotsc, v_{n}\} \ \text{y} \ \varphi_{1}\ in F^{\tau}$,
      únicas y tales que $ Li(\varphi_{1}) \subseteq \{v_{1}, \dotsc, v_{n}, v\}$.
      \item $\varphi = \exists v_{j} \varphi_{1}$, con $v_{j} \in \{v_{1}, \dotsc, v_{n}\} \ \text{y} \ \varphi_{1} \in
      F^{\tau}$, únicas y tales que $ Li(\varphi_{1}) \subseteq \{v_{1}, \dotsc, v_{n}\}$.
      \item $\varphi = \exists v \varphi_{1}$, con $v \in Var-\{v_{1}, \dotsc, v_{n}\} \ \text{y} \ \varphi_{1} \in F^{\tau}$,
      únicas y tales que $ Li(\varphi_{1}) \subseteq \{v_{1}, \dotsc, v_{n}, v\}$.
    \end{enumerate}
  \end{lemma}
  \begin{proof}
    Induccion en el k tal que $\varphi \in F_{k}^{\tau}$
  \end{proof}

  % Lemma 152. Con prueba. Lemma 62.
  \begin{lemma} \label{lemma_62}
    \PN Supongamos $\varphi =_{d} \varphi(v_{1}, \dotsc, v_{n})$. Sea $\mathbf{A} = (A, i)$ un modelo de tipo $\tau$ y
    sean $a_{1}, \dotsc, a_{n} \in A$, entonces:
    \begin{enumerate}
      \item Si $\varphi = (t \equiv s)$, entonces:
      \begin{center}
        $\mathbf{A} \models \varphi \lbrack a_{1}, \dotsc, a_{n}] \ \text{si y solo si} \ t^{\mathbf{A}}[a_{1}, \dotsc,
        a_{n}]=s^{\mathbf{A}}[a_{1}, \dotsc, a_{n}]$
      \end{center}
      \item Si $\varphi = r(t_{1}, \dotsc, t_{m})$, entonces:
      \begin{equation*}
        \mathbf{A} \models \varphi \lbrack a_{1}, \dotsc, a_{n}] \ \text{si y solo si} \ (t_{1}^{A}[a_{1}, \dotsc,
        a_{n}], \dotsc, t_{m}^{A}[a_{1}, \dotsc, a_{n}])\in r^{A}
      \end{equation*}
      \item Si $\varphi = (\varphi_{1} \wedge \varphi_{2})$ entonces:
      \[
        \mathbf{A} \models \varphi \lbrack a_{1}, \dotsc, a_{n}] \ \text{si y solo si} \ \mathbf{A} \models \varphi_{1}
        [a_{1}, \dotsc, a_{n}] \ \text{y} \ \mathbf{A} \models \varphi_{2}[a_{1}, \dotsc, a_{n}]
      \]
      \item Si $\varphi = (\varphi_{1} \vee \varphi_{2})$ entonces:
      \[
        \mathbf{A} \models \varphi \lbrack a_{1}, \dotsc, a_{n}] \ \text{si y solo si} \ \mathbf{A} \models \varphi_{1}
        [a_{1}, \dotsc, a_{n}] \ \text{o} \ \mathbf{A} \models \varphi_{2}[a_{1}, \dotsc, a_{n}]
      \]
      \item Si $\varphi = (\varphi_{1} \rightarrow \varphi_{2})$ entonces:
      \[
        \mathbf{A} \models \varphi \lbrack a_{1}, \dotsc, a_{n}] \ \text{si y solo si} \ \mathbf{A} \models \varphi_{2}
        [a_{1}, \dotsc, a_{n}] \ \text{o} \ \mathbf{A} \not\models \varphi_{1}[a_{1}, \dotsc, a_{n}]
      \]
      \item Si $\varphi = (\varphi_{1} \leftrightarrow \varphi_{2})$ entonces:
      \begin{eqnarray*}
        \mathbf{A} \models \varphi \lbrack a_{1}, \dotsc, a_{n}] \ \text{si y solo si ya sea} \ && \mathbf{A} \models
        \varphi_{1}[a_{1}, \dotsc, a_{n}] \ \text{y} \ \mathbf{A} \models \varphi_{2}[a_{1}, \dotsc, a_{n}] \ \text{o}
        \\
        && \mathbf{A} \not\models \varphi_{1}[a_{1}, \dotsc, a_{n}] \ \text{y} \ \mathbf{A} \not\models \varphi_{2}
        [a_{1}, \dotsc, a_{n}]
      \end{eqnarray*}
      \item Si $\varphi = \lnot \varphi_{1}$ entonces:
      \[
        \mathbf{A} \models \varphi \lbrack a_{1}, \dotsc, a_{n}] \ \text{si y solo si} \ \mathbf{A} \not\models
        \varphi_{1}[a_{1}, \dotsc, a_{n}]
      \]
      \item Si $\varphi = \forall v\varphi_{1}$ con $v\not\in \{v_{1}, \dotsc, v_{n}\} \ \text{y} \ \varphi_{1} =_{d}
      \varphi_{1}(v_{1}, \dotsc, v_{n}, v)$ entonces:
      \[
        \mathbf{A} \models \varphi \lbrack a_{1}, \dotsc, a_{n}] \ \text{si y solo si} \ \mathbf{A} \models \varphi_{1}
        [a_{1}, \dotsc, a_{n},a] \ \text{para todo} \ a \in A
      \]
      \item Si $\varphi = \forall v_{j}\varphi_{1}$ entonces:
      \[
        \mathbf{A} \models \varphi \lbrack a_{1}, \dotsc, a_{n}] \ \text{si y solo si} \ \mathbf{A} \models \varphi_{1}
        [a_{1}, \dotsc, a, \dotsc, a_{n}] \ \text{para todo} \ a \in A
      \]
      \item) Si $\varphi = \exists v\varphi_{1}$ con $v\not\in \{v_{1}, \dotsc, v_{n}\} \ \text{y} \ \varphi_{1} =_{d}
      \varphi_{1}(v_{1}, \dotsc, v_{n}, v)$ entonces:
      \[
        \mathbf{A} \models \varphi \lbrack a_{1}, \dotsc, a_{n}] \ \text{si y solo si} \ \mathbf{A} \models \varphi_{1}
        [a_{1}, \dotsc, a_{n}, a] \ \text{para algún} \ a \in A
      \]
      \item) Si $\varphi = \exists v_{j}\varphi_{1}$ entonces:
      \[
        \mathbf{A} \models \varphi \lbrack a_{1}, \dotsc, a_{n}] \ \text{si y solo si} \ \mathbf{A} \models \varphi_{1}
        [a_{1}, \dotsc, a, \dotsc, a_{n}] \ \text{para algún} \ a \in A
      \]
    \end{enumerate}
  \end{lemma}
  \begin{proof}
  \end{proof}

  % Lemma 153. Sin prueba. Lemma 63.
  \begin{lemma} \label{lemma_63}
    \PN Si $Qv$ ocurre en $\varphi$ a partir de $i$, entonces hay una única fórmula $\psi$ tal que $Qv\psi$ ocurre en
    $\varphi$ a partir de $i$.
  \end{lemma}

  % Lemma 154. Sin prueba. Lemma 64.
  \begin{lemma} \label{lemma_64}
    \PN Supongamos $\varphi =_{d} \varphi(w_{1}, \dotsc, w_{k}), t_{1} =_{d} t_{1}(v_{1}, \dotsc, v_{n}), \dotsc, t_{k}
    =_{d} t_{k}(v_{1}, \dotsc, v_{n})$ y supongamos además que cada $w_{j}$ es sustituible por $t_{j}$ en $\varphi$,
    entonces:
    \begin{enumerate}[(a)]
      \item $Li(\varphi(t_{1}, \dotsc, t_{k})) \subseteq \{v_{1}, \dotsc, v_{n}\}$
      \item Si declaramos $\varphi(t_{1}, \dotsc, t_{k}) =_{d} \varphi(t_{1}, \dotsc, t_{k})(v_{1}, \dotsc, v_{n})$,
      entonces para cada estructura $\mathbf{A} \ \text{y} \ \vec{a} \in A^{n}$ se tiene:
        \[
          \mathbf{A} \models \varphi(t_{1}, \dotsc, t_{k})[\vec{a}] \ \text{si y solo si} \ \mathbf{A} \models \varphi
          \lbrack t_{1}^{\mathbf{A}}[\vec{a}], \dotsc, t_{k}^{ \mathbf{A}}[\vec{a}]]
        \]
    \end{enumerate}
  \end{lemma}

	\section{Teorias de primer orden}

  % Lemma 155. Con prueba. Lemma 65.
  \begin{lemma} \label{lemma_65}
    \PN Si $(\varphi_{1}, \varphi_{2}) \in Generaliz^{\tau}$, entonces el nombre de constante $c$ del cual habla la
    definición de $Generaliz^{\tau}$ está unívocamente determinado por el par $(\varphi_{1}, \varphi_{2})$.
  \end{lemma}
  \begin{proof}
    \PN Recordemos la definición de $Generaliz^{\tau}$.
    \begin{eqnarray*}
      Generaliz^{\tau} &=& \{(\psi, \forall v \tilde{\psi}): \psi \in S^{\tau}, \ v \ \text{no ocurre en} \ \psi \
      \text{y existe } \ c \in \mathcal{C} \ \text{tal que} \ \tilde{\psi} = \text{resultado de} \\
      && \text{reemplazar} \ \text{en} \ \psi \ \text{cada ocurrencia de} \ c \ \text{por} \ v\}
    \end{eqnarray*}

    \PN Notese que $c$ es el único nombre de constante que ocurre en $\varphi_{1}$ y no ocurre en $\varphi_{2}$.
  \end{proof}

  % Lemma 156. Con prueba. Lemma 66.
  \begin{lemma} \label{lemma_66}
    \PN Si $(\varphi_{1}, \varphi_{2}) \in Elec^{\tau}$, entonces el nombre de constante $e$ del cual habla la
    definición de $Elec^{\tau}$ está unívocamente determinado por el par $(\varphi_{1}, \varphi_{2})$.
  \end{lemma}
  \begin{proof}
    \PN Recordemos la definición de $Elec^{\tau}$.
    \begin{eqnarray*}
      Elect^{\tau} = \{(\exists v \varphi(v), \varphi(e)): \varphi =_{d} \varphi(v), \ Li(\varphi) = \{v\} \ \text{y} \
      e \in \mathcal{C} \ \text{no ocurre en} \ \varphi\}
    \end{eqnarray*}

    \PN Notese que $e$ es el único nombre de constante que ocurre en $\varphi_{1}$ y no ocurre en $\varphi_{2}$.
  \end{proof}

  % Lemma 157. Con prueba. Lemma 67.
  \begin{lemma} \label{lemma_67}
    \PN Todas las reglas excepto las reglas de elección y generalización son universales en el sentido que si $\varphi$
    se deduce de $\psi_{1}, \dotsc, \psi_{k}$ por alguna de estas reglas, entonces $\left((\psi_{1} \wedge \dotsc \wedge
    \psi_{k}) \rightarrow \varphi \right)$ es una sentencia universalmente válida.
  \end{lemma}
  \begin{proof}
    \begin{itemize}
      \item \textbf{DISJUNCIÓN-INTRODUCCIÓN:}
      \item \textbf{CONJUNCIÓN-INTRODUCCIÓN:}
      \item \textbf{CONJUNCIÓN-ELIMINACIÓN:}
      \item \textbf{EQUIVALENCIA-INTRODUCCIÓN:}
      \item \textbf{EQUIVALENCIA-ELIMINACIÓN:}
      \item \textbf{PARTICULARIZACIÓN:}

      \item \textbf{EXISTENCIA:} Supongamos $\varphi =_{d}\varphi (v)$, $t\in T_{c}^{\tau }$ y $\mathbf{A}$ es una
      estructura de tipo $\tau $ tal que $\mathbf{A}\models \varphi (t)$. Sea $t^{\mathbf{A}}$ el valor que toma $t$ en
      $\mathbf{A}$. Por el Lema 148 tenemos que $\mathbf{A}\models \varphi \left\lbrack t^{\mathbf{A}}\right\rbrack $,
      por lo cual tenemos que $\mathbf{A}\models \exists v\varphi (v)$.

      \item \textbf{EVOCACIÓN:}
      \item \textbf{ABSURDO:}
      \item \textbf{CONMUTATIVIDAD:}
      \item \textbf{MODUS-PONENS:}
      \item \textbf{DIVISIÓN-POR-CASOS:}

      \item \textbf{REEMPLAZO:} Debemos probar que si $(\psi_{1},\psi_{2},\varphi )\in Reemp^{\tau }=Reemp1^{\tau }\cup Reemp2^{\tau }$ , entonces $\left( (\psi_{1}\wedge \psi_{2})\rightarrow \varphi \right) $ es una sentencia universalmente valida. El caso en el que $(\psi_{1},\psi_{2},\varphi )\in Reemp1^{\tau }$ es facil y lo dejaremos al lector. Para el caso en el que $(\psi_{1},\psi_{2},\varphi )\in Reemp2^{\tau }$ nos hara falta un resultado un poco mas general. Veamos por induccion en $k$ que si se dan las siguienes condiciones

      - $\alpha \in F_{k}^{\tau }$ y $\varphi ,\psi \in F^{\tau }$
      - $\mathbf{A}$ es una estructura de tipo $\tau $
      - $\overline{\alpha }=$ resultado de reemplazar en $\alpha $ una ocurrencia de $\varphi$ por $\psi$,
      - $\mathbf{A}\models \varphi \left\lbrack \vec{a}\right\rbrack $ si y solo si $ \mathbf{A}\models \psi \left\lbrack \vec{a}\right\rbrack $, para cada $\vec{a}\in A^{ \mathbb{N}}$
      entonces se da que

      - $\mathbf{A}\models \alpha \left\lbrack \vec{a}\right\rbrack $ si y solo si $ \mathbf{A}\models \overline{\alpha }\left\lbrack \vec{a}\right\rbrack $, para cada $\vec{ a}\in A^{\mathbb{N}}$.
      CASO $k=0.$

      Entonces $\alpha $ es atomica y por lo tanto ya que $\alpha $ es la unica subformula de $\alpha $, la situacion es facil de probar.

      CASO $\alpha =\forall x_{i}\alpha _{1}.$

      Si $\varphi =\alpha $, entonces la situacion es facil de probar. Si $\varphi \neq \alpha $, entonces la ocurrencia de $\varphi$ a reemplazar sucede en $\alpha _{1}$ y por lo tanto $\overline{\alpha }=\forall x_{i} \overline{\alpha _{1}}.$ Se tiene entonces que para un $\vec{a}$ dado,

      $\displaystyle \begin{array}{c} \mathbf{A}\models \alpha \left\lbrack \vec{a}\right\rbrack \\ \Updownarrow \\ \mathbf{A}\models \alpha _{1}\left\lbrack \downarrow _{i}^{a}\vec{a}\right\rbrack ,\text{ para cada }a\in A \\ \Updownarrow \\ \mathbf{A}\models \overline{\alpha _{1}}\left\lbrack \downarrow _{i}^{a}\vec{a} \right\rbrack ,\text{ para cada }a\in A \\ \Updownarrow \\ \mathbf{A}\models \overline{\alpha }\left\lbrack \vec{a}\right\rbrack \end{array} $

      CASO $\alpha =(\alpha _{1}\vee \alpha _{2})$.
      Si $\varphi =\alpha $, entonces la situacion es facil de probar. Supongamos $\varphi \neq \alpha $ y supongamos que la ocurrencia de $\varphi$ a reemplazar sucede en $\alpha _{1}$. Entonces $\overline{\alpha }=( \overline{\alpha _{1}}\vee \alpha _{2})$ y tenemos que

      $\displaystyle \begin{array}{c} \mathbf{A}\models \alpha \left\lbrack \vec{a}\right\rbrack \\ \Updownarrow \\ \mathbf{A}\models \alpha _{1}\left\lbrack \vec{a}\right\rbrack \text{ o }\mathbf{A} \models \alpha _{2}\left\lbrack \vec{a}\right\rbrack \\ \Updownarrow \\ \mathbf{A}\models \overline{\alpha _{1}}\left\lbrack \vec{a}\right\rbrack \text{ o } \mathbf{A}\models \alpha _{2}\left\lbrack \vec{a}\right\rbrack \\ \Updownarrow \\ \mathbf{A}\models \overline{\alpha }\left\lbrack \vec{a}\right\rbrack \end{array} $

      \item \textbf{TRANSITIVIDAD:}
    \end{itemize}
  \end{proof}

  % Lemma 158. Con prueba. Lemma 68.
  \begin{lemma} \label{lemma_68}
    \PN Sea $\pmb{\varphi} \in S^{\tau +}$, hay únicos $n \geq 1$ y $\varphi_{1}, \dotsc, \varphi_{n} \in S^{\tau}$
    tales que $\pmb{\varphi} = \varphi_{1} \dotsc \varphi_{n}$.
  \end{lemma}
  \begin{proof}
    \PN Debemos probar la unicidad de $\varphi_{1}, \dotsc, \varphi_{n}$, pero esto vale utilizando la
    \textbf{Proposition~\ref{proposition_1}} Solo hay que probar la unicidad la cual sigue de la Proposicion 128.
  \end{proof}

  % Lemma 159. Con prueba. Lemma 69.
  \begin{lemma} \label{lemma_69}
    \PN Sea $\mathbf{J} \in Just^{+}$, hay únicos $n \geq 1$ y $J_{1}, \dotsc, J_{n} \in Just$ tales que $\mathbf{J} =
    J_{1} \dotsc J_{n}$.
  \end{lemma}
  \begin{proof}
    \PN Supongamos $J_{1}, \dotsc, J_{n}$, $J_{1}^{\prime}, \dotsc, J_{m}^{\prime}$, con $ n,m\geq 1$, son
    justificaciones tales que $J_{1}\dotsc J_{n}=J_{1}^{\prime}\dotsc J_{m}^{\prime}$. Es facil ver que entonces tenemos
    $J_{1}=J_{1}^{\prime}$, por lo cual $J_{2}\dotsc J_{n}=J_{2}^{\prime}\dotsc J_{m}^{\prime}$. Un argumento inductivo
    nos dice que entonces $n=m$ y $J_{i}=J_{i}^{\prime}$, $ i=1, \dotsc, n$
  \end{proof}

  % Lemma 160. Con prueba. Lemma 70.
  \begin{lemma} \label{lemma_70}
    \PN Sea $(\pmb{\varphi}, \mathbf{J})$ una prueba de $\varphi$ en $(\Sigma, \tau)$.
    \begin{enumerate}
      \item Sea $m \in \mathbb{N}$ tal que $\mathbf{J}_{i} \neq \mathrm{HIPOTESIS}\bar{m}$, para cada $i = 1, \dotsc,
      n(\pmb{\varphi})$. Supongamos que $\mathbf{J}_{i} = \mathrm{HIPOTESIS}\bar{k}$ y que $\mathbf{J}_{j} =
      \mathrm{TESIS}\bar{k} \alpha$, con $\lbrack\alpha\rbrack_{1} \notin Num$. Sea $\mathbf{\tilde{J}}$ el resultado de
      reemplazar en $\mathbf{J}$ la justificación $\mathbf{J}_{i}$ por $\mathrm{HIPOTESIS}\bar{m}$ y reemplazar la
      justificación $\mathbf{J}_{j}$ por $\mathrm{TESIS}\bar{m}\alpha$, entonces $(\pmb{\varphi}, \mathbf{\tilde{J}})$
      es una prueba de $\varphi$ en $(\Sigma, \tau)$.
      \item Sea $\mathcal{C}_{1}$ el conjunto de nombres de constante que ocurren en alguna $\pmb{\varphi}_{i}$ y que
      no pertenecen a $\mathcal{C}$. Sea $e \in \mathcal{C}_{1}-\mathcal{C}$. Sea $\tilde{e} \notin \mathcal{C} \cup
      \mathcal{C}_{1}$ tal que $(\mathcal{C} \cup (\mathcal{C}_{1}-\{e\}) \cup \{\tilde{e}\}, \mathcal{F}, \mathcal{R},
      a)$ es un tipo. Sea $\tilde{\varphi}_{i} =$ resultado de reemplazar en $\pmb{\varphi}_{i}$ cada ocurrencia de
      $e$ por $\tilde{e}$, entonces $(\pmb{\tilde{\varphi}}_{1} \dotsc \pmb{\tilde{\varphi}}_{n(\mathbf{\varphi})},
      \mathbf{J})$ es una prueba de $\varphi$ en $(\Sigma, \tau)$.
    \end{enumerate}
  \end{lemma}
  \begin{proof}
    \begin{enumerate}
      \item Obvio.
      \item Sean

      $\displaystyle \begin{array}{rcl} \tau _{1} & =& (\mathcal{C}\cup \mathcal{C}_{1},\mathcal{F},\mathcal{R},a) \\ \tau _{2} & =& (\mathcal{C}\cup (\mathcal{C}_{1}-\{e\})\cup \{\tilde{e}\}, \mathcal{F},\mathcal{R},a) \end{array} $

      Para cada $c\in \mathcal{C}\cup (\mathcal{C}_{1}-\{e\})$ definamos $\tilde{c} =c$. Notese que el mapeo $c\rightarrow \tilde{c}$ es una biyeccion entre el conjunto de nombres de constante de $\tau _{1}$ y el conjunto de nombres de cte de $\tau _{2}$. Para cada $t\in T^{\tau _{1}}$ sea $\tilde{t}=$ resultado de reemplazar en $t$ cada ocurrencia de $c$ por $\tilde{c}$, para cada $c\in \mathcal{C}\cup \mathcal{C}_{1}$. Analogamente para una formula $ \psi \in F^{\tau _{1}}$, sea $\tilde{\psi}=$ resultado de reemplazar en $ \psi $ cada ocurrencia de $c$ por $\tilde{c}$, para cada $c\in \mathcal{C} \cup \mathcal{C}_{1}$. Notese que los mapeos $t\rightarrow \tilde{t}$ y $ \psi \rightarrow \tilde{\psi}$ son biyecciones naturales entre $T^{\tau _{1}} $ y $T^{\tau _{2}}$ y entre $F^{\tau _{1}}$ y $F^{\tau _{2}}$, respectivamente. Notese que cualesquiera sean $\psi_{1},\psi_{2}\in F^{\tau _{1}}$, tenemos que $\psi_{1}$ se deduce de $\psi_{2}$ por la regla de generalizacion con constante $c$ sii $\tilde{\psi}_{1}$ se deduce de $\tilde{\psi}_{2}$ por la regla de generalizacion con constante $\tilde{c} $. Para las otras reglas sucede lo mismo. Notese tambien que $c$ ocurre en $ \psi $ sii $\tilde{c}$ ocurre en $\tilde{\psi}.$ Mas aun notese que $c$ depende de $d$ en $(\mathbf{\varphi },\mathbf{J})$ sii $\tilde{c}$ depende de $\tilde{d}$ en $(\mathbf{\tilde{\varphi}},\mathbf{J})$, donde $\mathbf{ \tilde{\varphi}}=\widetilde{\mathbf{\varphi }_{1}}\dotsc\widetilde{\mathbf{ \varphi }_{n(\mathbf{\varphi })}}$. Ahora es facil chequear que $(\mathbf{ \tilde{\varphi}},\mathbf{J})$ es una prueba de $\varphi$ en $(\Sigma, \tau) $ basandose en que $(\mathbf{\varphi },\mathbf{J})$ es una prueba de $\varphi$ en $(\Sigma, \tau)$.
    \end{enumerate}
  \end{proof}

  % Lemma 161. Con prueba. Lemma 71.
  \begin{lemma} \label{lemma_71}
    \PN Sea $(\Sigma, \tau)$ una teoría.
    \begin{enumerate}
      \item Si $(\Sigma, \tau) \vdash \varphi_{1}, \dotsc, \varphi_{n}$ y $(\Sigma \cup \{\varphi_{1}, \dotsc,
      \varphi_{n}\},\tau ) \vdash \varphi$ entonces $(\Sigma, \tau) \vdash \varphi$.
      \item Si $(\Sigma, \tau) \vdash \varphi_{1}, \dotsc, \varphi_{n}$ y $\varphi$ se deduce por alguna regla universal
      a partir de $\varphi_{1}, \dotsc, \varphi_{n}$, entonces $(\Sigma, \tau) \vdash \varphi$.
      \item Si $(\Sigma, \tau)$ es inconsistente, entonces $(\Sigma, \tau) \vdash \varphi$, para toda sentencia
      $\varphi$.
      \item Si $(\Sigma, \tau)$ es consistente y $(\Sigma, \tau) \vdash \varphi$, entonces $(\Sigma \cup \{\varphi\},
      \tau)$ es consistente.
      \item $(\Sigma, \tau) \vdash (\varphi \rightarrow \psi)$ si y solo si $ (\Sigma \cup \{\varphi\}, \tau) \vdash
      \psi$.
      \item Si $(\Sigma, \tau) \not \vdash \lnot \varphi$, entonces $(\Sigma \cup \{\varphi\}, \tau)$ es consistente.
    \end{enumerate}
  \end{lemma}
  \begin{proof}
    \begin{enumerate}
      \item Haremos el caso $n=2.$ Supongamos entonces que $(\Sigma, \tau)\vdash \varphi_{1},\varphi_{2}$ y $(\Sigma \cup \{\varphi_{1},\varphi_{2}\},\tau )\vdash \varphi $. Para $i=1,2$, sea $(\varphi_{1}^{i}\dotsc\varphi_{n_{i}}^{i},J_{1}^{i}\dotsc J_{n_{i}}^{i})$ una prueba de $ \varphi_{i}$ en $(\Sigma, \tau)$. Sea $(\psi_{1}\dotsc\psi_{n},J_{1}\dotsc J_{n})$ una prueba de $\varphi$ en $(\Sigma \cup \{\varphi_{1},\varphi_{2}\},\tau )$. Notese que por el Lema 154 podemos suponer que estas tres pruebas no comparten ningun nombre de constante auxiliar y que tampoco comparten numeros asociados a hipotesis o tesis. Para cada $i=1, \dotsc, n$, definamos $\widetilde{J_{i}}$ de la siguiente manera.

      - Si $\psi_{i}=\varphi_{1}$ y $J_{i}=\mathrm{AXIOMAPROPIO}$, entonces $\widetilde{J_{i}}=\mathrm{EVOCACION}(\overline{n_{1}})$
      - Si $\psi_{i}=\varphi_{2}$ y $J_{i}=\mathrm{AXIOMAPROPIO}$, entonces $\widetilde{J_{i}}=\mathrm{EVOCACION}(\overline{n_{1}+n_{2}})$.
      - Si $\psi_{i}\notin \{\varphi_{1},\varphi_{2}\}$ y $J_{i}=\mathrm{ AXIOMAPROPIO}$, entonces $\widetilde{J_{i}}=\mathrm{AXIOMAPROPIO}$.
      - Si $J_{i}=\mathrm{AXIOMALOGICO}$, entonces $\widetilde{J_{i}}= \mathrm{AXIOMALOGICO}$
      - Si $J_{i}=\mathrm{CONCLUSION}$, entonces $\widetilde{J_{i}}=\mathrm{ CONCLUSION}$.
      - Si $J_{i}=\mathrm{HIPOTESIS}\bar{k}$, entonces $\widetilde{J_{i}}= \mathrm{HIPOTESIS}\bar{k}$
      - Si $J_{i}=\alpha P(\overline{l_{1}}, \dotsc, \overline{l_{k}})$, con $ \alpha \in \{\varepsilon \}\cup \{\mathrm{TESIS}\bar{k}:k\in \mathbb{N}\}$, entonces $\widetilde{J_{i}}=\alpha P(\overline{l_{1}+n_{1}+n_{2}}, \dotsc,  \overline{l_{k}+n_{1}+n_{2}})$
      Para cada $i=1, \dotsc, n_{2}$, definamos $\widetilde{J_{i}^{2}}$ de la siguiente manera.

      - Si $J_{i}^{2}=\mathrm{AXIOMAPROPIO}$, entonces $\widetilde{J_{i}^{2} }=\mathrm{AXIOMAPROPIO}$
      - Si $J_{i}^{2}=\mathrm{AXIOMALOGICO}$, entonces $\widetilde{J_{i}^{2} }=\mathrm{AXIOMALOGICO}$
      - Si $J_{i}^{2}=\mathrm{CONCLUSION}$, entonces $\widetilde{J_{i}^{2}}= \mathrm{CONCLUSION}$.
      - Si $J_{i}^{2}=\mathrm{HIPOTESIS}\bar{k}$, entonces $\widetilde{ J_{i}^{2}}=\mathrm{HIPOTESIS}\bar{k}$
      - Si $J_{i}^{2}=\alpha P(\overline{l_{1}}, \dotsc, \overline{l_{k}})$, con $\alpha \in \{\varepsilon \}\cup \{\mathrm{TESIS}\bar{k}:k\in \mathbb{N}\}$, entonces $\widetilde{J_{i}^{2}}=\alpha P(\overline{l_{1}+n_{1}}, \dotsc,  \overline{l_{k}+n_{1}})$
      Es facil chequear que

      $\displaystyle (\varphi_{1}^{1}\dotsc\varphi_{n_{1}}^{1}\varphi_{1}^{2}\dotsc\varphi_{n_{2}}^{2}\psi_{1}\dotsc\psi_{n},J_{1}^{1}\dotsc J_{n_{1}}^{1}\widetilde{ J_{1}^{2}}\dotsc\widetilde{J_{n_{2}}^{2}}\widetilde{J_{1}}\dotsc\widetilde{J_{n}}) $

      es una prueba de $\varphi$ en $(\Sigma, \tau)$

      \item Supongamos que $(\Sigma, \tau)\vdash \varphi_{1}, \dotsc, \varphi_{n}$ y que $\varphi$ se deduce por regla R a partir de $\varphi_{1}, \dotsc, \varphi_{n}$, con R universal. Notese que

      $\displaystyle \begin{array}{llll} 1.\; & \varphi_{1} & & \text{AXIOMAPROPIO} \\ 2.\; & \varphi_{2} & & \text{AXIOMAPROPIO} \\ \vdots & \vdots & & \vdots \\ n. & \varphi_{n} & & \text{AXIOMAPROPIO} \\ n+1 & \varphi & & \text{R}(\bar{1}, \dotsc, \bar{n}) \end{array} $

      es una prueba de $\varphi$ en $(\Sigma \cup \{\varphi_{1}, \dotsc, \varphi_{n}\},\tau )$, lo cual por (1) nos dice que $(\Sigma, \tau)\vdash \varphi $ .
      \item Si $(\Sigma, \tau)$ es inconsistente, entonces por definicion tenemos que $(\Sigma, \tau)\vdash \psi \wedge \lnot \psi $ para alguna sentencia $ \psi $. Dada una sentencia cualquiera $\varphi$ tenemos que $\varphi$ se deduce por la regla del absurdo a partir de $\psi \wedge \lnot \psi $ con lo cual (2) nos dice que $(\Sigma, \tau)\vdash \varphi $

      \item Supongamos $(\Sigma, \tau)$ es consistente y $(\Sigma, \tau)\vdash \varphi $. Si $(\Sigma \cup \{\varphi \},\tau )$ fuera inconsistente, entonces $(\Sigma \cup \{\varphi \},\tau )\vdash \psi \wedge \lnot \psi $, para alguna sentencia $\psi$, lo cual por (1) nos diria que $(\Sigma, \tau)\vdash \psi \wedge \lnot \psi $, es decir nos diria que $(\Sigma, \tau)$ es inconsistente.

      \item Supongamos $(\Sigma, \tau)\vdash (\varphi \rightarrow \psi )$. Entonces tenemos que $(\Sigma \cup \{\varphi \},\tau )\vdash (\varphi \rightarrow \psi ),\varphi $, lo cual por (2) nos dice que $(\Sigma \cup \{\varphi \},\tau )\vdash \psi $. Supongamos ahora que $(\Sigma \cup \{\varphi \},\tau )\vdash \psi $. Sea $(\varphi_{1}\dotsc\varphi_{n},J_{1}\dotsc,J_{n})$ una prueba de $\psi$ en $(\Sigma \cup \{\varphi \},\tau )$. Notese que podemos suponer que $J_{n}$ es de la forma $P(\overline{l_{1}}, \dotsc, \overline{l_{k}})$ . Definimos $\widetilde{J_{i}}=$ $\mathrm{TESIS}\bar{m}P(\overline{l_{1}+1} , \dotsc, \overline{l_{k}+1}$, donde $m$ es tal que ninguna $J_{i}$ es igual a $ \mathrm{HIPOTESIS}\bar{m}$. Para cada $i=1, \dotsc, n-1$, definamos $\widetilde{ J_{i}}$ de la siguiente manera.

      - Si $\varphi_{i}=\varphi $ y $J_{i}=\mathrm{AXIOMAPROPIO}$, entonces $\widetilde{J_{i}}=\mathrm{EVOCACION}(1)$
      - Si $\varphi_{i}\neq \varphi $ y $J_{i}=\mathrm{AXIOMAPROPIO}$, entonces $\widetilde{J_{i}}=\mathrm{AXIOMAPROPIO}$
      - Si $J_{i}=\mathrm{AXIOMALOGICO}$, entonces $\widetilde{J_{i}}= \mathrm{AXIOMALOGICO}$
      - Si $J_{i}=\mathrm{CONCLUSION}$, entonces $\widetilde{J_{i}}=\mathrm{ CONCLUSION}$
      - Si $J_{i}=\mathrm{HIPOTESIS}\bar{k}$ entonces $\widetilde{J_{i}}= \mathrm{HIPOTESIS}\bar{k}$
      - Si $J_{i}=\alpha P(\overline{l_{1}}, \dotsc, \overline{l_{k}})$, con $ \alpha \in \{\varepsilon \}\cup \{\mathrm{TESIS}\bar{k}:k\in \mathbb{N}\}$, entonces $\widetilde{J_{i}}=\alpha P(\overline{l_{1}+1}, \dotsc, \overline{l_{k}+1 })$
      Es facil chequear que

      $\displaystyle (\varphi \varphi_{1}\dotsc\varphi_{n}(\varphi \rightarrow \psi ),\text{ HIPOTESIS}\bar{m}\widetilde{J_{1}}\dotsc\widetilde{J_{n}}\text{CONCLUSION}) $

      es una prueba de $(\varphi \rightarrow \psi )$ en $(\Sigma, \tau)$
    \end{enumerate}
  \end{proof}

  % Lemma 162. Sin prueba. Theorem 72.
  \begin{theorem} \label{theorem_72}
    \PN \textbf{(Corrección)} $(\Sigma, \tau) \vdash \varphi$ implica $(\Sigma, \tau) \models \varphi$.
  \end{theorem}

  % Lemma 163. Con prueba. Corollary 73.
  \begin{corollary} \label{corollary_73}
    \PN Si $(\Sigma, \tau)$ tiene un modelo, entonces $(\Sigma, \tau)$ es consistente.
  \end{corollary}
  \begin{proof}
    Supongamos $\mathbf{A}$ es un modelo de $(\Sigma, \tau).$ Si $(\Sigma, \tau)$ fuera inconsistente, tendriamos que hay una $\varphi \in S^{t}$ tal que $ (\Sigma, \tau)\vdash (\varphi \wedge \lnot \varphi )$, lo cual por el Teorema de Correccion nos diria que $\mathbf{A}\models (\varphi \wedge \lnot \varphi )$
  \end{proof}

  % Lemma 164. Con prueba. Lemma 74.
  \begin{lemma} \label{lemma_74}
    \PN $\dashv \vdash$ es una relación de equivalencia.
  \end{lemma}
  \begin{proof}
    \PN \newline
    \begin{itemize}
      \item \textit{Reflexiva:} La relacion es reflexiva ya que $(\varphi \leftrightarrow \varphi )$ es un axioma logico.
      \item \textit{Simétrica:} Supongamos que $\varphi \dashv \vdash \psi $, es decir $(\Sigma, \tau)\vdash \left( \varphi \leftrightarrow \psi \right) $. Ya que $((\varphi \leftrightarrow \psi )\leftrightarrow (\psi \leftrightarrow \varphi ))$ es un axioma logico, tenemos que $(\Sigma, \tau)\vdash ((\varphi \leftrightarrow \psi )\leftrightarrow (\psi \leftrightarrow \varphi ))$. Notese que $\left( \psi \leftrightarrow \varphi \right) $ se deduce de $((\varphi \leftrightarrow \psi )\leftrightarrow (\psi \leftrightarrow \varphi ))$ y $(\varphi \leftrightarrow \psi )$ por la regla de reemplazo, lo cual por (2) del Lema 155 nos dice que $(\Sigma, \tau)\vdash \left( \psi \leftrightarrow \varphi \right) $.

      \item \textit{Transitiva:}   Analogamente, usando la regla de transitividad se puede probar que $\dashv \vdash $ es transitiva.
    \end{itemize}
  \end{proof}

  % Lemma 165. Con prueba. Lemma 75.
  \begin{lemma} \label{lemma_75}
    \PN Dada una teoria $T = (\Sigma, \tau)$, se tiene que:
    \begin{enumerate}[(1)]
      \item $\{\varphi \in S^{\tau}: \varphi \ \text{es un teorema de T}\} \in S^{\tau}/ \dashv \vdash_{T}$
      \item $\{\varphi \in S^{\tau}: \varphi \ \text{es refutable en T}\} \in S^{\tau}/ \dashv \vdash_{T}$
    \end{enumerate}
  \end{lemma}
  \begin{proof}
  \end{proof}

  % Lemma 166. Con prueba. Lemma 76.
  \begin{lemma} \label{lemma_76}
    \PN Sea $T = (\Sigma, \tau)$ una teoría, entonces $(S^{\tau}/\mathrm{\dashv \vdash}, \SU^{T}, \IN^{T}, 0^{T},
    1^{T})$ es un álgebra de Boole.
  \end{lemma}
  \begin{proof}
  \end{proof}

  % Lemma 167. Con prueba. Lemma 77.
  \begin{lemma} \label{lemma_77}
    \PN Sea T una teoría y sea $\leq^{T}$ el orden parcial asociado al álgebra de Boole $\mathcal{A}_{T}$ (es decir
    $[\varphi]_{T} \leq^{T} [\psi]_{T}$ si y solo si $[\varphi]_{T} \ \SU^{T} \ [\psi]_{T} = [\psi]_{T})$, entonces se
    tiene que:
    \[
      [\varphi]_{T} \leq^{T} [\psi]_{T} \text{ si y solo si} \ T \vdash (\varphi \rightarrow \psi)
    \]
  \end{lemma}
  \begin{proof}
    En virtud de los lemas anteriores solo falta probar que

    $\displaystyle \begin{array}{rcl} \lbrack \varphi \rbrack\;\mathsf{s}\;\lbrack\varphi \rbrack^{c} & =& 1 \\ \lbrack \varphi \rbrack\;\mathsf{i}\;\lbrack\varphi \rbrack^{c} & =& 0 \end{array} $

    Dejamos al lector la prueba de estas igualdades.
  \end{proof}

  % Lemma 168. Con prueba. Lemma 78.
  \begin{lemma} \label{lemma_78}
    \PN Sean $\tau = (\mathcal{C}, \mathcal{F}, \mathcal{R}, a)$ y $\tau^{\prime} = (\mathcal{C}^{\prime},
    \mathcal{F}^{\prime}, \mathcal{R}^{\prime}, a^{\prime})$ tipos.
    \begin{enumerate}
      \item Si $\mathcal{C} \subseteq \mathcal{C}^{\prime}, \mathcal{F} \subseteq \mathcal{F}^{\prime}, \mathcal{R}
      \subseteq \mathcal{R}^{\prime}$ y $a^{\prime}\mid_{\mathcal{F} \cup \mathcal{R}} = a$, entonces $(\Sigma, \tau)
      \vdash \varphi$ implica $(\Sigma, \tau^{\prime}) \vdash \varphi$.
      \item Si $\mathcal{C} \subseteq \mathcal{C}^{\prime}, \mathcal{F} = \mathcal{F}^{\prime}, \mathcal{R} =
      \mathcal{R}^{\prime}$ y $a^{\prime} = a$, entonces $(\Sigma, \tau^{\prime}) \vdash \varphi$ implica $(\Sigma,
      \tau) \vdash \varphi$, cada vez que $\Sigma \cup \{\varphi\} \subseteq S^{\tau}$.
    \end{enumerate}
  \end{lemma}
  \begin{proof}
    \PN \newline
    \begin{enumerate}[(1)]
      \item Supongamos $(\Sigma, \tau)\vdash \varphi $. Entonces hay una prueba $ (\varphi_{1}\dotsc\varphi_{n},J_{1}\dotsc J_{n})$ de $\varphi$ en $(\Sigma, \tau)$. Sea $\mathcal{C}_{1}$ el conjunto de nombres de constante que ocurren en alguna $\varphi_{i}$ y que no pertenecen a $\mathcal{C}.$ Notese que aplicando varias veces el Lema 154 podemos obtener una prueba $ (\tilde{\varphi}_{1}\dotsc\tilde{\varphi}_{n},J_{1}\dotsc J_{n})$ de $\varphi$ en $(\Sigma, \tau)$ la cual cumple que los nombres de constante que ocurren en alguna $\psi_{i}$ y que no pertenecen a $\mathcal{C}$ no pertenecen a $ \mathcal{C}^{\prime}$. Pero entonces $(\tilde{\varphi}_{1}\dotsc\tilde{\varphi} _{n},J_{1}\dotsc J_{n})$ es una prueba de $\varphi$ en $(\Sigma ,\tau ^{\prime})$, con lo cual $(\Sigma ,\tau ^{\prime})\vdash \varphi $

      \item Supongamos $(\Sigma ,\tau ^{\prime})\vdash \varphi $. Entonces hay una prueba $(\mathbf{\varphi },\mathbf{J})$ de $\varphi$ en $(\Sigma ,\tau ^{\prime})$. Veremos que $(\mathbf{\varphi },\mathbf{J})$ es una prueba de $\varphi$ en $(\Sigma, \tau)$. Ya que $(\mathbf{\varphi },\mathbf{J})$ es una prueba de $\varphi$ en $(\Sigma ,\tau ^{\prime})$ hay un conjunto finito $\mathcal{C}_{1}$, disjunto con $\mathcal{C}^{\prime}$, tal que $( \mathcal{C}^{\prime}\cup \mathcal{C}_{1},\mathcal{F},\mathcal{R},a)$ es un tipo y cada $\mathbf{\varphi }_{i}$ es una sentencia de tipo $(\mathcal{C} ^{\prime}\cup \mathcal{C}_{1},\mathcal{F},\mathcal{R},a)$. Notese que $ \widetilde{\mathcal{C}_{1}}=\mathcal{C}_{1}\cup (\mathcal{C}^{\prime}- \mathcal{C})$ es tal que $(\mathcal{C}\cup \widetilde{\mathcal{C}_{1}}, \mathcal{F},\mathcal{R},a)$ es un tipo y cada $\mathbf{\varphi }_{i}$ es una sentencia de tipo $(\mathcal{C}\cup \widetilde{\mathcal{C}_{1}},\mathcal{F}, \mathcal{R},a)$, con lo cual $(\mathbf{\varphi },\mathbf{J})$ cunple el punto 1. de la definicion de prueba. Todos los otros puntos se cumplen en forma directa, exepto los puntos 4(f) y 4(g)i para los cuales es necesario notar que $\mathcal{C}\subseteq \mathcal{C}^{\prime}$.
    \end{enumerate}
  \end{proof}

  % Lemma 169. Con prueba. Lemma 79.
  \begin{lemma} \label{lemma_79}
    \PN Sea $(\Sigma, \tau)$ una teoría y supongamos que $\tau$ tiene una cantidad infinita de nombres de constante que
    no ocurren en las sentencias de $\Sigma$, entonces para cada formula $\varphi =_{d} \varphi(v)$, se tiene que
    $\lbrack \forall v\varphi (v) \rbrack_{T} = \inf(\{\lbrack \varphi(t) \rbrack_{T}: t$ es un término cerrado$\})$.
  \end{lemma}
  \begin{proof}
    Primero notese que $\lbrack\forall v\;\varphi (v)\rbrack\leq \lbrack \varphi (t)\rbrack$, para todo termino cerrado $t$, ya que podemos dar la siguiente prueba:

    $\displaystyle \begin{array}{cllll} 1. & \forall v\;\varphi (v) & & & \text{HIPOTESIS}1 \\ 2. & \varphi (t) & & & \text{TESIS}1\text{PARTICULARIZACION}(1) \\ 3. & (\forall v\;\varphi (v)\rightarrow \varphi (t)) & & & \text{CONCLUSION } \end{array} $

    Supongamos ahora que $\lbrack\psi \rbrack\leq \lbrack \varphi (t)\rbrack$, para todo termino cerrado $t.$ Por hipotesis hay un nombre de cte $c\in \mathcal{C}$ el cual no ocurre en los elementos de $\Sigma \cup \{\psi ,\varphi (v)\}.$ Ya que $ \lbrack\psi \rbrack\leq \lbrack \varphi (c)\rbrack$, hay una prueba $(\varphi_{1}\dotsc\varphi_{n},J_{1}\dotsc J_{n})$ de $\left( \psi \rightarrow \varphi (c)\right) $ en $(\Sigma, \tau)$. Pero entonces es facil de chequear que la siguiente es una prueba en $(\Sigma ,(\mathcal{C}-\{c\},\mathcal{F},\mathcal{R},a))$ de $ \left( \psi \rightarrow \forall v\;\varphi (v)\right) $:
    $\displaystyle \begin{array}{rlcl} 1. & \varphi_{1} & & J_{1} \\ 2. & \varphi_{2} & & J_{2} \\ \vdots & \vdots & & \vdots \\ n. & \varphi_{n}=\left( \psi \rightarrow \varphi (c)\right) & & J_{n} \\ n+1. & \psi & & \text{HIPOTESIS}\bar{m} \\ n+2. & \varphi (c) & & \text{MODUSPONENS}(\bar{n},\overline{n+1}) \\ n+3. & \forall v\varphi (v) & & \text{TESIS}\bar{m}\text{GENERALIZACION}( \overline{n+2}) \\ n+4. & \left( \psi \rightarrow \forall v\varphi (v)\right) & & \text{ CONCLUSION} \end{array} $

    (con $m$ elejido suficientemente grande). Por el Lema 162 tenemos entonces que $(\Sigma, \tau)\vdash \left( \psi \rightarrow \forall v\;\varphi (v)\right) $
  \end{proof}

  % Lemma 170. Con prueba. Lemma 80.
  \begin{lemma} \label{lemma_80}
    \PN \textbf{(Coincidencia)} Sean $\tau$ y $\tau^{\prime}$ dos tipos cualesquiera y sea $\tau_{\cap}$ dado por:
    \begin{itemize}
      \item $\mathcal{C}_{\cap} = \mathcal{C} \cap \mathcal{C}^{\prime}$
      \item $\mathcal{F}_{\cap} = \{f \in \mathcal{F} \cap \mathcal{F}^{\prime}: a(f) = a^{\prime}(f)\}$
      \item $\mathcal{R}_{\cap} = \{r \in \mathcal{R} \cap \mathcal{R}^{\prime}: a(r) = a^{\prime}(r)\}$
      \item $a_{\cap} = a\mid_{\mathcal{F}_{\cap} \cup \mathcal{R}_{\cap}}$
    \end{itemize}

    \PN Sean $\mathbf{A}$ y $\mathbf{A}^{\prime}$ modelos de tipo $\tau$ y $\tau^{\prime}$ respectivamente. Supongamos
    que $A = A^{\prime}$ y que $c^{\mathbf{A}} = c^{\mathbf{A}^{\prime}}$, para cada $c \in \mathcal{C}_{\cap},
    f^{\mathbf{A}} = f^{\mathbf{A}^{\prime}}$, para cada $f \in \mathcal{F}_{\cap}$ y $r^{\mathbf{A}} =
    r^{\mathbf{A}^{\prime}}$, para cada $r \in \mathcal{R}_{\cap}$, entonces:
    \begin{enumerate}[(a)]
      \item Para cada $t =_{d} t(\vec{v}) \in T^{\tau_{\cap}}$ se tiene que $t^{\mathbf{A}} \lbrack \vec{a} \rbrack =
      t^{\mathbf{A}^{\prime}} \lbrack \vec{a} \rbrack$, para cada $\vec{a} \in A^{n}$.
      \item Para cada $\varphi =_{d} \varphi (\vec{v}) \in F^{\tau_{\cap}}$ se tiene que:
      \[
        \mathbf{A} \models \varphi \lbrack \vec{a} \rbrack \ \text{si y solo si} \ \mathbf{A}^{\prime} \models \varphi
        \lbrack \vec{a} \rbrack
      \]
      \item Si $\Sigma \cup \{\varphi\} \subseteq S^{\tau_{\cap}}$, entonces:
      \[
        (\Sigma, \tau) \models \varphi \ \text{si y solo si} \ (\Sigma, \tau^{\prime}) \models \varphi
      \]
    \end{enumerate}
  \end{lemma}
  \begin{proof}
    (a) y (b) son directos por induccion.

    (c) Supongamos que $(\Sigma, \tau)\models \varphi $. Sea $\mathbf{A} ^{\prime}$ un modelo de $\tau ^{\prime}$ tal que $\mathbf{A}^{\prime}\models \Sigma $. Sea $a\in A^{\prime}$ un elemento fijo. Sea $\mathbf{A}$ el modelo de tipo $\tau $ definido de la siguiente manera

    - universo de $\mathbf{A}=$ $A^{\prime}$
    - $c^{\mathbf{A}}=c^{\mathbf{A}^{\prime}}$, para cada $c\in \mathcal{ C}_{\cap }$,
    - $f^{\mathbf{A}}=f^{\mathbf{A}^{\prime}}$, para cada $f\in \mathcal{ F}_{\cap }$
    - $r^{\mathbf{A}}=r^{\mathbf{A}^{\prime}}$, para cada $r\in \mathcal{ R}_{\cap }$
    - $c^{\mathbf{A}}=a$, para cada $c\in \mathcal{C}-\widetilde{\mathcal{ C}}$
    - $f^{\mathbf{A}}(a_{1}, \dotsc, a_{a(f)})=a$, para cada $f\in \mathcal{F}- \mathcal{F}_{\cap }$, $a_{1}, \dotsc, a_{a(f)}\in A^{\prime}$
    - $r^{\mathbf{A}}=\varnothing $, para cada $r\in \mathcal{R-R}_{\cap }$
    Ya que $\mathbf{A}^{\prime}\models \Sigma $, (b) nos dice que $ \mathbf{A}\models \Sigma $, lo cual nos dice que $\mathbf{A}\models \varphi $ . Nuevamente por (b) tenemos que $\mathbf{A}^{\prime}\models \varphi $, con lo cual hemos probado que $(\Sigma ,\tau ^{\prime})\models \varphi $
  \end{proof}

  % Lemma 171. Con prueba. Lemma 81.
  \begin{lemma} \label{lemma_81}
    \PN Sea $\tau$ un tipo. Hay una infinitupla $(\gamma_{1}, \gamma_{2}, \dotsc) \in F^{\tau_{\mathbb{N}}}$ tal que:
    \begin{enumerate}
      \item $\lvert Li(\gamma_{j}) \rvert \leq 1$, para cada $j = 1, 2, \dotsc$
      \item Si $\lvert Li(\gamma )\rvert \leq 1$, entonces $\gamma = \gamma_{j}$, para algún $j \in \mathbb{N}$
    \end{enumerate}
  \end{lemma}
  \begin{proof}
    Notese que las formulas de tipo $\tau $ son palabras de algun alfabeto finito $A$. Dado un orden total estricto $< $ para $A$, podemos definir

    $\displaystyle \begin{array}{rcl} \gamma _{1} & =& \min\nolimits_{\alpha }^{< }\left( \alpha \in F^{\tau }\wedge \lvert Li(\alpha )\rvert \leq 1\right) \\ \gamma _{t+1} & =& \min\nolimits_{\alpha }^{< }\left( \alpha \in F^{\tau }\wedge \lvert Li(\alpha )\rvert \leq 1\wedge (\forall i\in \omega )_{i\leq t}\alpha \neq \gamma _{i}\right) \end{array} $

    Claramente esta sucesion cumple (1) y es facil ver que tambien se cumple la propiedad (2).
  \end{proof}

  % Theorem 172. Con prueba. Theorem 82.
  \begin{theorem} \label{theorem_82}
    \PN \textbf{(Completitud) (G{\"o}del)} Sea $T = (\Sigma, \tau)$ una teoría de primer orden. Si $T \models \varphi$
    entonces $T \vdash \varphi$.
  \end{theorem}
  % \begin{proof} % TODO error de compilación.
  %   Primero probaremos completitud para el caso en que $\tau $ tiene una cantidad infinita de nombres de cte que no ocurren en las sentencias de $ \Sigma $. Lo probaremos por el absurdo, es decir supongamos que $\varphi_{0} $ es tal que $(\Sigma, \tau)\models \varphi_{0}$ y $(\Sigma, \tau)\not\vdash \varphi_{0}.$ Notese que ya que $(\Sigma, \tau)\not\vdash \varphi_{0}$, tenemos que $\lbrack\lnot \varphi_{0}\rbrack\not=0^{\mathcal{A}_{(\Sigma, \tau)}}.$ Para cada $j\in \mathbb{N}$, sea $w_{j}\in Var$ tal que $ Li(\gamma _{j})\subseteq \{w_{j}\}$. Para cada $j$, declaremos $\gamma _{j}=_{d}\gamma _{j}(w_{j})$. Notese que por el Lema 163 tenemos que $\inf \{\lbrack\gamma _{j}(t)\rbrack:t\in T_{c}^{\tau }\}=\lbrack\forall w_{j}\gamma _{j}(w_{j})\rbrack$, para cada $j=1,2,\dotsc$. Por el Teorema de Rasiova y Sikorski tenemos que hay un filtro primo de $\mathcal{A}_{(\Sigma, \tau)}$ , $\mathcal{U}$ el cual cumple:
  %
  %   (a) $\lbrack\lnot \varphi_{0}\rbrack\in \mathcal{U}$
  %   (b) para cada $j\in \mathbb{N}$, $\{\lbrack\gamma _{j}(t)\rbrack:t\in T_{c}^{\tau }\}\subseteq \mathcal{U}$ implica que $\lbrack\forall w_{j}\gamma _{j}(w_{j})\rbrack\in \mathcal{U}$
  %   Ya que la sucesion de las $\gamma _{i}$ cubre todas las formulas con a lo sumo una variable libre, podemos reescribir la propiedad (b) de la siguiente manera
  %
  %   (b)$^{\prime}$ para cada $\varphi =_{d}\varphi (v)\in F^{\tau }$, si $\{\lbrack\varphi (t)\rbrack:t\in T_{c}^{\tau }\}\subseteq \mathcal{U}$ entonces $ \lbrack\forall v\varphi (v)\rbrack\in \mathcal{U}$
  %   Definamos sobre $T_{c}^{\tau }$ la siguiente relacion:
  %
  %   $\displaystyle t\bowtie s\text{ si y solo si }\lbrack(t\equiv s)\rbrack\in \mathcal{U}\text{.} $
  %
  %   Veamos entonces que:
  %   (1) $\bowtie $ es de equivalencia.
  %   (2) Para cada $\varphi =_{d}\varphi (v_{1}, \dotsc, v_{n})\in F^{\tau }$, $ t_{1}, \dotsc, t_{n},s_{1}, \dotsc, s_{n}\in T_{c}^{\tau }$, si $t_{1}\bowtie s_{1}$, $ t_{2}\bowtie s_{2}$, $\dotsc$, $t_{n}\bowtie s_{n}$, entonces $\lbrack\varphi (t_{1}, \dotsc, t_{n})\rbrack\in \mathcal{U}$ si y solo si $\lbrack\varphi (s_{1}, \dotsc, s_{n})\rbrack\in \mathcal{U}$.
  %   (3) Para cada $f\in \mathcal{F}_{n}$, $ t_{1}, \dotsc, t_{n},s_{1}, \dotsc, s_{n}\in T_{c}^{\tau }$,
  %   $\displaystyle t_{1}\bowtie s_{1},t_{2}\bowtie s_{2}, \dotsc, \;t_{n}\bowtie s_{n}\text{ implica }f(t_{1}, \dotsc, t_{n})\bowtie f(s_{1}, \dotsc, s_{n}). $
  %
  %   Probaremos (2). Notese que
  %
  %   $\displaystyle (\Sigma, \tau)\vdash \left( (t_{1}\equiv s_{1})\wedge (t_{2}\equiv s_{2})\wedge \dotsc\wedge (t_{n}\equiv s_{n})\wedge \varphi (t_{1}, \dotsc, t_{n})\right) \rightarrow \varphi (s_{1}, \dotsc, s_{n}) $
  %
  %   lo cual nos dice que
  %   $\displaystyle \lbrack (t_{1}\equiv s_{1})\rbrack \ \IN \ \lbrack(t_{2}\equiv s_{2})\rbrack \ \IN \  \dotsc \ \IN \ \lbrack(t_{n}\equiv s_{n})\rbrack \ \IN \ \lbrack\varphi (t_{1}, \dotsc, t_{n})\rbrack\leq \lbrack \varphi (s_{1}, \dotsc, s_{n})\rbrack $
  %
  %   de lo cual se desprende que
  %   $\displaystyle \lbrack \varphi (t_{1}, \dotsc, t_{n})\rbrack\in \mathcal{U}\text{ implica }\lbrack\varphi (s_{1}, \dotsc, s_{n})\rbrack\in \mathcal{U} $
  %
  %   ya que $\mathcal{U}$ es un filtro. La otra implicacion es analoga.
  %   Para probar (3) podemos tomar $\varphi =\left( f(v_{1}, \dotsc, v_{n})\equiv f(s_{1}, \dotsc, s_{n})\right) $ y aplicar (2).
  %
  %   Definamos ahora un modelo $\mathbf{A}_{\mathcal{U}}$ de tipo $\tau $ de la siguiente manera:
  %
  %   - Universo de $\mathbf{A}_{\mathcal{U}}=T_{c}^{\tau }/\mathrm{\bowtie }$
  %   - $f^{\mathbf{A}_{\mathcal{U}}}(t_{1}/\mathrm{\bowtie }, \dotsc, t_{n}/ \mathrm{\bowtie })=f(t_{1}, \dotsc, t_{n})/\mathrm{\bowtie }$, $f\in \mathcal{F} _{n}$, $t_{1}, \dotsc, t_{n}\in T_{c}^{\tau }\;$
  %   - $r^{\mathbf{A}_{\mathcal{U}}}=\{(t_{1}/\mathrm{\bowtie }, \dotsc, t_{n}/ \mathrm{\bowtie }):\lbrackr(t_{1}, \dotsc, t_{n})\rbrack\in \mathcal{U}\}$, $r\in \mathcal{R} _{n}.$
  %   Notese que la definicion de $f^{\mathbf{A}_{\mathcal{U}}}$ es inambigua por (3). Probaremos las siguientes propiedades basicas:
  %
  %   (4) Para cada $t=_{d}t(v_{1}, \dotsc, v_{n})\in \TAU$, $ t_{1}, \dotsc, t_{n}\in T_{c}^{\tau }$, tenemos que
  %   $\displaystyle t^{\mathbf{A}_{\mathcal{U}}}\lbrack t_{1}/\mathrm{\bowtie }, \dotsc, t_{n}/\mathrm{ \bowtie }\rbrack=t(t_{1}, \dotsc, t_{n})/\mathrm{\bowtie } $
  %
  %   (5) Para cada $\varphi =_{d}\varphi (v_{1}, \dotsc, v_{n})\in F^{\tau }$, $ t_{1}, \dotsc, t_{n}\in T_{c}^{\tau }$, tenemos que
  %   $\displaystyle \mathbf{A}_{\mathcal{U}}\models \varphi \lbrack t_{1}/\mathrm{\bowtie } , \dotsc, t_{n}/\mathrm{\bowtie }\rbrack\text{ si y solo si }\lbrack\varphi (t_{1}, \dotsc, t_{n})\rbrack\in \mathcal{U}. $
  %
  %   La prueba de (4) es directa por induccion. Probaremos (5) por induccion en el $k$ tal que $\varphi \in F_{k}^{\tau }$. El caso $k=0$ es dejado al lector. Supongamos (5) vale para $\varphi \in F_{k}^{\tau }$. Sea $\varphi =_{d}\varphi (v_{1}, \dotsc, v_{n})\in F_{k+1}^{\tau }-F_{k}^{\tau }.$ Hay varios casos:
  %
  %   CASO $\varphi (v_{1}, \dotsc, v_{n})=\left( \varphi_{1}(v_{1}, \dotsc, v_{n})\vee \varphi_{2}(v_{1}, \dotsc, v_{n})\right) .$
  %
  %   Tenemos
  %
  %   $\displaystyle \begin{array}{c} \mathbf{A}_{\mathcal{U}}\models \varphi \lbrack t_{1}/\mathrm{\bowtie } , \dotsc, t_{n}/\mathrm{\bowtie }\rbrack \\ \Updownarrow \\ \mathbf{A}_{\mathcal{U}}\models \varphi_{1}\lbrack t_{1}/\mathrm{\bowtie } , \dotsc, t_{n}/\mathrm{\bowtie }\rbrack\text{ o }\mathbf{A}_{\mathcal{U}}\models \varphi_{2}\lbrack t_{1}/\mathrm{\bowtie }, \dotsc, t_{n}/\mathrm{\bowtie }\rbrack \\ \Updownarrow \\ \lbrack \varphi_{1}(t_{1}, \dotsc, t_{n})\rbrack\in \mathcal{U}\text{ o }\lbrack\varphi_{2}(t_{1}, \dotsc, t_{n})\rbrack\in \mathcal{U} \\ \Updownarrow \\ \lbrack \varphi_{1}(t_{1}, \dotsc, t_{n})\rbrack\ \mathsf{s\ }\lbrack\varphi_{2}(t_{1}, \dotsc, t_{n})\rbrack\in \mathcal{U} \\ \Updownarrow \\ \lbrack \left( \varphi_{1}(t_{1}, \dotsc, t_{n})\vee \varphi_{2}(t_{1}, \dotsc, t_{n})\right) \rbrack\in \mathcal{U} \\ \Updownarrow \\ \lbrack \varphi (t_{1}, \dotsc, t_{n})\rbrack\in \mathcal{U}. \end{array} $
  %
  %   CASO $\varphi (v_{1}, \dotsc, v_{n})=\forall v\varphi_{1}(v_{1}, \dotsc, v_{n},v).$
  %   Tenemos
  %
  %   $\displaystyle \begin{array}{c} \mathbf{A}_{\mathcal{U}}\models \varphi \lbrack t_{1}/\mathrm{\bowtie } , \dotsc, t_{n}/\mathrm{\bowtie }\rbrack \\ \Updownarrow \\ \mathbf{A}_{\mathcal{U}}\models \varphi_{1}\lbrack t_{1}/\mathrm{\bowtie } , \dotsc, t_{n}/\mathrm{\bowtie },t/\mathrm{\bowtie }\rbrack\text{, para todo }t\in T_{c}^{\tau } \\ \Updownarrow \\ \lbrack \varphi_{1}(t_{1}, \dotsc, t_{n},t)\rbrack\in \mathcal{U}\text{, para todo } t\in T_{c}^{\tau } \\ \Updownarrow \\ \lbrack \forall v\varphi_{1}(t_{1}, \dotsc, t_{n},v)\rbrack\in \mathcal{U} \\ \Updownarrow \\ \lbrack \varphi (t_{1}, \dotsc, t_{n})\rbrack\in \mathcal{U}. \end{array} $
  %
  %   CASO $\varphi (v_{1}, \dotsc, v_{n})=\exists v\varphi_{1}(v_{1}, \dotsc, v_{n},v).$
  %   Tenemos
  %
  %   $\displaystyle \begin{array}{c} \mathbf{A}_{\mathcal{U}}\models \varphi \lbrack t_{1}/\mathrm{\bowtie } , \dotsc, t_{n}/\mathrm{\bowtie }\rbrack \\ \Updownarrow \\ \mathbf{A}_{\mathcal{U}}\models \varphi_{1}\lbrack t_{1}/\mathrm{\bowtie } , \dotsc, t_{n}/\mathrm{\bowtie },t/\mathrm{\bowtie }\rbrack\text{, para algun }t\in T_{c}^{\tau } \\ \Updownarrow \\ \lbrack \varphi_{1}(t_{1}, \dotsc, t_{n},t)\rbrack\in \mathcal{U}\text{, para algun } t\in T_{c}^{\tau } \\ \Updownarrow \\ \lbrack \varphi_{1}(t_{1}, \dotsc, t_{n},t)\rbrack^{c}\not\in \mathcal{U}\text{, para algun }t\in T_{c}^{\tau } \\ \Updownarrow \\ \lbrack \lnot \varphi_{1}(t_{1}, \dotsc, t_{n},t)\rbrack\not\in \mathcal{U}\text{, para algun }t\in T_{c}^{\tau } \\ \Updownarrow \\ \lbrack \forall v\;\lnot \varphi_{1}(t_{1}, \dotsc, t_{n},v)\rbrack\not\in \mathcal{U} \\ \Updownarrow \\ \lbrack \forall v\;\lnot \varphi_{1}(t_{1}, \dotsc, t_{n},v)\rbrack^{c}\in \mathcal{U} \\ \Updownarrow \\ \lbrack \lnot \forall v\;\lnot \varphi_{1}(t_{1}, \dotsc, t_{n},v)\rbrack\in \mathcal{U } \\ \Updownarrow \\ \lbrack \varphi (t_{1}, \dotsc, t_{n})\rbrack\in \mathcal{U}. \end{array} $
  %
  %   Pero ahora notese que (5) en particular nos dice que para cada sentencia $ \psi \in S^{\tau }$, $\mathbf{A}_{\mathcal{U}}\models \psi $ si y solo si $ \lbrack\psi \rbrack\in \mathcal{U}.$ De esta forma llegamos a que $\mathbf{A}_{\mathcal{U }}\models \Sigma $ y $\mathbf{A}_{\mathcal{U}}\models \lnot \varphi_{0}$, lo cual contradice la suposicion de que $(\Sigma, \tau)\models \varphi_{0}. $
  %   Ahora supongamos que $\tau $ es cualquier tipo. Sean $s_{1}$ y $s_{2}$ un par de simbolos no pertenecientes a la lista
  %
  %   $\displaystyle \forall \ \ \exists \ \ \lnot \ \ \vee \ \ \wedge \ \ \rightarrow \ \ \leftrightarrow \ \ (\ \ )\ \ ,\ \equiv \ \ \mathsf{X}\ \ \mathit{0}\ \ \mathit{1}\ \ \dotsc\ \ \mathit{9}\ \ \mathbf{0}\ \ \mathbf{1}\ \ \dotsc\ \ \mathbf{9} $
  %
  %   y tales que ninguno ocurra en alguna palabra de $\mathcal{C}\cup \mathcal{F} \cup \mathcal{R}.$ Si $(\Sigma, \tau)\models \varphi $, entonces usando el Lema de Coincidencia se puede ver que $(\Sigma ,(\mathcal{C}\cup \{s_{1}s_{2}s_{1},s_{1}s_{2}s_{2}s_{1},\dotsc\},\mathcal{F},\mathcal{R} ,a))\models \varphi $, por lo cual
  %   $\displaystyle (\Sigma ,(\mathcal{C}\cup \{s_{1}s_{2}s_{1},s_{1}s_{2}s_{2}s_{1},\dotsc\}, \mathcal{F},\mathcal{R},a))\vdash \varphi . $
  %
  %   Pero por Lema 162, tenemos que $(\Sigma, \tau)\vdash \varphi .$
  % \end{proof}

  % Corollary 173. Con prueba. Corollary 83.
  \begin{corollary} \label{corollary_83}
    \PN Toda teoría consistente tiene un modelo.
  \end{corollary}
  \begin{proof}
    Supongamos $(\Sigma, \tau)$ es consistente y no tiene modelos. Entonces $ (\Sigma, \tau)\models \left( \varphi
    \wedge \lnot \varphi \right) $, con lo cual por completitud $(\Sigma, \tau)\vdash \left( \varphi \wedge \lnot
    \varphi \right) $, lo cual es absurdo.
  \end{proof}

  % Corollary 174. Con prueba. Corollary 84.
  \begin{corollary} \label{corollary_84}
    \PN \textbf{(Teorema de Compacidad)}
    \begin{enumerate}[(a)]
      \item Si $(\Sigma, \tau)$ es tal que $(\Sigma_{0}, \tau)$ tiene un modelo, para cada subconjunto finito
      $\Sigma_{0} \subseteq \Sigma$, entonces $(\Sigma, \tau)$ tiene un modelo.
      \item Si $(\Sigma, \tau) \models \varphi$, entonces hay un subconjunto finito $\Sigma_{0} \subseteq \Sigma$ tal
      que $(\Sigma_{0}, \tau) \models \varphi$.
    \end{enumerate}
  \end{corollary}
  \begin{proof}
    (a) Si $(\Sigma, \tau)$ fuera inconsistente habria un subconjunto finito $ \Sigma _{0}\subseteq \Sigma $ tal que la teoria $(\Sigma _{0},\tau )$ es inconsistente ($\Sigma _{0}$ puede ser formado con los axiomas de $\Sigma $ usados en una prueba que atestigue que $(\Sigma, \tau)\vdash \left( \varphi \wedge \lnot \varphi \right) $). O sea que $(\Sigma, \tau)$ es consistente por lo cual tiene un modelo.

    (b) Si $(\Sigma, \tau)\models \varphi $, entonces por completitud, $(\Sigma, \tau)\vdash \varphi $. Pero entonces hay un subconjunto finito $\Sigma _{0}\subseteq \Sigma $ tal que $(\Sigma _{0},\tau )\vdash \varphi $, es decir tal que $(\Sigma _{0},\tau )\models \varphi $ (correccion).
  \end{proof}

  % El número 175 es usado para un ejemplo.

  % Lemma 176. Con prueba. Lemma 85.
  \begin{lemma} \label{lemma_85}
    \PN Dados $t_{1}, \dotsc, t_{n}, t =_{d} t(x_{1}, \dotsc, x_{n}) \in \TAU$, se tiene que $t^{\mathbf{T}^{\tau}}
    \lbrack t_{1}, \dotsc, t_{n} \rbrack = t(t_{1}, \dotsc, t_{n})$.
  \end{lemma}
  \begin{proof}
    Para cada $k\geq 0$, sea

    - Teo$_{k}$: Dados $t_{1}, \dotsc, t_{n}\in \TAU$ y $ t=_{d}t(x_{1}, \dotsc, x_{n})\in T_{k}^{\tau }$, se tiene que $t^{\mathbf{T} ^{\tau }}\lbrack t_{1}, \dotsc, t_{n}\rbrack=t(t_{1}, \dotsc, t_{n})$.
    Veamos que es cierto Teo$_{0}$. Hay dos casos

    Caso $t=_{d}t(x_{1}, \dotsc, x_{n})=c\in \mathcal{C}$.

    Entonces tenemos

    $\displaystyle \begin{array}{cll} t^{\mathbf{T}^{\tau }}\lbrack t_{1}, \dotsc, t_{n}\rbrack & = & c^{\mathbf{T}^{\tau }} \\ & = & c \\ & = & t(t_{1}, \dotsc, t_{n}) \end{array} $

    Caso $t=_{d}t(x_{1}, \dotsc, x_{n})=x_{i}$, para algun $i$.

    Entonces tenemos

    $\displaystyle \begin{array}{cll} t^{\mathbf{T}^{\tau }}\lbrack t_{1}, \dotsc, t_{n}\rbrack & = & t_{i} \\ & = & t(t_{1}, \dotsc, t_{n}) \end{array} $

    Veamos que Teo$_{k}$ implica Teo$_{k+1}$. Supongamos que vale Teo$_{k}$. Sean $t_{1}, \dotsc, t_{n}\in \TAU$ y $t=_{d}t(x_{1}, \dotsc, x_{n})\in T_{k+1}^{\tau }-T_{k}^{\tau }$. Hay $f\in \mathcal{F}_{m}$, con $m\geq 1$, y terminos $s_{1}, \dotsc, s_{m}\in T_{k}^{\tau }$ tales que $t=f(s_{1}, \dotsc, s_{m})$ . Notese que $s_{i}=_{d}s_{i}(x_{1}, \dotsc, x_{n})$, $i=1, \dotsc, m$. Tenemos entonces que
    $\displaystyle \begin{array}{lll} t^{\mathbf{T}^{\tau }}\lbrack t_{1}, \dotsc, t_{n}\rbrack & = & f(s_{1}, \dotsc, s_{m})^{\mathbf{T} ^{\tau }}\lbrack t_{1}, \dotsc, t_{n}\rbrack \\ & = & f^{\mathbf{T}^{\tau }}(s_{1}^{\mathbf{T}^{\tau }}\lbrack t_{1}, \dotsc, t_{n}\rbrack, \dotsc, s_{m}^{\mathbf{T}^{\tau }}\lbrack t_{1}, \dotsc, t_{n}\rbrack) \\ & = & f^{\mathbf{T}^{\tau }}(s_{1}(t_{1}, \dotsc, t_{n}), \dotsc, s_{m}(t_{1}, \dotsc, t_{n})) \\ & = & f(s_{1}(t_{1}, \dotsc, t_{n}), \dotsc, s_{m}(t_{1}, \dotsc, t_{n})) \\ & = & t(t_{1}, \dotsc, t_{n}) \end{array} $

    con lo cual vale Teo$_{k+1}$
  \end{proof}

  % Lemma 177. Con prueba. Lemma 86.
  \begin{lemma} \label{lemma_86}
    \PN \textbf{(Universal Mapping Property)} Si $\mathbf{A}$ es cualquier $\tau$-álgebra y $F: Var \rightarrow A$, es
    una función cualquiera, entonces $F$ puede ser extendida a un homomorfismo $\bar{F}: \mathbf{T}^{\tau} \rightarrow
    \mathbf{A}$.
  \end{lemma}
  \begin{proof}
    Definamos $\bar{F}$ de la siguiente manera:

    $\displaystyle \bar{F}(t)=t^{\mathbf{A}}\lbrack(F(x_{1}),F(x_{2}),\dotsc)\rbrack $

    Es claro que $\bar{F}$ extiende a $F$. Veamos que es un homomorfismo. Dada $ c\in \mathcal{C}$, tenemos que
    $\displaystyle \begin{array}{lll} \bar{F}(c^{\mathbf{T}^{\tau }}) & = & \bar{F}(c) \\ & = & c^{\mathbf{A}}\lbrack(F(x_{1}),F(x_{2}),\dotsc)\rbrack \\ & = & c^{\mathbf{A}} \end{array} $

    Dados $f\in \mathcal{F}_{n}$, $t_{1}, \dotsc, t_{n}\in \TAU$ tenemos que
    $\displaystyle \begin{array}{lll} \bar{F}(f^{\mathbf{T}^{\tau }}(t_{1}, \dotsc, t_{n})) & = & \bar{F} (f(t_{1}, \dotsc, t_{n})) \\ & = & f(t_{1}, \dotsc, t_{n})^{\mathbf{A}}\lbrack(F(x_{1}),F(x_{2}),\dotsc)\rbrack \\ & = & f^{\mathbf{A}}(t_{1}^{\mathbf{A}}\lbrack(F(x_{1}),F(x_{2}),\dotsc)\rbrack, \dotsc, t_{n}^{ \mathbf{A}}\lbrack(F(x_{1}),F(x_{2}),\dotsc)\rbrack) \\ & = & f^{\mathbf{A}}(\bar{F}(t_{1}), \dotsc, \bar{F}(t_{n})) \end{array} $

    con lo cual hemos probado que $\bar{F}$ es un homomorfismo.
  \end{proof}

  % Lemma 178. Con prueba. Lemma 87.
  \begin{lemma} \label{lemma_87}
    \PN Todas las reglas excepto las reglas anteriores son universales en el sentido que si $\varphi$ se deduce de
    $\psi_{1}, \dotsc, \psi_{k}$ por alguna de estas reglas, entonces $\left((\psi_{1} \wedge \dotsc \wedge \psi_{k})
    \rightarrow \varphi \right)$ es una sentencia universalmente válida.
  \end{lemma}
  \begin{proof}
    Veamos que la regla de reemplazo es universal. Basta con ver por induccion en $k$ que

    - Teo$_{k}$: Sean $t,s\in \TAU$, $r\in T_{k}^{\tau }$ y sea $ \mathbf{A}$ una $\tau $-algebra tal que $t^{\mathbf{A}}\lbrack\vec{a}\rbrack=s^{\mathbf{A }}\lbrack\vec{a}\rbrack$, para cada $\vec{a}\in A^{\mathbb{N}}$. Entonces $r^{\mathbf{A} }\lbrack\vec{a}\rbrack=\tilde{r}^{\mathbf{A}}\lbrack\vec{a}\rbrack$, para cada $\vec{a}\in A^{ \mathbb{N}}$, donde $\tilde{r}$ es el resultado de reemplazar algunas ocurrencias de $t$ en $r$ por $s.$
    La prueba de Teo$_{0}$ es dejada al lector. Asumamos que vale Teo$_{k}$ y probemos que vale Teo$_{k+1}$. Sean $t,s\in \TAU$, $r\in T_{k+1}^{\tau }-T_{k}^{\tau }$ y sea $\mathbf{A}$ una $\tau $-algebra tal que $t^{\mathbf{A }}\lbrack\vec{a}\rbrack=s^{\mathbf{A}}\lbrack\vec{a}\rbrack$, para cada $\vec{a}\in A^{\mathbb{N}}$. Sea $\tilde{r}$ el resultado de reemplazar algunas ocurrencias de $t$ en $r$ por $s$. El caso $t=r$ es trivial. Supongamos entonces que $t\neq r$. Supongamos $r=f(r_{1}, \dotsc, r_{n})$, con $r_{1}, \dotsc, r_{n}\in T_{k}^{\tau }$ y $ f\in \mathcal{F}_{n}$. Notese que por Lema 123 tenemos que $\tilde{r}=f(\tilde{r}_{1}, \dotsc, \tilde{r}_{n})$, donde cada $\tilde{r} _{i} $ es el resultado de reemplazar algunas ocurrencias de $t$ en $r_{i}$ por $s$. Para $\vec{a}\in A^{\mathbb{N}}$ se tiene que

    $\displaystyle \begin{array}{cclll} r^{\mathbf{A}}\lbrack\vec{a}\rbrack & = & f(r_{1}, \dotsc, r_{n})^{\mathbf{A}}\lbrack\vec{a}\rbrack & & \\ & = & f^{\mathbf{A}}(r_{1}^{\mathbf{A}}\lbrack\vec{a}\rbrack, \dotsc, r_{n}^{\mathbf{A}}\lbrack\vec{ a}\rbrack) & & \\ & = & f^{\mathbf{A}}(\tilde{r}_{1}^{\mathbf{A}}\lbrack\vec{a}\rbrack, \dotsc, \tilde{r}_{n}^{ \mathbf{A}}\lbrack\vec{a}\rbrack) & & \text{por Teo}_{k} \\ & = & f(\tilde{r}_{1}, \dotsc, \tilde{r}_{n})^{\mathbf{A}}\lbrack\vec{a}\rbrack & & \\ & = & \tilde{r}^{\mathbf{A}}\lbrack\vec{a}\rbrack & & \end{array} $

    lo cual prueba Teo$_{k+1}$
    Veamos que la regla de substitucion es universal. Supongamos $\mathbf{A} \models t\approx s$, con $t=_{d}t(x_{1}, \dotsc, x_{n})$ y $ s=_{d}s(x_{1}, \dotsc, x_{n})$. Veremos que entonces $\mathbf{A}\models t(p_{1}, \dotsc, p_{n})\approx s(p_{1}, \dotsc, p_{n}).$ Supongamos que $ p_{i}=_{d}p_{i}(x_{1}, \dotsc, x_{m})$, para cada $i=1, \dotsc, n.$ Por (a) del Lema 146, tenemos que

    $\displaystyle \begin{array}{rcl} t(p_{1}, \dotsc, p_{n}) & =& _{d}t(p_{1}, \dotsc, p_{n})(x_{1}, \dotsc, x_{m}) \\ s(p_{1}, \dotsc, p_{n}) & =& _{d}s(p_{1}, \dotsc, p_{n})(x_{1}, \dotsc, x_{m}) \end{array} $

    Sea $\vec{a}\in A^{m}$. Tenemos que
    $\displaystyle \begin{array}{rcl} t(p_{1}, \dotsc, p_{n})^{\mathbf{A}}\left\lbrack \vec{a}\right\rbrack & = & t^{\mathbf{A}} \left\lbrack p_{1}^{\mathbf{A}}\left\lbrack \vec{a}\right\rbrack , \dotsc, p_{n}^{\mathbf{A}}\left\lbrack \vec{a}\right\rbrack \right\rbrack \\ & = & s^{\mathbf{A}}\left\lbrack p_{1}^{\mathbf{A}}\left\lbrack \vec{a}\right\rbrack , \dotsc, p_{n}^{\mathbf{A}}\left\lbrack \vec{a}\right\rbrack \right\rbrack \\ & = & s(p_{1}, \dotsc, p_{n})^{\mathbf{A}}\left\lbrack \vec{a}\right\rbrack \end{array} $

    lo cual nos dice que $\mathbf{A}\models t(p_{1}, \dotsc, p_{n})\approx s(p_{1}, \dotsc, p_{n})$.
  \end{proof}

  % Theorem 179. Con prueba. Theorem 88.
  \begin{theorem} \label{theorem_88}
    \PN \textbf{(Corrección)} Si $(\Sigma, \tau) \vdash_{ec} p \approx q$, entonces $(\Sigma, \tau) \models p \approx q$.
  \end{theorem}
  \begin{proof}
    Sea

    $\displaystyle p_{1}\approx q_{1}, \dotsc, p_{n}\approx q_{n} $

    una prueba ecuacional de $p\approx q$ en $(\Sigma, \tau).$ Usando el lema anterior se puede probar facilmente por induccion en $i$ que $(\Sigma, \tau)\models p_{i}\approx q_{i}$, por lo cual $(\Sigma, \tau)\models p\approx q. $
  \end{proof}

  % Theorem 180. Con prueba. Theorem 89.
  \begin{theorem} \label{theorem_89}
    \PN \textbf{(Completitud) (Birkhoff)} Sea $(\Sigma, \tau)$ una teoría ecuacional. Si $(\Sigma, \tau) \models p
    \approx q$, entonces $(\Sigma, \tau) \vdash_{ec} p \approx q$.
  \end{theorem}
  % \begin{proof} % TODO error de compilación.
  %   Supongamos $(\Sigma, \tau)\models p\approx q.$ Sea $\theta $ la siguiente relacion binaria sobre $\TAU$:
  %
  %   $\displaystyle \theta =\{(t,s):(\Sigma, \tau)\vdash _{ec}t\approx s\}. $
  %
  %   Dejamos al lector probar que $\theta $ es una congruencia de $\mathbf{T} ^{\tau }$. Veamos que
  %   (*) $t^{\mathbf{T}^{\tau }/\theta }\lbrack t_{1}/\theta , \dotsc, t_{n}/\theta \rbrack=t(t_{1}, \dotsc, t_{n})/\theta $, para todo $t_{1}, \dotsc, t_{n}$, $ t=_{d}t(x_{1}, \dotsc, x_{n})$
  %   Por Corolario 142 tenemos que
  %
  %   $\displaystyle t^{\mathbf{T}^{\tau }/\theta }\lbrack t_{1}/\theta , \dotsc, t_{n}/\theta \rbrack=t^{\mathbf{T} ^{\tau }}\lbrack t_{1}, \dotsc, t_{n}\rbrack/\theta $
  %
  %   Pero por Lema 170 tenemos que $t^{\mathbf{T}^{\tau }}\lbrack t_{1}, \dotsc, t_{n}\rbrack=t(t_{1}, \dotsc, t_{n})$ por lo cual (*) es verdadera.
  %   Veamos que $\mathbf{T}^{\tau }/\theta \models \Sigma .$ Sea $t\approx s$ un elemento de $\Sigma $, con $t=_{d}t(x_{1}, \dotsc, x_{n})$ y $ s=_{d}s(x_{1}, \dotsc, x_{n}).$ Veremos que $\mathbf{T}^{\tau }/\theta \models t\approx s$, es decir veremos que
  %
  %   $\displaystyle t^{\mathbf{T}^{\tau }/\theta }\lbrack t_{1}/\theta , \dotsc, t_{n}/\theta \rbrack=s^{\mathbf{T} ^{\tau }/\theta }\lbrack t_{1}/\theta , \dotsc, t_{n}/\theta \rbrack $
  %
  %   para todo $t_{1}/\theta , \dotsc, t_{n}/\theta \in \TAU/\theta $. Notese que
  %   $\displaystyle (\Sigma, \tau)\vdash _{ec}t(t_{1}, \dotsc, t_{n})\approx s(t_{1}, \dotsc, t_{n}) $
  %
  %   por lo cual $t(t_{1}, \dotsc, t_{n})/\theta =s(t_{1}, \dotsc, t_{n})/\theta .$ Por (*) tenemos entonces
  %   $\displaystyle t^{\mathbf{T}^{\tau }/\theta }\lbrack t_{1}/\theta , \dotsc, t_{n}/\theta \rbrack=t(t_{1}, \dotsc, t_{n})/\theta =s(t_{1}, \dotsc, t_{n})/\theta =s^{\mathbf{T}^{\tau }/\theta }\lbrack t_{1}/\theta , \dotsc, t_{n}/\theta \rbrack, $
  %
  %   lo cual nos dice que $\mathbf{T}^{\tau }/\theta $ satisface la identidad $ t\approx s.$
  %   Ya que $\mathbf{T}^{\tau }/\theta \models \Sigma $, por hipotesis tenemos que $\mathbf{T}^{\tau }/\theta \models p\approx q.$ Es decir que si $ p=_{d}p(x_{1}, \dotsc, x_{n})$ y $q=_{d}q(x_{1}, \dotsc, x_{n})$ tenemos que $p^{ \mathbf{T}^{\tau }/\theta }\lbrack t_{1}/\theta , \dotsc, t_{n}/\theta \rbrack=q^{\mathbf{T} ^{\tau }/\theta }\lbrack t_{1}/\theta , \dotsc, t_{n}/\theta \rbrack$, para todo $ t_{1}, \dotsc, t_{n}\in \TAU$. En particular, tomando $t_{i}=x_{i}$, $ i=1, \dotsc, n$ tenemos que
  %
  %   $\displaystyle p^{\mathbf{T}^{\tau }/\theta }\lbrackx_{1}/\theta , \dotsc, x_{n}/\theta \rbrack=q^{\mathbf{T} ^{\tau }/\theta }\lbrackx_{1}/\theta , \dotsc, x_{n}/\theta \rbrack $
  %
  %   lo cual por (*) nos dice que $p/\theta =q/\theta $, produciendo $(\Sigma, \tau)\vdash _{ec}p\approx q$.
  % \end{proof}

	\section{La aritmética de Peano}

  % Lemma 182. Con prueba. Lemma 91.
  \begin{lemma} \label{lemma_91}
    \PN $\mathbf{\omega}$ es un modelo de $Arit$.
  \end{lemma}
  \begin{proof}
    Sea $\psi =_{d}\psi (v_{1},...,v_{n},v)$, f\'{o}rmula de $\tau _{A}$. Veremos que $\mathbf{\omega }\vDash Ind_{\psi }$. Sea

    $\displaystyle \varphi =((\psi (\vec{v},0)\wedge \forall v\ (\psi (\vec{v},v)\rightarrow \psi (\vec{v},v+1)))\rightarrow \forall v\ \psi (\vec{v},v)) $

    Declaremos $\varphi =_{d}\varphi (v_{1},...,v_{n})$. Notese que $\mathbf{ \omega }\vDash Ind_{\psi }$ si y solo si para cada $a_{1},...,a_{n}\in \omega $ se tiene que $\mathbf{\omega }\vDash \varphi \lbrack \vec{a}]$. Sean $a_{1},...,a_{n}\in \omega $ fijos. Probaremos que $\mathbf{\omega } \vDash \varphi \lbrack \vec{a}]$. Notar que si
    $\displaystyle \mathbf{\omega }\nvDash (\psi (\vec{v},0)\wedge \forall v\ (\psi (\vec{v} ,v)\rightarrow \psi (\vec{v},v+1)))[\vec{a}] $

    entonces $\mathbf{\omega }\vDash \varphi \lbrack \vec{a}]$ por lo cual podemos hacer solo el caso en que
    $\displaystyle \mathbf{\omega }\vDash (\psi (\vec{v},0)\wedge \forall v\ (\psi (\vec{v} ,v)\rightarrow \psi (\vec{v},v+1)))[\vec{a}] $

    Sea $S=\{a\in \omega :\mathbf{\omega }\vDash \psi (\vec{v},v)[\vec{a},a]\}$. Ya que $\mathbf{\omega }\vDash \psi (\vec{v},0)[\vec{a}]$, es facil ver usando el lema de reemplazo que $\mathbf{\omega }\vDash \psi (\vec{v},v)[ \vec{a},0]$, lo cual nos dice que $0\in S$. Ya que $\mathbf{\omega }\vDash (\forall v\ (\psi (\vec{v},v)\rightarrow \psi (\vec{v},v+1)))[\vec{a}]$, tenemos que
    (1) Para cada $a\in \omega $, si $\mathbf{\omega }\vDash \psi (\vec{v} ,v)[\vec{a},a]$, entonces $\mathbf{\omega }\vDash \psi (\vec{v},v+1)[\vec{a} ,a]$.
    Pero por el lema de reemplazo, tenemos que $\mathbf{\omega }\vDash \psi ( \vec{v},v+1)[\vec{a},a]$ sii $\mathbf{\omega }\vDash \psi (\vec{v},v)[\vec{a} ,a+1]$, lo cual nos dice que

    (2) Para cada $a\in \omega $, si $\mathbf{\omega }\vDash \psi (\vec{v} ,v)[\vec{a},a]$, entonces $\mathbf{\omega }\vDash \psi (\vec{v},v)[\vec{a} ,a+1]$.
    Ya que (2) nos dice que $a\in S$ implica $a+1\in S$, tenemos que $S=\omega $ ya que $0\in S$. Es decir que para cada $a\in \omega $, se da que $\mathbf{ \omega }\vDash \psi (\vec{v},v)[\vec{a},a]$ lo cual nos dice que $\mathbf{ \omega }\vDash \forall v\ \psi (\vec{v},v)[\vec{a}]$.

    Es rutina probar que $\mathbf{\omega }$ satisface los otros 15 axiomas de $ Arit$. $\Box$
  \end{proof}

  % Proposition 183. Con prueba. Proposition 92.
  \begin{proposition} \label{proposition_92}
    \PN Hay un modelo de $Arit$ el cual no es isomorfo a $\mathbf{\omega}$.
  \end{proposition}
  \begin{proof}
    Sea $\tau =(\{0,1,\blacktriangle \},\{+^{2},.^{2}\},\{\leq ^{2}\},a)$ y sea $ \Sigma =\Sigma _{A}\cup \{\lnot (\widehat{n}\equiv \blacktriangle ):n\in \omega \}$. Por el Teorema de Compacidad la teoria $(\Sigma ,\tau )$ tiene un modelo $\mathbf{A}=(A,i)$. Ya que

    $\displaystyle \mathbf{A}\vDash \lnot (\widehat{n}\equiv \blacktriangle )\text{, para cada } n\in \omega $

    tenemos que
    $\displaystyle i(\blacktriangle )\neq \widehat{n}^{\mathbf{A}}\text{, para cada }n\in \omega $

    Por el Lema de Coincidencia la estructura $\mathbf{B}=(A,i\mid _{\{0,1,+,.,\leq \}})$ es un modelo de $Arit$. Ademas dicho lema nos garantiza que $\widehat{n}^{\mathbf{B}}=\widehat{n}^{\mathbf{A}}$, para cada $n\in \omega $, por lo cual tenemos que
    $\displaystyle i(\blacktriangle )\neq \widehat{n}^{\mathbf{B}}\text{, para cada }n\in \omega $

    Veamos que $\mathbf{B}$ no es isomorfo a $\mathbf{\omega }$. Supongamos $ F:\omega \rightarrow A$ es un isomorfismo de $\mathbf{\omega }$ en $\mathbf{B }$. Es facil de probar por induccion en $n$ que $F(n)=\widehat{n}^{\mathbf{B} }$, para cada $n\in \omega $. Pero esto produce un absurdo ya que nos dice que $i(\blacktriangle )$ no esta en la imagen de $F$. $\Box$
  \end{proof}

  % Lemma 184. Con prueba. Lemma 93.
  \begin{lemma} \label{lemma_93}
    \PN Las siguientes sentencias son teoremas de la aritmética de Peano:
    \begin{enumerate}[(1)]
      \item $\forall x \; 0 \leq x$
      \item $\forall x \; (x \leq 0 \rightarrow x \equiv 0)$
      \item $\forall x \forall y \; (x + y \equiv 0 \rightarrow (x \equiv 0 \wedge y \equiv 0))$
      \item $\forall x \; (\lnot (x \equiv 0) \rightarrow \exists z \ (x \equiv z + 1))$
      \item $\forall x \forall y \; (x < y \rightarrow x + 1 \leq y)$
      \item $\forall x \forall y \; (x < y + 1 \rightarrow x \leq y)$
      \item $\forall x \forall y \; (x \leq y+1 \rightarrow (x \leq y \vee x \equiv y+1))$
      \item $\forall x \forall y \; (\lnot y \equiv 0 \rightarrow \exists q \exists r \; x \equiv q.y+r \wedge r < y)$
    \end{enumerate}
  \end{lemma}
  \begin{proof}
    \begin{enumerate}[(1)]
      \item Prueba para $\forall x \; 0 \leq x$:

      \item Prueba para $\forall x \; (x \leq 0 \rightarrow x \equiv 0)$:

        $\begin{array}{lllll}
          1. & x_{0} \leq 0 &&& \text{HIPOTESIS}1 \\
          2. & \forall x \; 0 \leq x &&& \text{TEOREMA} \\
          3. & 0 \leq x_{0} &&& \text{PARTICULARIZACION}(2) \\
          4. & x_{0} \leq 0 \wedge 0 \leq x_{0} &&& \text{CONJUNCIONINTRODUCCION}(1,3) \\
          5. & \forall x_{1} \forall x_{2} \; ((x_{1} \leq x_{2} \wedge x_{2} \leq x_{1}) \rightarrow x_{1} \equiv
            x_{2}) &&& \text{AXIOMAPROPIO} \\
          6. & \forall x_{2} \; ((x_{0} \leq x_{2} \wedge x_{2} \leq x_{0}) \rightarrow x_{0} \equiv x_{2}) &&&
            \text{PARTICULARIZACION}(5) \\
          7. & ((x_{0}\leq 0\wedge 0\leq x_{0})\rightarrow x_{0}\equiv 0) &&& \text{PARTICULARIZACION}(6) \\
          8. & x_{0}\equiv 0 &&& \text{TESIS}1\text{MODUSPONENS}(4,7) \\
          9. & x_{0}\leq 0\rightarrow x_{0}\equiv 0 &&& \text{CONCLUSION} \\
          10. & \forall x\ (x\leq 0\rightarrow x\equiv 0) &&& \text{GENERALIZACION} (9)
        \end{array}$

        \item Prueba para $\forall x \forall y \; (x + y \equiv 0 \rightarrow (x \equiv 0 \wedge y \equiv 0))$:

          $\begin{array}{lllll}
            1. & x_{0} + y_{0} \equiv 0 &&& \text{HIPOTESIS}1 \\
            2. & (x_{0} + y_{0} \equiv 0 \leftrightarrow 0 \equiv x_{0} + y_{0}) &&& \text{AXIOMALOGICO} \\
            3. & 0\equiv x_{0} + y_{0} &&& \text{REEMPLAZO(1,2)} \\
            4. & \exists X_{3} \ (0 \equiv x_{0} + X_{3}) &&& \text{EXISTENCIAL(3)} \\
            5. & \forall x_{1}\forall x_{2}\;(x_{1}\leq x_{2}\leftrightarrow \exists x_{3}\;x_{2}\equiv x_{1}+x_{3}) &&& \text{PARTICULARIZACION}(2) \\
            6. & x_{0}\leq 0 &&& \text{CONJUNCIONINTRODUCCION}(1,3) \\
            7. & \forall x_{1}\forall x_{2}\;((x_{1}\leq x_{2}\wedge x_{2}\leq x_{1})\rightarrow x_{1}\equiv x_{2}) &&& \text{AXIOMAPROPIO} \\
            8. & \forall x_{2}\;((x_{0}\leq x_{2}\wedge x_{2}\leq x_{0})\rightarrow x_{0}\equiv x_{2}) &&& \text{PARTICULARIZACION}(5) \\
            9. & ((x_{0}\leq 0\wedge 0\leq x_{0})\rightarrow x_{0}\equiv 0) &&& \text{PARTICULARIZACION}(6) \\
            10. & x_{0}\equiv 0 &&& \text{TESIS}1\text{MODUSPONENS}(4,7) \\
            11. & x_{0} \leq 0 \rightarrow x_{0} \equiv 0 &&& \text{CONCLUSION} \\
            12. & \forall x \ (x \leq 0 \rightarrow x \equiv 0) &&& \text{GENERALIZACION}(9)
          \end{array}$

        \item Prueba para $\forall x \; (\lnot (x \equiv 0) \rightarrow \exists z \ (x \equiv z + 1))$:
        \item Prueba para $\forall x \forall y \; (x < y \rightarrow x + 1 \leq y)$:
        \item Prueba para $\forall x \forall y \; (x < y + 1 \rightarrow x \leq y)$:
        \item Prueba para $\forall x \forall y \; (x \leq y+1 \rightarrow (x \leq y \vee x \equiv y+1))$:
        \item Prueba para $\forall x \forall y \; (\lnot y \equiv 0 \rightarrow \exists q \exists r \; x \equiv q.y+r
          \wedge r < y)$:
    \end{enumerate}

  \end{proof}

  % Lemma 185. Sin prueba. Lemma 94.
  \begin{lemma} \label{lemma_94}
    \PN Sean $n, m \in \omega$. Las siguientes sentencias son teoremas de la aritmética de Peano:
    \begin{enumerate}
      \item $(+(\widehat{n},\widehat{m})\equiv \widehat{n+m})$
      \item $(.(\widehat{n},\widehat{m})\equiv \widehat{n.m})$
      \item $\forall x \; (x \leq \widehat{n} \rightarrow (x \equiv \widehat{0} \vee x \equiv \widehat{1} \vee \dotsc
        \vee x \equiv \widehat{n}))$
    \end{enumerate}
  \end{lemma}

  % Lemma 186. Sin prueba. Lemma 95.
  \begin{lemma} \label{lemma_95}
    \PN Para cada término cerrado $t$, tenemos que $Arit \vdash (t \equiv \widehat{t^{\mathbf{\omega}}})$.
  \end{lemma}

  % Lemma 187. Sin prueba. Lemma 96.
  \begin{lemma} \label{lemma_96}
    \PN Si $\varphi$ es una sentencia atómica o negación de atómica y $\mathbf{\omega} \models \varphi$, entonces $Arit
    \vdash \varphi$.
  \end{lemma}
  \begin{proof}
    Hay dos casos. Supongamos $\varphi =(t\equiv s)$, con $t,s$ terminos cerrados. Ya que $\mathbf{\omega }\models \varphi $, tenemos que $t^{\mathbf{ \omega }}=s^{\mathbf{\omega }}$ y por lo tanto $\widehat{t^{\mathbf{\omega }} }=\widehat{s^{\mathbf{\omega }}}$. Por el lema anterior tenemos que $ Arit\vdash (t\equiv \widehat{t^{\mathbf{\omega }}}),(s\equiv \widehat{s^{ \mathbf{\omega }}})$ lo cual, ya que $\widehat{t^{\mathbf{\omega }}}$ y $ \widehat{s^{\mathbf{\omega }}}$ son el mismo termino nos dice por la regla de transitividad que $Arit\vdash (t\equiv s)$. Supongamos $\varphi =(t\leq s) $, con $t,s$ terminos cerrados. Ya que $\mathbf{\omega }\models \varphi $ , tenemos que $t^{\mathbf{\omega }}\leq s^{\mathbf{\omega }}$ y por lo tanto hay un $k\in \omega $ tal que $t^{\mathbf{\omega }}+k=s^{\mathbf{\omega }}$. Se tiene entonces que $\widehat{t^{\mathbf{\omega }}+k}=\widehat{s^{\mathbf{ \omega }}}$. Por el lema anterior tenemos que $Arit\vdash +(\widehat{t^{ \mathbf{\omega }}},\widehat{k})\equiv \widehat{t^{\mathbf{\omega }}+k}$ lo cual nos dice que

    $\displaystyle Arit\vdash +(\widehat{t^{\mathbf{\omega }}},\widehat{k})\equiv \widehat{s^{ \mathbf{\omega }}} $

    Pero el lema anterior nos dice que
    $\displaystyle Arit\vdash (t\equiv \widehat{t^{\mathbf{\omega }}}),(s\equiv \widehat{s^{ \mathbf{\omega }}}) $

    y por lo tanto la regla de reemplazo nos asegura que $Arit\vdash +(t, \widehat{k})\equiv s$. Ya que
    $\displaystyle \forall x_{1}\forall x_{2}\;(x_{1}\leq x_{2}\leftrightarrow \exists x_{3}\;x_{2}\equiv x_{1}+x_{3}) $

    es un axioma de $Arit$, tenemos que $Arit\vdash (t\leq s)$.
  \end{proof}

  % Lemma 188. Con prueba. Lemma 97.
  \begin{lemma} \label{lemma_97}
    Sea $\varphi =_{d} \varphi(\vec{v}, v) \in F^{\tau_{A}}$. Supongamos $v$ es sustituible por $w$ en $\varphi$ y $w
    \notin \{v_{1}, \dotsc, v_{n}\}$, entonces:
    \[
      Arit \vdash \forall \vec{v}((\varphi(\vec{v}, 0) \wedge \forall v(\forall w(w < v \rightarrow \varphi(\vec{v}, w))
      \rightarrow \varphi(\vec{v} ,v))) \rightarrow \forall v\varphi(\vec{v}, v))
    \]
  \end{lemma}
  \begin{proof}
    Sea $\tilde{\varphi}=\forall w(w\leq v\rightarrow \varphi (\vec{v},w))$. Notar que $\tilde{\varphi}=_{d}\tilde{\varphi}(\vec{v},v)$. Sea $ IndCom_{\varphi }$ la sentencia

    $\displaystyle \forall \vec{v}((\varphi (\vec{v},0)\wedge \forall v(\forall w(w< v\rightarrow \varphi (\vec{v},w))\rightarrow \varphi (\vec{v} ,v)))\rightarrow \forall v\varphi (\vec{v},v)) $

    Salvo por el uso de algunos teoremas simples y el uso simultaneo de las reglas de particularizacion y generalizacion, la siguiente es la prueba buscada
    $\displaystyle \begin{array}{lllll} \;1. & (\varphi (\vec{c},0)\wedge \forall v(\forall w(w< v\rightarrow \varphi (\vec{c},w))\rightarrow \varphi (\vec{c},v)) &&& \text{HIPOTESIS}1 \\ \;2. & \;\;\;w_{0}\leq 0 &&& \text{HIPOTESIS}2 \\ \;3. & \;\;\;\forall x\;(x\leq 0\rightarrow x\equiv 0) &&& \text{TEOREMA} \\ \;4. & \;\;\;w_{0}\leq 0\rightarrow w_{0}\equiv 0 &&& \text{ PARTICULARIZACION}(3) \\ \;5. & \;\;\;w_{0}\equiv 0 &&& \text{MODUSPONENS}(2,4) \\ \;6. & \;\;\;\varphi (\vec{c},0) &&& \text{CONJUNCIONELIMINACION}(1) \\ \;7. & \;\;\;\varphi (\vec{c},w_{0}) &&& \text{TESIS}2\text{REEMPLAZO} (5,6) \\ \;8. & w_{0}\leq 0\rightarrow \varphi (\vec{c},w_{0}) &&& \text{ CONCLUSION} \\ \;9. & \tilde{\varphi}(\vec{c},0) &&& \text{GENERALIZACION}(8) \\ 10. & \;\;\;\tilde{\varphi}(\vec{c},v_{0}) &&& \text{HIPOTESIS}3 \\ 11. & \;\;\;\;\;\;w_{0}< v_{0}+1 &&& \text{HIPOTESIS}4 \\ 12. & \;\;\;\;\;\;\forall x,y\;x< y+1\rightarrow x\leq y &&& \text{TEOREMA } \\ 13. & \;\;\;\;\;\;w_{0}< v_{0}+1\rightarrow w_{0}\leq v_{0} &&& \text{ PARTICULARIZACION}(12) \\ 14. & \;\;\;\;\;\;w_{0}\leq v_{0} &&& \text{MODUSPONENS}(11,13) \\ 15. & \;\;\;\;\;\;w_{0}\leq v_{0}\rightarrow \varphi (\vec{c},w_{0}) &&& \text{PARTICULARIZACION}(10) \\ 16. & \;\;\;\;\;\;\varphi (\vec{c},w_{0}) &&& \text{TESIS}4\text{ MODUSPONENS}(14,15) \\ 17. & \;\;\;w_{0}< v_{0}+1\rightarrow \varphi (\vec{c},w_{0}) &&& \text{ CONCLUSION} \\ 18. & \;\;\;\forall w\;w< v_{0}+1\rightarrow \varphi (\vec{c},w) &&& \text{GENERALIZACION}(17) \\ 19. & \;\;\;\forall v(\forall w(w< v\rightarrow \varphi (\vec{c} ,w))\rightarrow \varphi (\vec{c},v)) &&& \text{CONJUNCIONELIMINACION}(1) \\ 20. & \;\;\;(\forall w(w< v_{0}+1\rightarrow \varphi (\vec{c},w))\rightarrow \varphi (\vec{c},v_{0}+1)) &&& \text{PARTICULARIZACION}(19) \\ 21. & \;\;\;\varphi (\vec{c},v_{0}+1) &&& \text{MODUSPONENS}(18,20) \\ 22. & \;\;\;\;\;\;w_{0}\leq v_{0}+1 &&& \text{HIPOTESIS}5 \\ 23. & \;\;\;\;\;\;\forall x,y\;x\leq y+1\rightarrow (x\leq y\vee x\equiv y+1) &&& \text{TEOREMA} \\ 24. & \;\;\;\;\;\;w_{0}\leq v_{0}+1\rightarrow (w_{0}\leq v_{0}\vee w_{0}\equiv v_{0}+1) &&& \text{PARTICULARIZACION}(23) \\ 25. & \;\;\;\;\;\;(w_{0}\leq v_{0}\vee w_{0}\equiv v_{0}+1) &&& \text{ MODUSPONENS}(22,24) \\ 26. & \;\;\;\;\;\;w_{0}\leq v_{0}\rightarrow \varphi (\vec{c},w_{0}) &&& \text{PARTICULARIZACION}(10) \\ 27. & \;\;\;\;\;\;\;\;\;w_{0}\equiv v_{0}+1 &&& \text{HIPOTESIS}6 \\ 28. & \;\;\;\;\;\;\;\;\;\varphi (\vec{c},w_{0}) &&& \text{TESIS}6\text{ REEMPLAZO}(21,27) \\ 29. & \;\;\;w_{0}\equiv v_{0}+1\rightarrow \varphi (\vec{c},w_{0}) &&& \text{CONCLUSION} \\ 30. & \;\;\;\;\;\;\varphi (\vec{c},w_{0}) &&& \text{TESIS}5\text{ DISJUNCIONELIMINACION}(25,26,29) \\ 31. & \;\;\;w_{0}\leq v_{0}+1\rightarrow \varphi (\vec{c},w_{0}) &&& \text{CONCLUSION} \\ 32. & \;\;\;\tilde{\varphi}(\vec{c},v_{0}+1) &&& \text{TESIS}3\text{ GENERALIZACION}(31) \\ 33. & \tilde{\varphi}(\vec{c},v_{0})\rightarrow \tilde{\varphi}(\vec{c} ,v_{0}+1) &&& \text{CONCLUSION} \\ 34. & \forall v\tilde{\varphi}(\vec{c},v)\rightarrow \tilde{\varphi}(\vec{c} ,v+1) &&& \text{GENERALIZACION}(33) \\ 35. & \tilde{\varphi}(\vec{c},0)\wedge \forall v\tilde{\varphi}(\vec{c} ,v)\rightarrow \tilde{\varphi}(\vec{c},v+1) &&& \text{ CONJUNCIONINTRODUCCION}(9,34) \\ 36. & Ind_{\tilde{\varphi}} &&& \text{AXIOMAPROPIO} \\ 37. & (\tilde{\varphi}(\vec{c},0)\wedge \forall v(\tilde{\varphi}(\vec{c} ,v)\rightarrow \tilde{\varphi}(\vec{c},v+1))\rightarrow \forall v\tilde{ \varphi}(\vec{c},v) &&& \text{PARTICULARIZACION}(36) \\ 38. & \forall v\tilde{\varphi}(\vec{c},v) &&& \text{MODUSPONENS}(35,37) \\ 39. & \tilde{\varphi}(\vec{c},v_{0}) &&& \text{PARTICULARIZACION}(38) \\ 40. & v_{0}\leq v_{0}\rightarrow \varphi (\vec{c},v_{0}) &&& \text{ PARTICULARIZACION}(39) \\ 41. & \forall x\;x\leq x &&& \text{AXIOMAPROPIO} \\ 42. & v_{0}\leq v_{0} &&& \text{PARTICULARIZACION}(41) \\ 43. & \varphi (\vec{c},v_{0}) &&& \text{MODUSPONENS}(40,42) \\ 44. & \forall v\varphi (\vec{c},v) &&& \text{TESIS}1\text{GENERALIZACION} (43) \\ 45. & (\varphi (\vec{c},0)\wedge \forall v(\forall w(w< v\rightarrow \varphi ( \vec{c},w))\rightarrow \varphi (\vec{c},v)))\rightarrow \forall v\varphi ( \vec{c},v) &&& \text{CONCLUSION} \\ 46. & IndCom_{\varphi } &&& \text{GENERALIZACION}(45) \end{array} $
  \end{proof}

  % Lemma 189. Con prueba. Lemma 98.
  \begin{lemma} \label{lemma_98}
    \PN Los conjuntos $T^{\tau_{A}^{e}}, F^{\tau_{A}^{e}}, T^{\tau_{A}}$ y $F^{\tau_{A}}$ son $\mathcal{A}$-recursivos.
  \end{lemma}
  \begin{proof}
    \PN Notese que los conjuntos \(T^{\tau _{A}^{e}}\), \(F^{\tau _{A}^{e}}\), \(T^{\tau _{A}}\) y \(F^{\tau _{A}}\) son
    \(\mathcal{A}\)-efectivamente computables (justifique). Entonces la Tesis de Church nos garantiza que dichos conjuntos
    son \(\mathcal{A}\)-recursivos.

    \PN A continuacion daremos una prueba de que dichos conjuntos son en realidad \( \mathcal{A}\)-primitivos recursivos.

    Veamos que $T^{\tau_{A}^{e}}$ es $A$-p.r. Fijemos un orden total estricto $ < $ sobre $A$. Sea $P=\lambda x[\ast ^{< }(x)\in T^{\tau_{A}^{e}}]$. Notese que $P(0)=0$ y $P(x+1)=1$ si y solo si se da alguna de las siguientes

    - $\ast ^{< }(x+1)\in \{0,1\}\cup Aux$
    - $(\exists u,v\in \omega )\ast ^{< }(x+1)=+(\mathrm{\ast }^{< }(u), \mathrm{\ast }^{< }(v))\wedge (P^{\downarrow }(x))_{u+1}\wedge (P^{\downarrow }(x))_{v+1}$
    - $(\exists u,v\in \omega )\ast ^{< }(x+1)=\mathrm{.}(\mathrm{\ast } ^{< }(u),\mathrm{\ast }^{< }(v))\wedge (P^{\downarrow }(x))_{u+1}\wedge (P^{\downarrow }(x))_{v+1}$
    Por el Lema 48 tenemos que $P$ es $A$-p.r., por lo cual $\chi _{T^{\tau_{A}^{e}}}=P\circ \#^{< }$ lo es. Notese que

    $\displaystyle t\in T^{\tau _{A}}\text{ sii }t\in T^{\tau_{A}^{e}}\wedge \triangle \text{ no ocurre en }t\wedge \Box \text{ no ocurre en }t $

    por lo cual $T^{\tau _{A}}$ es $A$-p.r.
  \end{proof}

  % Lemma 190. Con prueba. Lemma 99.
  \begin{lemma} \label{lemma_99}
    \PN Los siguientes predicados son $\mathcal{A}$-recursivos:
    \begin{enumerate}[(1)]
      \item \CL $v$ ocurre libremente en $\varphi$ a partir de $i$\CR: $\omega \times Var \times F^{\tau_{A}^{e}}
        \rightarrow \omega$
      \item \CL $v \in Li(\varphi)$\CR: $Var \times F^{\tau_{A}^{e}} \rightarrow \omega$
      \item \CL $v$ es sustituible por $t$ en $\varphi$\CR: $Var \times T^{\tau_{A}^{e}} \times F^{\tau_{A}^{e}}
        \rightarrow \omega$
    \end{enumerate}
  \end{lemma}
  \begin{proof}
    (1). Veamos que $P:\omega \times Var\times F^{\tau_{A}^{e}}\rightarrow \omega $, dado por

    $\displaystyle P(i,v,\varphi )=\left\{ \begin{array}{ccl} 1 & & \text{si }v\mathit{\ }\text{ocurre libremente en}\mathit{\ }\varphi \text{ a partir de }i \\ 0 & & \text{caso contrario} \end{array} \right. $

    es $B$-p.r.. Sea $R:\omega \times Var\rightarrow \omega $ el predicado dado por $R(x,v)=1$ si y solo si $\ast ^{< }((x)_{1})\in F^{\tau_{A}^{e}}$ y $v \mathit{\ }$ocurre libremente en$\mathit{\ }\ast ^{< }((x)_{1})$ a partir de $ (x)_{2}$. Sea $\mathrm{Nex}=\{\wedge ,\vee ,\rightarrow ,\leftrightarrow \}$ . Notese que $F_{0}^{\tau_{A}^{e}}$ es $A$-p.r. ya que
    $\displaystyle F_{0}^{\tau_{A}^{e}}=F^{\tau_{A}^{e}}\cap (A-\{\forall ,\exists ,\lnot ,\vee ,\wedge ,\rightarrow ,\leftrightarrow \})^{\ast } $

    Notese que $R(0,v)=0$, para cada $v\in Var$ y que $R(x+1,v)=1$ si y solo si $ (x+1)_{2}\geq 1$ y se da alguna de las siguientes:
    - $\ast ^{< }((x+1)_{1})\in F_{0}^{\tau_{A}^{e}}\wedge v$ ocurre en $ \ast ^{< }((x+1)_{1})$ a partir de $(x+1)_{2}$
    - $(\exists \varphi _{1},\varphi _{2}\in F^{\tau_{A}^{e}})(\exists \eta \in \mathrm{Nex})\ast ^{< }((x+1)_{1})=(\varphi _{1}\eta \varphi _{2})\wedge $
    $\ \ \ \ \ \ \ \ \ \ \ \ \ \ \ \ \ \ \ \ \ \ \ \ \ \left( (R^{\downarrow }(x,v))_{\left\langle \#^{< }(\varphi _{1}),(x+1)_{2}-1\right\rangle +1}\vee (R^{\downarrow }(x,v))_{\left\langle \#^{< }(\varphi _{2}),(x+1)_{2}-\left\vert (\varphi _{1}\eta \right\vert \right\rangle +1}\right) $

    - $(\exists \varphi _{1}\in F^{\tau_{A}^{e}})\ast ^{< }((x+1)_{1})=\lnot \varphi _{1}\wedge (R^{\downarrow }(x,v))_{\left\langle \#^{< }(\varphi _{1}),(x+1)_{2}-1\right\rangle +1}$
    - $(\exists \varphi _{1}\in F^{\tau_{A}^{e}})(\exists w\in Var)(Q\in \{\forall ,\exists \})\;w\neq v\wedge $
    $\ \ \ \ \ \ \ \ \ \ \ \ \ \ \ \ \ \ \ \ \ \ \ \ \ \ \ \ \ \ \ \ \ \ \ \ \ \ast ^{< }((x+1)_{1})=Qw\varphi _{1}\wedge (R^{\downarrow }(x,v))_{\left\langle \#^{< }(\varphi _{1}),(x+1)_{2}-\left\vert (Qw\right\vert \right\rangle +1}$

    Es decir que por el Lema 48 tenemos que $R$ es $A$ -p.r.. Notese que para $(i,v,\varphi )\in \omega \times Var\times F^{\tau_{A}^{e}}$, tenemos $P(i,v,\varphi )=R(\left\langle \#^{< }(\varphi ),i\right\rangle ,v)$. Ahora es facil obtener la funcion $P$ haciendo composiciones adecuadas con $R$.
  \end{proof}

  % Lemma 191. Con prueba. Lemma 100.
  \begin{lemma} \label{lemma_100}
    \PN Las funciones $\lambda svt[\downarrow_{v}^{t}(s)]$ y $\lambda \varphi vt[\downarrow_{v}^{t}(\varphi)]$ son
    $\mathcal{A}$-recursivas.
  \end{lemma}
  \begin{proof}
    Sea $< $ un orden total estricto sobre $A$. Sea $h:\omega \times Var\times T^{\tau_{A}^{e}}\rightarrow \omega $ dada por

    $\displaystyle h(x,v,t)=\left\{ \begin{array}{ccc} \#^{< }(\downarrow _{v}^{t}(\ast ^{< }(x))) & & \text{si }\ast ^{< }(x)\in T^{\tau_{A}^{e}} \\ 0 & & \text{caso contrario} \end{array} \right. $

    Sea $P:\omega \times \omega \times Var\times T^{\tau_{A}^{e}}\times A^{\ast }\rightarrow \omega $ tal que $P(A,x,v,t,\alpha )=1$ si y solo si se da alguna de las siguientes
    - $\ast ^{< }(x+1)\notin T^{\tau_{A}^{e}}\wedge \alpha =\varepsilon $
    - $\ast ^{< }(x+1)=v\wedge \alpha =t$
    - $\ast ^{< }(x+1)\in (\{0,1\}\cup Aux)-\{v\}\wedge \alpha =\ast ^{< }(x+1)$
    - $(\exists r,s\in T^{\tau_{A}^{e}})\ast ^{< }(x+1)=+(r,s)\wedge \alpha =+(\ast ^{< }((A)_{\#^{< }(r)+1}),\ast ^{< }((A)_{\#^{< }(s)+1}))$
    - $(\exists r,s\in T^{\tau_{A}^{e}})\ast ^{< }(x+1)=\mathrm{.} (r,s)\wedge \alpha =\mathrm{.}(\ast ^{< }((A)_{\#^{< }(r)+1}),\ast ^{< }((A)_{\#^{< }(s)+1}))$
    Notese que $P(h^{\downarrow }(x,v,t),x,v,t,\alpha )=1$ si y solo si ya sea $\ast ^{< }(x+1)\notin T^{\tau }$ y $\alpha =\varepsilon $ o $\ast ^{< }(x+1)\in T^{\tau }$ y $\alpha =\mathrm{\downarrow }_{v}^{t}(\ast ^{< }(x+1))$. Tenemos entonces

    $\displaystyle \begin{array}{rcl} h(0,v,t) & =& 0 \\ h(x+1,v,t) & =& \#^{< }(\min_{\alpha }^{< }P(h^{\downarrow }(x,v,t),x,v,t,\alpha )), \end{array} $

    por lo cual el Lema 48 nos dice que $h$ es $A$-p.r. Ahora es facil obtener la funcion $\downarrow _{v}^{t}(s):T^{\tau_{A}^{e}}\times Var\times T^{\tau_{A}^{e}}\rightarrow T^{\tau_{A}^{e}}$ haciendo composiciones adecuadas con $h$. $\Box$
  \end{proof}

  % Lemma 192. Con prueba. Lemma 101.
  \begin{lemma} \label{lemma_101}
    \PN El predicado $R: F^{\tau_{A}^{e}} \times F^{\tau_{A}^{e}} \times F^{\tau_{A}^{e}} \times F^{\tau_{A}^{e}}
    \rightarrow \omega$, dado por:
    \[
      R(\varphi, \tilde{\varphi}, \psi_{1}, \psi_{2}) = \left\{\begin{array}{cccl}
                                                                \begin{array}{c}
                                                                  1 \\
                                                                  \ \\
                                                                \end{array}
                                                                &&&
                                                                \begin{array}{cl}
                                                                  \text{si } \tilde{\varphi} = &\text{resultado de
                                                                    reemplazar algunas} \\
                                                                  &\text{(posiblemente } 0 \text{) ocurrencias de} \\
                                                                  &\psi_{1} \text{ en } \varphi \text{ por } \psi_{2}
                                                                \end{array} \\ 0 &&& \text{ caso contrario}
                                                              \end{array}\right.
    \]
    \PN es $\mathcal{A}$-recursivo.
  \end{lemma}
  \begin{proof}
    Sea $\mathrm{Nex}=\{\wedge ,\vee ,\rightarrow ,\leftrightarrow \}$. Sea $< $ un orden total estricto sobre $A$. Notese que $R(\varphi ,\tilde{\varphi} ,\psi_{1},\psi_{2})=1$ sii se da alguna de las siguientes

    - $\varphi =\tilde{\varphi}$
    - $(\varphi =\psi_{1}\wedge \tilde{\varphi}=\psi_{2})$
    - $(\exists \varphi _{1},\varphi _{2},\tilde{\varphi}_{1},\tilde{ \varphi}_{2}\in F^{\tau_{A}^{e}})(\exists \eta \in \mathrm{Nex})\varphi =(\varphi _{1}\eta \varphi _{2})\wedge \tilde{\varphi}=(\tilde{\varphi} _{1}\eta \tilde{\varphi}_{2})\wedge $
    $\;\;\;\;\;\;\;\;\;\;\;\;\;\;\;\;\;\;\;\;\;\;\;\;\;\;\;\;\;\;\;\;\;\;\;R( \varphi _{1},\tilde{\varphi}_{1},\psi_{1},\psi_{2})\wedge R(\varphi _{2}, \tilde{\varphi}_{2},\psi_{1},\psi_{2})$

    - $(\exists \varphi _{1},\tilde{\varphi}_{1}\in F^{\tau_{A}^{e}})\varphi =\lnot \varphi _{1}\wedge \tilde{\varphi}=\lnot \tilde{ \varphi}_{1}\wedge R(\varphi _{1},\tilde{\varphi}_{1},\psi_{1},\psi_{2})$
    - $(\exists \varphi _{1},\tilde{\varphi}_{1}\in F^{\tau_{A}^{e}})(\exists v\in Var)(\exists Q\in \{\forall ,\exists \})\varphi =Qv\varphi _{1}\wedge $
    $\ \ \ \ \ \ \ \ \ \ \ \ \ \ \ \ \ \ \ \ \ \ \ \ \ \ \ \ \ \ \ \ \ \ \ \ \ \ \ \ \ \ \ \ \ \ \ \ \ \ \ \ \ \ \ \ \ \ \ \ \ \ \tilde{\varphi}=Qv\tilde{ \varphi}_{1}\wedge R(\varphi _{1},\tilde{\varphi}_{1},\psi_{1},\psi_{2})$

    Se puede usar lo anterior para ver que $R^{\prime }:\omega \times F^{\tau_{A}^{e}}\times F^{\tau_{A}^{e}}\rightarrow \omega $, dado por

    $\displaystyle R^{\prime }(x,\psi_{1},\psi_{2})=\left\{ \begin{array}{cc} \begin{array}{c} 1 \\ \; \end{array} & \begin{array}{c} \text{si }\ast ^{< }((x)_{1}),\ast ^{< }((x)_{2})\in F^{\tau_{A}^{e}}\text{ y }\ast ^{< }((x)_{2})=\text{resultado de} \\ \text{reemplazar algunas ocurrencias de }\psi_{1}\text{ en }\ast ^{< }((x)_{1})\text{ por }\psi_{2} \end{array} \\ 0 & \text{caso contrario} \end{array} \right. $

    es $A$-p.r., via el Lema 48. Finalmente $R$ puede obtenerse haciendo composiciones adecuadas con $R^{\prime }$.
  \end{proof}

  % Lemma 193. Con prueba. Lemma 102.
  \begin{lemma} \label{lemma_102}
    \PN Sea $\Sigma$ un alfabeto finito. Sea $S \subseteq \Sigma^{\ast}$ un conjunto $\Sigma$-recursivo. El conjunto
    $S^{+}$ es $\Sigma$-recursivo.
  \end{lemma}
  \begin{proof}
    Notese que $\alpha \in S^{+}$ si y solo si

    $\displaystyle (\exists z\in \mathbb{N})(\forall i\in \mathbb{N})_{i\leq Lt(z)}\ast ^{< }((z)_{i})\in S\wedge \alpha =\mathrm{\subset }_{i=1}^{Lt(z)}\ast ^{< }((z)_{i}) $

    Dejamos al lector completar los detalles faltantes.
  \end{proof}

  % Lemma 194. Con prueba. Lemma 103.
  \begin{lemma} \label{lemma_103}
    \PN Los conjuntos $ModPon^{\tau^{e}_{A}}, Elect^{\tau^{e}_{A}}, Reemp^{\tau^{e}_{A}}, ConjInt^{\tau^{e}_{A}},
    ConjElim^{\tau^{e}_{A}}, EquivInt^{\tau^{e}_{A}}, \linebreak DisjElim^{\tau^{e}_{A}}, DisjInt^{\tau^{e}_{A}},
    EquivElim^{\tau^{e}_{A}}, Generaliz^{\tau^{e}_{A}}, Commut^{\tau^{e}_{A}}, Trans^{\tau^{e}_{A}},
    Exist^{\tau^{e}_{A}}, Evoc^{\tau^{e}_{A}}, \linebreak Absur^{\tau^{e}_{A}}, DivPorCas^{\tau^{e}_{A}},
    Partic^{\tau^{e}_{A}}$, son $\mathcal{A}$-recursivos.
  \end{lemma}
  \begin{proof}
    Veamos que $Reem_{2}^{\tau_{A}^{e}}$ es $A$-p.r.. Sea $Q:F^{\tau_{A}^{e}}\times F^{\tau_{A}^{e}}\times F^{\tau_{A}^{e}}\rightarrow \omega $ el predicado tal que $Q(\varphi ,\psi ,\sigma )=1$ si y solo si

    $(\exists \alpha \in (\forall Var)^{+})(\exists \psi_{1},\psi_{2}\in F^{\tau_{A}^{e}})\ \psi =\alpha (\psi_{1}\leftrightarrow \psi_{2})\wedge $
    $\ \ \ \ \ \ \ \ \ \ \ \ \ \ \ \ \ \ \ \ \ \ \ Li(\psi_{1})=Li(\psi_{2})\wedge \left( (\forall v\in Var)\ v\notin Li(\psi_{1})\vee v\text{ ocurre en }\alpha \right) $
    $\ \ \ \ \ \ \ \ \ \ \ \ \ \ \ \ \ \ \ \ \ \ \ \ \ \ \ \ \ \ \ \ \ \ \ \ \ \ \ \ \ \ \ \ \ \ \ \ \ \ \ \ \ \ \ \ \ \ \ \ \ \ \ \ \ \ \ \ \ \ \ \ \ \ \ \ \ \ \ \ \ \ \ \ \ \ \ \ \ \ \ \ \ \ \ \ \ \ \ \ \ \ \ \ \ \ \ \wedge R(\varphi ,\sigma ,\psi_{1},\psi_{2})$
    ($R$ es el predicado dado por el Lema 185). Es facil ver que $Q$ es $A$-p.r. y que $Reem_{2}^{\tau_{A}^{e}}=Q\mid _{S^{\tau_{A}^{e}}\times S^{\tau_{A}^{e}}\times S^{\tau_{A}^{e}}}$. $\Box$
  \end{proof}

  % Lemma 195. Con prueba. Lemma 104.
  \begin{lemma} \label{lemma_104}
    \PN El predicado \CL$\psi$ se deduce de $\varphi$ por generalización con constante $c$, con respecto a
    $\tau^{e}_{A}$\CR: $S^{\tau_{A}^{e}} \times S^{\tau_{A}^{e}} \times Aux \rightarrow \omega$ es $\mathcal{A}$-recursivo.
  \end{lemma}
  \begin{proof}
    Notese que $\psi $ se deduce de $\varphi $ por generalizacion con constante $ c$ si y solo si hay una formula $\gamma $ y una variable $v$ tales que

    - $Li(\gamma )=\{v\}$
    - cada ocurrencia de $v$ en $\gamma $ es libre
    - $c$ no ocurre en $\gamma $
    - $\varphi =\mathrm{\downarrow }_{v}^{c}(\gamma )\wedge \psi =\forall v\gamma $
    El lector podra usando esta equivalencia facilmente justificar que el predicado en cuestion es $A$-p.r..
  \end{proof}

  % Lemma 196. Sin prueba. Lemma 105.
  \begin{lemma} \label{lemma_105}
    \PN El predicado \CL$\psi$ se deduce de $\varphi$ por elección con constante $e$, con respecto a $\tau^{e}_{A}$\CR:
    $S^{\tau_{A}^{e}} \times S^{\tau_{A}^{e}} \times Aux \rightarrow \omega$ es $\mathcal{A}$-PR.
  \end{lemma}

  % Lemma 197. Sin prueba. Lemma 106.
  \begin{lemma} \label{lemma_106}
    \PN $AxLog^{\tau_{A}^{e}}$ es $\mathcal{A}$-recursivo.
  \end{lemma}

  % Lemma 198. Sin prueba. Lemma 107.
  \begin{lemma} \label{lemma_107}
    \PN Las funciones:
    \begin{equation*}
      \begin{aligned}
        S^{\tau_{A}^{e}+} &\rightarrow& \omega \\
        \mathbf{\varphi} &\rightarrow& n(\mathbf{\varphi}) \\
      \end{aligned}
      \qquad\qquad\qquad
      \begin{aligned}
        \omega \times S^{\tau_{A}^{e}+} &\rightarrow& S^{\tau_{A}^{e}} \cup \{\varepsilon\} \\
        (i, \mathbf{\varphi}) &\rightarrow& \mathbf{\varphi}_{i} \\
      \end{aligned}
    \end{equation*}
    \PN son $\mathcal{A}$-recursivas.
  \end{lemma}

  % Lemma 199. Con prueba. Lemma 108.
  \begin{lemma} \label{lemma_108}
    \PN $Just$ es $\mathcal{B}$-PR. Las funciones:
    \begin{equation*}
      \begin{aligned}
        Just^{+} &\rightarrow& \omega \\
        \mathbf{J} &\rightarrow& n(\mathbf{J}) \\
      \end{aligned}
      \qquad\qquad\qquad
      \begin{aligned}
        \omega \times Just^{+} &\rightarrow& Just \cup \{\varepsilon\} \\
        (i, \mathbf{J}) &\rightarrow& \mathbf{J}_{i} \\
      \end{aligned}
    \end{equation*}
    \PN son $\mathcal{B}$-recursivas.
  \end{lemma}
  \begin{proof}
    Sean

    $\displaystyle \begin{array}{rcl} R & =& \{\gamma :\gamma \text{ es el nombre de alguna regla}\} \\ T & =& \{\varepsilon \}\cup \{\text{TESIS}\bar{k}:k\in \mathbb{N}\} \end{array} $

    Notese que $Just$ es la union de una cantidad finita de conjuntos. Uno de ellos es el conjunto
    $\displaystyle L=\{\beta \gamma (\overline{l_{1}},...,\overline{l_{k}}):\gamma \in R\text{, cada }l_{j}\in \mathbb{N}\text{ y }\beta \in T\} $

    Veremos que $L$ es $B$-p.r., y dejamos al lector lor restantes. Dejamos al lector la prueba de que $T$ es $B$-p.r.. Notese que $\alpha \in L$ sii
    $\displaystyle (\exists \beta \in T)(\exists \gamma \in R)(\exists z\in \mathbb{N})\ \alpha =\beta \gamma (\left( \subset _{i=1}^{Lt(z)-1}\overline{(z)_{i}},\right) \overline{(z)_{Lt(z)}}) $

    Se deja al lector dar cotas de los cuantificadores, para poder aplicar el lema de cuantificacion acotada.
  \end{proof}

  % Lemma 200. Sin prueba. Lemma 109.
  \begin{lemma} \label{lemma_109}
    \PN El conjunto $\{\mathbf{J} \in Just^{+}: \mathbf{J}$ es balanceada $\}$ es $\mathcal{B}$-recursivo.
  \end{lemma}

  % Lemma 201. Sin prueba. Lemma 110.
  \begin{lemma} \label{lemma_110}
    \PN El predicado
    \[
      \begin{array}{rcl}
        \omega \times S^{\tau_{A}^{e}} \times S^{\tau_{A}^{e}+} \times Just^{+} &\rightarrow& \omega \\
        (i, \varphi, \mathbf{\varphi}, \mathbf{J}) &\rightarrow& \left\{
          \begin{array}{ccl}
            1 && \text{si }(\mathbf{\varphi}, \mathbf{J}) \text{ es adecuado y } \varphi \text{ es hipótesis de }
              \mathbf{\varphi}_{i} \text{ en }(\mathbf{\varphi}, \mathbf{J}) \\
            0 && \text{caso contrario}
          \end{array}\right.
      \end{array}
    \]
    \PN es $(\mathcal{A} \cup \mathcal{B})$-recursivo.
  \end{lemma}

  % Lemma 202. Sin prueba. Lemma 111.
  \begin{lemma} \label{lemma_111}
    \PN El predicado
    \[
      \begin{array}{rcl}
        \mathcal{C} \times \mathcal{C} \times S^{\tau_{A}^{e}+} \times Just^{+} &\rightarrow& \omega \\
        (e, d, \mathbf{\varphi}, \mathbf{J}) &\rightarrow& \left\{
          \begin{array}{ccl}
            1 && \text{si }(\mathbf{\varphi}, \mathbf{J}) \text{ es adecuado y } e \text{ depende de } d \text{ en }
              (\mathbf{\varphi}, \mathbf{J}) \\
            0 && \text{caso contrario}
          \end{array}\right.
      \end{array}
    \]
    \PN es $(\mathcal{A} \cup \mathcal{B})$-recursivo.
  \end{lemma}

  % Lemma 203. Sin prueba. Lemma 112.
  \begin{lemma} \label{lemma_112}
    \PN Sea $(\Sigma, \tau_{A})$ una teoría tal que $\Sigma$ es $\mathcal{A}$-recursivo (resp. $\mathcal{A}$-RE),
    entonces $Pruebas_{(\Sigma, \tau_{A})}$ es $(\mathcal{A} \cup \mathcal{B})$-recursivo (resp. $(\mathcal{A} \cup
    \mathcal{B})$-RE).
  \end{lemma}

  % Proposition 204. Con prueba. Proposition 113.
  \begin{lemma}
    \PN Si $(\Sigma, \tau_{A})$ es una teoría tal que $\Sigma$ es $\mathcal{A}$-RE., entonces $Teo_{(\Sigma, \tau_{A})}$
    es $\mathcal{A}$-RE.
  \end{lemma}
  \begin{proof}
    Ya que $Pruebas_{(\Sigma ,\tau _{A})}$ es $(A\cup B)$-r.e. tenemos que hay una funcion $F:\omega \rightarrow S^{\tau_{A}^{e}+}\times Just^{+}$ la cual cumple que $p_{1}^{0,2}\circ F$ y $p_{2}^{0,2}\circ F$ son $(A\cup B)$-r. y ademas $I_{F}=Pruebas_{(\Sigma ,\tau _{A})}$. Sea

    $\displaystyle \begin{array}{ccc} g:S^{\tau_{A}^{e}+} & \rightarrow & S^{\tau_{A}^{e}} \\ \mathbf{\varphi } & \rightarrow & \mathbf{\varphi }_{n(\mathbf{\varphi })} \end{array} $

    Por lemas anteriores $g$ es $A$-p.r.. Notese que $I_{(g\circ p_{1}^{0,2}\circ F)}=Teo_{(\Sigma ,\tau _{A})}$, lo cual dice que $ Teo_{(\Sigma ,\tau _{A})}$ es $(A\cup B)$-r.e.. Por el teorema de independencia del alfabeto tenemos que $Teo_{(\Sigma ,\tau _{A})}$ es $A$ -r.e.. $\Box$
  \end{proof}

  % Lemma 205. Con prueba. Proposition 114.
  \begin{lemma} \label{lemma_114}
    \PN Cualesquiera sean $z_{0}, \dotsc, z_{n} \in \omega, n \geq 0$, hay $x, y \in \omega$, tales que $\beta(x, y, i)
    = z_{i}, i = 0, \dotsc, n$.
  \end{lemma}
  \begin{proof}
    Dados $x,y,m\in \omega $ con $m\geq 1$, usaremos $x\equiv y(m)$ para expresar que $x$ es congruente a $y$ modulo $m$, es decir para expresar que $ x-y$ es divisible por $m$. Usaremos en esta prueba el Teorema Chino del Resto:

    - Dados $m_{0},...,m_{n},z_{0},...,z_{n}\in \omega $ tales que $ m_{0},...,m_{n}$ son coprimos de a pares, hay un $x\in \omega $ tal que $ x\equiv z_{i}(m_{i})$, para $i=0,...,n.$
    Sea $y=\max (z_{0},...,z_{n},n)!$. Sean $m_{i}=y(i+1)+1$, $i=0,...,n$. Veamos que $m_{0},...,m_{n}$ son coprimos de a pares. Supongamos $p$ divide a $m_{i}$ y a $m_{j}$ con $i< j$. Entonces $p$ divide a $m_{j}-m_{i}=y(j-i)$ y ya que $p$ no puede dividir a $y$, tenemos que $p$ divide a $j-i$. Pero ya que $j-i< n$ tenemos que $p< n$ lo cual es absurdo ya que implicaria que $p$ divide $y$.

    Por el Teorema Chino del Resto hay un $x$ tal que $x\equiv z_{i}(m_{i})$, para $i=0,...,n$. Ya que $z_{i}< m_{i}$, tenemos que

    $\displaystyle \beta (x,y,i)=r(x,y(i+1)+1)=r(x,m_{i})=z_{i}\text{, }i=0,...,n\text{.} $
  \end{proof}

  % Proposition 206. Con prueba. Proposition 115.
  \begin{proposition} \label{proposition_115}
    \PN Si $h$ es $\emptyset$-recursiva, entonces $h$ es representable.
  \end{proposition}
  \begin{proof}
    Supongamos $f:S_{1}\times ...\times S_{n}\rightarrow \omega $ y $g:\omega \times \omega \times S_{1}\times ...\times S_{n}\rightarrow \omega $ son representables, con $S_{1},...,S_{n}\subseteq \omega $. Probaremos que $ R(f,g):\omega \times S_{1}\times ...\times S_{n}\rightarrow \omega $ lo es. Para esto primero notese que para $t,x_{1},...,x_{n},z\in \omega $, las siguientes son equivalentes

    (1) $R(f,g)(t,\vec{x})=z$
    (2) hay $z_{0},...,z_{t}\in \omega $ tales que
    $\displaystyle \begin{array}{rcl} z_{0} & =& f(\vec{x}) \\ z_{i+1} & =& g(z_{i},i,\vec{x})\text{, }i=0,...,t-1 \\ z_{t} & =& z \end{array} $

    (3) hay $x,y\in \omega $ tales que
    $\displaystyle \begin{array}{rcl} \beta (x,y,0) & =& f(\vec{x}) \\ \beta (x,y,i+1) & =& g(\beta (x,y,i),i,\vec{x})\text{, }i=0,...,t-1 \\ \beta (x,y,t) & =& z \end{array} $

    Sean

    $\displaystyle \begin{array}{rcl} \varphi _{\beta } & =& _{d}\varphi _{\beta }(v_{1},v_{2},v_{3},v) \\ \varphi _{f} & =& _{d}\varphi _{f}(v_{1},...,v_{n},v) \\ \varphi _{g} & =& _{d}\varphi _{g}(v_{1},...,v_{n+2},v) \end{array} $

    formulas que representen a las funciones $\beta $, $f$ y $g$, respectivamente. Sean $w_{1},...,w_{n+1},w$, $ y_{1},y_{2},y_{3},y_{4},z_{1},z_{2}$ variables todas distintas y tales que cada una de las variebles libres de $\varphi _{\beta }$, $\varphi _{f}$ y $ \varphi _{g}$ es sustituible por cada una de las variables $ w_{1},...,w_{n+1},w$, $y_{1},y_{2},y_{3},y_{4},z_{1},z_{2}$. Sea $\varphi _{R(f,g)}=\varphi _{R(f,g)}(w_{1},...,w_{n+1},w)$ la siguiente formula
    $\exists z_{1},z_{2}\;(\exists y_{1}\varphi _{\beta }(z_{1},z_{2},0,y_{1})\wedge \varphi _{f}(w_{2},...,w_{n+1},y_{1}))\wedge $
    $\ \ \ \ \ \ \ \ \ \ \ \ \ \ \ \ \ \ \ \varphi _{\beta }(z_{1},z_{2},w_{1},w)\wedge \forall y_{2}(y_{2}< w_{1}\rightarrow \exists y_{3},y_{4}\;\varphi _{\beta }(z_{1},z_{2},y_{2}+1,y_{3})\wedge $
    $\ \ \ \ \ \ \ \ \ \ \ \ \ \ \ \ \ \ \ \ \ \ \ \ \ \ \ \ \ \ \ \ \ \ \ \ \ \ \ \ \ \ \ \ \ \ \ \ \ \ \ \ \ \ \ \ \ \ \ \ \ \ \ \ \ \ \ \ \ \ \ \ \ \varphi _{\beta }(z_{1},z_{2},y_{2},y_{4})\wedge \varphi _{g}(y_{4},y_{2},w_{2},...,w_{n+1},y_{3}))$
    Es facil usando (3) ver que la formula $\varphi _{R(f,g)}$ representa a $R(f,g)$.

    En forma analoga se puede probar que las reglas de composicion y minimizacion preservan representabilidad por lo cual ya que los elementos de $\mathrm{R}_{0}^{\emptyset }$ son representables, tenemos que lo es toda funcion $\emptyset $-r. $\Box$
  \end{proof}

  % Lemma 207. Con prueba. Lemma 116.
  \begin{lemma} \label{lemma_116}
    \PN Hay un predicado $P: \omega \times \omega \rightarrow \omega$ el cual es $\emptyset$-PR y tal que el predicado
    $Q = \lambda x\left[(\exists t \in \omega) P(t, x)\right]: \omega \rightarrow \omega$ no es $\emptyset$-recursivo.
  \end{lemma}
  \begin{proof}
    Sea $\Sigma =\Sigma _{p}$. Recordemos que el predicado

    $\displaystyle P_{1}=\lambda t\mathcal{P}\left[ i^{0,1}(t,\mathcal{P},\mathcal{P})=n( \mathcal{P})+1\right] $

    es $\Sigma _{p}$-p.r. ya que la funcion $i^{0,1}$ lo es. Notese que el dominio de $P_{1}$ es $\omega \times \mathrm{Pro}^{\Sigma _{p}}$. Por Lema 69 (de la materia de lenguajes) tenemos que
    $\displaystyle Halt^{\Sigma _{p}}=\lambda \mathcal{P}\left[ (\exists t\in \omega )\;P_{1}(t, \mathcal{P})\right] $

    no es $\Sigma _{p}$-recursivo. Sea $< $ un orden total sobre $\Sigma _{p}$. Definamos $P:\omega \times \omega \rightarrow \omega $ de la siguiente manera
    $\displaystyle P(t,x)=\left\{ \begin{array}{ccc} P_{1}(t,\ast ^{< }(x)) & \text{si} & \ast ^{< }(x)\in \mathrm{Pro}^{\Sigma _{p}} \\ 0 & \text{si} & \ast ^{< }(x)\notin \mathrm{Pro}^{\Sigma _{p}} \end{array} \right. $

    Claramente $P$ es $\Sigma _{p}$-p.r., por lo cual es $\emptyset $-p.r.. Sea $ Q=\lambda x\left[ (\exists t\in \omega )P(t,x)\right] .$ Notese que
    $\displaystyle Halt^{\Sigma _{p}}=Q\circ \#^{< }\mathrm{\mid }_{\mathrm{Pro}^{\Sigma _{p}}} $

    lo cual dice que $Q$ no es $\Sigma _{p}$-r. ya que de serlo, el predicado $ Halt^{\Sigma _{p}}$ lo seria. Por el teorema de independencia del alfabeto tenemos entonces que $Q$ no es $\emptyset $-recursivo.
  \end{proof}

  % Lemma 208. Con prueba. Lemma 117.
  \begin{lemma} \label{lemma_117}
    \PN No toda función representable es $\emptyset$-recursiva.
  \end{lemma}
  \begin{proof}
    HACER!
  \end{proof}

  % Lemma 209. Con prueba. Lemma 118.
  \begin{lemma} \label{lemma_118}
    \PN Si $Verd_{\mathbf{\omega}}$ es $\mathcal{A}$-RE, entonces es $\mathcal{A}$-recursivo.
  \end{lemma}
  \begin{proof}
    Supongamos $Verd_{\mathbf{\omega }}$ es $A$-r. e. Sea $f:\omega \rightarrow Verd_{\mathbf{\omega }}$ una funcion sobre y $A$-r. Sea $g:S^{\tau _{A}}\rightarrow S^{\tau _{A}}$, dada por

    $\displaystyle g(\varphi )=\left\{ \begin{array}{ccc} ^{\curvearrowright }\varphi & \;\; & \text{si }\left[ \varphi \right] _{1}=\lnot \\ \lnot \varphi & \;\; & \text{caso contrario} \end{array} \right. $

    Notar que $g$ es $A$-p.r. por lo cual $g\circ f$ es $A$-r. Ya que $I_{g\circ f}=S^{\tau _{A}}-Verd_{\mathbf{\omega }}$ (justifique), tenemos que $S^{\tau _{A}}-Verd_{\mathbf{\omega }}$ es $A$-r. e., por lo cual
    $\displaystyle A^{\ast }-Verd_{\mathbf{\omega }}=(A^{\ast }-S^{\tau _{A}})\cup (S^{\tau _{A}}-Verd_{\mathbf{\omega }}) $

    lo es. Es decir que $Verd_{\mathbf{\omega }}$ y su complemento son $A$-r.e. por lo cual $Verd_{\mathbf{\omega }}$ es $A$-r.
  \end{proof}

  % Lemma 210. Con prueba. Lemma 119.
  \begin{lemma} \label{lemma_119}
    \PN $Verd_{\mathbf{\omega}}$ no es $\mathcal{A}$-RE.
  \end{lemma}
  \begin{proof}
    Por el Lema 200 hay un predicado $\emptyset $-p.r., $ P:\omega \times \omega \rightarrow \omega $ tal que el predicado $Q=\lambda x \left[ (\exists t\in \omega )P(t,x)\right] :\omega \rightarrow \omega $ no es $\emptyset $-recursivo. Notese que $Q$ tampoco es $A$-recursivo. Ya que $P $ es representable, hay una formula $\varphi =_{d}\varphi (v_{1},v_{2},v)\in F^{\tau _{A}}$ la cual cumple

    $\displaystyle \mathbf{\omega }\models \varphi \left[ t,x,k\right] \text{ si y solo si } P(t,x)=k, $

    cualesquiera sean $t,x,k\in \omega .$ Sea $\psi =\varphi (v_{1},v_{2},1)$. Notese que $\psi =_{d}\psi (v_{1},v_{2})$ y que
    $\displaystyle \mathbf{\omega }\models \psi \left[ t,x\right] \text{ si y solo si }P(t,x)=1 \text{,} $

    cualesquiera sean $t,x\in \omega .$ Sea $\psi_{0}=\exists v_{1}\ \psi (v_{1},v_{2})$. Notese que $\psi_{0}=_{d}\psi_{0}(v_{2})$ y que
    $\displaystyle \mathbf{\omega }\models \psi_{0}\left[ x\right] \text{ si y solo si }Q(x)=1 $

    cualesquiera sea $x\in \omega $. Por el lema de reemplazo tenemos que para $ x\in \omega $,
    $\displaystyle \mathbf{\omega }\models \psi_{0}\left[ x\right] \text{ si y solo si } \mathbf{\omega }\models \psi_{0}(\widehat{x}) $

    (justifique), por lo cual
    $\displaystyle \mathbf{\omega }\models \psi_{0}(\widehat{x})\text{ si y solo si }Q(x)=1 $

    cualesquiera sea $x\in \omega $. Ya que $\psi_{0}(\widehat{x})$ es una sentencia,
    $\displaystyle \psi_{0}(\widehat{x})\in Verd_{\mathbf{\omega }}\text{ si y solo si }Q(x)=1 $

    Sea $h:\omega \rightarrow A^{\ast }$, dada por $h(x)=\psi_{0}(\widehat{x})$ . Es facil ver que $h$ es $A$-recursiva. Ya que $Q=\chi _{Verd_{\mathbf{ \omega }}}\circ h$ y $Q$ no es $A$-recursivo, tenemos que $\chi _{Verd_{ \mathbf{\omega }}}$ no es $A$-recursiva, es decir que $Verd_{\mathbf{\omega } }$ es un conjunto no $A$-recursivo. El lema anterior nos dice entonces que es $Verd_{\mathbf{\omega }}$ no es $A$-r.e..
  \end{proof}

  % Theorem 211. Con prueba. Theorem 120.
  \begin{theorem} \label{theorem_120}
    \PN (Godel) Si $\Sigma \subseteq Verd_{\mathbf{\omega}}$ es $\mathcal{A}$-RE, entonces $Teo_{(\Sigma, \tau_{A})}
    \subsetneq Verd_{\mathbf{\omega}}$.
  \end{theorem}
  \begin{proof}
    Por el Teorema de Correccion, tenemos que $Teo_{(\Sigma ,\tau _{A})}\subseteq Verd_{\mathbf{\omega }}$. Ya que $Teo_{(\Sigma ,\tau _{A})}$ es $A$-r.e y $Verd_{\mathbf{\omega }}$ no lo es, tenemos que $Teo_{(\Sigma ,\tau _{A})}\neq Verd_{\mathbf{\omega }}$. $\Box$
  \end{proof}

  % Corollary 212. Sin prueba. Corollary 121.
  \begin{corollary} \label{corollary_121}
    \PN Existe $\varphi \in S^{\tau_{A}}$ tal que $Arit \nvdash \varphi$ y $Arit \nvdash \lnot \varphi$.
  \end{corollary}

\end{document}
